\section{Direct and Semidirect Products and Abelian Groups}

\subsection{Direct Products}

\begin{exercise}\label{ex5.1.1}
    Show that the center of a direct product is the direct product of the centers:
    \[Z(G_1 \times G_2 \times \cdots \times G_n) = Z(G_1) \times Z(G_2) \times \cdots \times Z(G_n)\]
    Deduce that the direct product of two groups is abelian if and only if each of the factors is abelian.
\end{exercise}

\begin{sol}
    Let $G = G_1 \times G_2 \times \cdots \times G_n$. Let $z = (z_1, z_2, \ldots, z_n) \in Z(G)$. Then for any $g = (g_1, g_2, \ldots, g_n) \in G$, we have $zg = gz$, or equivalently,
    \[(z_1g_1, z_2g_2, \ldots, z_ng_n) = (g_1z_1, g_2z_2, \ldots, g_nz_n)\]
    This implies that $z_i g_i = g_i z_i$ for all $i$, so $z_i \in Z(G_i)$ for all $i$. Thus, $z \in Z(G_1) \times Z(G_2) \times \cdots \times Z(G_n)$, and we have shown that $Z(G) \subseteq Z(G_1) \times Z(G_2) \times \cdots \times Z(G_n)$.

    Conversely, let $z = (z_1, z_2, \ldots, z_n) \in Z(G_1) \times Z(G_2) \times \cdots \times Z(G_n)$. Then for any $g = (g_1, g_2, \ldots, g_n) \in G$, we have
    \[zg = (z_1g_1, z_2g_2, \ldots, z_ng_n) = (g_1z_1, g_2z_2, \ldots, g_nz_n) = gz\]
    Thus, $z \in Z(G)$, and we have shown that $Z(G_1) \times Z(G_2) \times \cdots \times Z(G_n) \subseteq Z(G)$.
\end{sol}

\begin{exercise}\label{ex5.1.2}
    Let $G_1, G_2, \ldots, G_n$ be groups and $G = G_1 \times \cdots \times G_n$. Let $I$ be a proper, nonempty subset of $\set{1, \ldots, n}$ and let $J = \set{1, \ldots, n} - I$. Define $G_I$ to be the set of elements of $G$ that have the identity $G_j$ in position $j$ for all $j \in J$.
    \begin{subproblems}
        \item Prove that $G_I$ is isomorphic to the direct product of the $G_i$ for $i \in I$.
        \item Prove that $G_I$ is a normal subgroup of $G$ and that $G/G_I \cong G_J$.
        \item Prove that $G \cong G_I \times G_J$.
    \end{subproblems}
\end{exercise}

\begin{solalph}
    \item Identifying $G_i$ with the subgroup of $G$ defined in Proposition 5.2, let $H$ be the product of the $G_i$ for $i \in I$ in $G$. Let $\phi : G_I \to H$ be the map defined by mapping the non-identity coordinates of $G_I$ to the corresponding coordinates of $H$. It is clear that $\phi$ is a homomorphism. Moreover, $\phi$ is surjective by definition of $H$. Finally, the kernel of $\phi$ is the set of all elements of $G_I$ that have the identity in all coordinates, which is just the identity element of $G$. Thus, $\phi$ is an isomorphism from $G_I$ to $H$, and we are done.
    \item Consider the map $\phi : G \to G_J$ defined by mapping the non-identity coordinates of $G$ to the corresponding coordinates of $G_J$. It is clear that $\phi$ is a homomorphism. Moreover, $\phi$ is surjective by definition of $G_J$. Finally, the kernel of $\phi$ is the set of all elements of $G$ that have the identity in all coordinates corresponding to $J$, which is exactly $G_I$. Thus, $\phi$ induces an isomorphism from $G/G_I$ to $G_J$, and we are done.
    \item Any element of $G_I \times G_J$ can be written as the product of an element of $G_I$ and an element of $G_J$ by taking the first and second coordinates, respectively. Thus, the map $\phi : G_I \times G_J \to G$ defined by $\phi((g_i)_{i \in I}, (g_j)_{j \in J}) = (g_1, g_2, \ldots, g_n)$ is a bijection. Injectivity is clear, since $\ker\phi$ is trivial. It is also a homomorphism since the group operation in $G$ is defined coordinate-wise. Thus, $\phi$ is an isomorphism from $G_I \times G_J$ to $G$, and we are done.
\end{solalph}

\newpage

\begin{exercise}\label{ex5.1.3}
    Under the notation of the preceding exercise, let $I$ and $K$ be any disjoint nonempty subsets of $\set{1, 2, \ldots, n}$, and let $G_I$ and $G_K$ be the subgroups of $G$ defined above. Prove that $xy = yx$ for all $x \in G_I$ and $y \in G_K$.
\end{exercise}

\begin{sol}
    Let $x \in G_I$ and $y \in G_K$. Then $x$ has the identity in all coordinates corresponding to $J = \set{1, 2, \ldots, n} - I$, and $y$ has the identity in all coordinates corresponding to $L = \set{1, 2, \ldots, n} - K$. Since $I$ and $K$ are disjoint, we have $J \cup L = \set{1, 2, \ldots, n}$, so every coordinate of $x$ and $y$ is the identity in at least one of them. Thus, when we compute the product $xy$, we have
    \[xy = (x_1y_1, x_2y_2, \ldots, x_ny_n) = (y_1x_1, y_2x_2, \ldots, y_nx_n) = yx\]
    as desired.
\end{sol}

\begin{exercise}\label{ex5.1.4}
    Let $A$ and $B$ be finite groups, and let $p$ be a prime. Prove that any Sylow $p$-subgroup of $A \times B$ is of the form $P \times Q$, where $P \in \syl_p(A)$ and $Q \in \syl_p(B)$. Prove that $n_p(A \times B) = n_p(A)n_p(B)$. Generalize both of these results to a direct product of any finite number of finite groups (so that the number of Sylow $p$-subgroups of a direct product is the product of the numbers of Sylow $p$-subgroups of the factors).
\end{exercise}

\begin{sol}
    Let $|A| = p^\alpha m$ and $|B| = p^\beta n$, where $p \nmid m$ and $p \nmid n$. Then $|A \times B| = |A||B| = p^{\alpha + \beta}mn$. Let $R \in \syl_p(A \times B)$. Then $|R| = p^{\alpha + \beta}$, so $R$ is a Sylow $p$-subgroup of $A \times B$. Let $\pi_A : A \times B \to A$ and $\pi_B : A \times B \to B$ be the projection maps. Then $\pi_A(R)$ is a $p$-subgroup of $A$, so $\pi_A(R) \subseteq P$ for some $P \in \syl_p(A)$. Similarly, $\pi_B(R) \subseteq Q$ for some $Q \in \syl_p(B)$. Thus, $R \subseteq P \times Q$. Since $|P \times Q| = |P||Q| = p^{\alpha + \beta}$, we must have $R = P \times Q$.

    To show $n_p(A \times B) = n_p(A)n_p(B)$, we note that the Sylow $p$-subgroups of $A \times B$ are precisely the subgroups of the form $P \times Q$, where $P \in \syl_p(A)$ and $Q \in \syl_p(B)$. Since there are $n_p(A)$ choices for $P$ and $n_p(B)$ choices for $Q$, we have $n_p(A \times B) = n_p(A)n_p(B)$.
\end{sol}

\begin{exercise}\label{ex5.1.5}
    Exhibit a non-normal subgroup of $Q_8 \times Z_4$ (note that every subgroup of each factor is normal).
\end{exercise}

\begin{sol}
    Note that $Q_8$ is non-abelian, so we may use the idea of a diagonal subgroup from \hyperref[ex3.1.38]{Exercise 3.1.38}. In particular, let $H = \gen{(i, 1)} = \set{(1, 0), (i, 1), (-1, 2), (-i, 3)}$. Then $H$ is a subgroup of $Q_8 \times Z_4$ that is isomorphic to $\gen{i}$. However,
    \[(j, 0)(i, 1)(j, 0)\inv = (ji(-j), 0 + 1 + 0) = (-i, 1)\]
    is not in $H$, so $H$ is not normal in $Q_8 \times Z_4$.
\end{sol}

\begin{exercise}\label{ex5.1.6}
    Show that all subgroups of $Q_8 \times E_{2^n}$ are normal.
\end{exercise}

\begin{sol}
    Let $H \leq Q_8 \times E_{2^n}$. Suppose $(a, b) \in H$ and $(x, y) \in Q_8 \times E_{2^n}$. Note that $(a\inv, b) \in H$, since every $b \in E_{2^n}$ is its own inverse. Then
    \[(x, y)(a, b)(x, y)\inv = (xax\inv, yby\inv) = (xax\inv, b)\]
    since $E_{2^n}$ is abelian. If $a \in Z(Q_8)$, then $xax\inv = a$, so $(x, y)(a, b)(x, y)\inv = (a, b) \in H$. If $a \notin Z(Q_8)$, then $a$ is one of $i, j, k$ or their inverses. In all cases, we have $xax\inv$ is one of $i, j, k$ or their inverses, so $xax\inv \in H$. Thus, $(x, y)(a, b)(x, y)\inv \in H$ for all $(a, b) \in H$ and $(x, y) \in Q_8 \times E_{2^n}$, so $H$ is normal in $Q_8 \times E_{2^n}$.
\end{sol}

\begin{exercise}\label{ex5.1.7}
    Let $G_1, G_2, \ldots, G_n$ be groups and $\pi$ be a fixed element of $S_n$. Prove that the map 
    \[\phi_\pi : G_1 \times G_2 \times \cdots \times G_n \to G_{\pi\inv(1)} \times G_{\pi\inv(2)} \times \cdots \times G_{\pi\inv(n)}\]
    defined by
    \[\phi_\pi(g_1, g_2, \ldots, g_n) = (g_{\pi\inv(1)}, g_{\pi\inv(2)}, \ldots, g_{\pi\inv(n)})\]
    is an isomorphism (so that changing the order of the factors in a direct product does not change the isomorphism type).
\end{exercise}

\begin{sol}
    Let $\pi \in S_n$ and let $\phi_\pi$ be the map defined in the problem statement. Let $g = (g_1, g_2, \ldots, g_n)$ and $h = (h_1, h_2, \ldots, h_n)$ be elements of $G_1 \times G_2 \times \cdots \times G_n$. Then
    \begin{align*}
        \phi_\pi(gh) & = \phi_\pi(g_1h_1, g_2h_2, \ldots, g_nh_n) \\
        & = (g_{\pi\inv(1)}h_{\pi\inv(1)}, g_{\pi\inv(2)}h_{\pi\inv(2)}, \ldots, g_{\pi\inv(n)}h_{\pi\inv(n)}) \\
        & = (g_{\pi\inv(1)}, g_{\pi\inv(2)}, \ldots, g_{\pi\inv(n)})(h_{\pi\inv(1)}, h_{\pi\inv(2)}, \ldots, h_{\pi\inv(n)}) \\
        & = \phi_\pi(g)\phi_\pi(h)
    \end{align*}
    Noting that $\pi \in S_n$ is a bijection, consider the map $\phi_{\pi\inv} : G_{\pi\inv(1)} \times G_{\pi\inv(2)} \times \cdots \times G_{\pi\inv(n)} \to G_1 \times G_2 \times \cdots \times G_n$ defined by
    \[\phi_{\pi\inv}(g_1, g_2, \ldots, g_n) = (g_{\pi(1)}, g_{\pi(2)}, \ldots, g_{\pi(n)})\]
    Then $\phi_{\pi\inv}$ is a homomorphism by the same argument as above. Moreover, $\phi_{\pi\inv}$ is the inverse of $\phi_\pi$, since
    \[\phi_\pi(\phi_{\pi\inv}(g_1, g_2, \ldots, g_n)) = \phi_\pi(g_{\pi(1)}, g_{\pi(2)}, \ldots, g_{\pi(n)}) = (g_1, g_2, \ldots, g_n)\]
    and
    \[\phi_{\pi\inv}(\phi_\pi(g_1, g_2, \ldots, g_n)) = \phi_{\pi\inv}(g_{\pi\inv(1)}, g_{\pi\inv(2)}, \ldots, g_{\pi\inv(n)}) = (g_1, g_2, \ldots, g_n).\]
    Therefore, $\phi_\pi$ is an isomorphism.
\end{sol}

\begin{exercise}\label{ex5.1.8}
    Let $G_1 = G_2 = \cdots = G_n$ and let $G = G_1 \times \cdots \times G_n$. Under the notation of the preceding exercise show that $\phi_\pi \in \text{Aut}(G)$. Show also that the map $\pi \mapsto \phi_\pi$ is an injective homomorphism of $S_n$ into $\text{Aut}(G)$. (In particular, $\phi_{\pi_1} \circ \phi_{\pi_2} = \phi_{\pi_1 \pi_2}$. It is at this point that the $\pi^{-1}$'s in the definition of $\phi_\pi$ are needed. The underlying reason for this is because if $e_i$ is the $n$-tuple with 1 in position $i$ and zeros elsewhere, $1 \leq i \leq n$, then $S_n$ acts on $\{e_1, \ldots, e_n\}$ by $\pi \cdot e_i = e_{\pi(i)}$; this is a left group action. If the $n$-tuple $(g_1, \ldots, g_n)$ is represented by $g_1 e_1 + \cdots + g_n e_n$, then this left group action on $\{e_1, \ldots, e_n\}$ extends to a left group action on sums by
    \[\pi \cdot (g_1 e_1 + g_2 e_2 + \cdots + g_n e_n) = g_1 e_{\pi(1)} + g_2 e_{\pi(2)} + \cdots + g_n e_{\pi(n)}.\]
    The coefficient of $e_{\pi(i)}$ on the right hand side is $g_i$, so the coefficient of $e_i$ is $g_{\pi^{-1}(i)}$. Thus the right hand side may be rewritten as $g_{\pi^{-1}(1)} e_1 + g_{\pi^{-1}(2)} e_2 + \cdots + g_{\pi^{-1}(n)} e_n$, which is precisely the sum attached to the $n$-tuple $(g_{\pi^{-1}(1)}, g_{\pi^{-1}(2)}, \ldots, g_{\pi^{-1}(n)})$. In other words, any permutation of the ``position vectors'' $e_1, \ldots, e_n$ (which fixes their coefficients) is the same as the inverse permutation on the coefficients (fixing the $e_i$'s). If one uses $\pi$'s in place of $\pi^{-1}$'s in the definition of $\phi_\pi$ then the map $\pi \mapsto \phi_\pi$ is not necessarily a homomorphism --- it corresponds to a \textit{right} group action.)
\end{exercise}

\begin{sol}
    Note that $\phi_\pi \in \aut(G)$ since $\phi_\pi$ is an isomorphism from $G$ to itself. Let $\pi_1, \pi_2 \in S_n$. Then
    \begin{align*}
        \phi_{\pi_1}(\phi_{\pi_2}(g_1, g_2, \ldots, g_n)) & = \phi_{\pi_1}(g_{\pi_2\inv(1)}, g_{\pi_2\inv(2)}, \ldots, g_{\pi_2\inv(n)}) \\
        & = (g_{\pi_2\inv(\pi_1\inv(1))}, g_{\pi_2\inv(\pi_1\inv(2))}, \ldots, g_{\pi_2\inv(\pi_1\inv(n))}) \\
        & = (g_{(\pi_1 \pi_2)\inv(1)}, g_{(\pi_1 \pi_2)\inv(2)}, \ldots, g_{(\pi_1 \pi_2)\inv(n)}) \\
        & = \phi_{\pi_1 \pi_2}(g_1, g_2, \ldots, g_n)
    \end{align*}
    Therefore, the map $\pi \mapsto \phi_\pi$ is a homomorphism. To see that it is injective, note that if $\phi_\pi$ is the identity automorphism, then for all $i$, $g_{\pi\inv(i)} = g_i$ for all $g_i \in G_i$, which implies $\pi$ is the identity permutation. Hence, the map is injective.
\end{sol}

\begin{exercise}\label{ex5.1.9}
    Let $G_i$ be a field $F$ for all $i$ and use the preceding exercise to show that the set of $n \times n$ matrices with one 1 in each row and each column is a subgroup of $\gl_n(F)$ isomorphic to $S_n$ (these matrices are the \textit{permutation matrices} since they simply permute the standard basis $e_1, \ldots, e_n$ (as above) of the $n$-dimensional vector space $F \times F \times \cdots \times F$).
\end{exercise}

\begin{sol}
    Let $G = F \times F \times \cdots \times F$ ($n$ times) and let $\pi \in S_n$. Then $\phi_\pi$ is an automorphism of $G$ that permutes the standard basis vectors $e_1, e_2, \ldots, e_n$. In particular, $\phi_\pi(e_i) = e_{\pi(i)}$ for all $i$. Thus, the matrix representation of $\phi_\pi$ with respect to the standard basis is a permutation matrix. Since the map $\pi \mapsto \phi_\pi$ is an injective homomorphism from $S_n$ to $\aut(G)$, the set of permutation matrices forms a subgroup of $\gl_n(F)$ that is isomorphic to $S_n$.
\end{sol}

\begin{exercise}\label{ex5.1.10}
    Let $p$ be a prime. Let $A$ and $B$ be two cyclic groups of order $p$ with generators $x$ and $y$ respectively. Set $E = A \times B$ so that $E$ is the elementary abelian group of order $p^2$: $E_{p^2}$ Prove that the distinct subgroups of $E$ of order $p$ are
    \[\langle x \rangle, \quad \langle xy \rangle, \quad \langle xy^2 \rangle, \quad \ldots, \quad \langle xy^{p-1} \rangle, \quad \langle y \rangle\] (note that there are $p + 1$ of them).
\end{exercise}

\begin{sol}
    Note that any proper, non-trivial subgroup of $E$ must have order $p$, and every element of $E$ has order $p$. The text shows that there are $p + 1$ subgroups of $E$ of order $p$, so it suffices to show that the subgroups listed in the problem statement are distinct. Note that $\gen x$ and $\gen y$ are distinct since they have different generators. Moreover, $\gen{xy^i}$ is distinct from $\gen x$ and $\gen y$ for all $i$, since $xy^i$ is not a power of $x$ or $y$ (as can be seen because $\gen x \cap \gen y$ is trivial). Finally, $\gen{xy^i}$ is distinct from $\gen{xy^j}$ for $i \neq j$, since if $\gen{xy^i} = \gen{xy^j}$, then $xy^i$ is a power of $xy^j$, which implies $y^{i-j}$ is a power of $x$, which is impossible since $\gen x \cap \gen y$ is trivial. Thus, all the subgroups listed in the problem statement are distinct.
\end{sol}

\begin{exercise}\label{ex5.1.11}
    Let $p$ be a prime, and let $n \in \zp$. Find a formula for the number of subgroups of order $p$ in the elementary abelian group $E_{p^n}$.
\end{exercise}

\begin{sol}
    Every element of $E_{p^n}$ has order $p$, so every subgroup of order $p$ is generated by a single non-identity element. Since there are $p^n - 1$ non-identity elements in $E_{p^n}$, and each subgroup of order $p$ has $p - 1$ generators (the non-identity elements of the subgroup), the number of subgroups of order $p$ in $E_{p^n}$ is
    \[\frac{p^n - 1}{p - 1} = 1 + p + p^2 + \cdots + p^{n-1}\]
    as desired.
\end{sol}

\begin{specialexercise}\label{ex5.1.12}
    Let $A$ and $B$ be groups. Assume $Z(A)$ contains a subgroup $Z_1$ and $Z(B)$ contains a subgroup $Z_2$ with $Z_1 \cong Z_2$. Let this isomorphism be given by the map $x_i \mapsto y_i$ for all $x_i \in Z_1$. A \textbf{\textit{central product}} of $A$ and $B$ is a quotient
    \[(A \times B)/Z \quad \text{where} \quad Z = \{(x_i, y_i^{-1}) \mid x_i \in Z_1\}\]
    and is denoted by $A * B$ --- it is not unique since it depends on $Z_1$, $Z_2$ and the isomorphism between them. (Think of $A * B$ as the direct product of $A$ and $B$ ``collapsed'' by identifying each element $x_i \in Z_1$ with its corresponding element $y_i \in Z_2$.)
    \begin{subproblems}
        \item Prove that the images of $A$ and $B$ in the quotient group $A * B$ are isomorphic to $A$ and $B$, respectively, and that these images intersect in a central subgroup isomorphic to $Z_1$. Find $|A * B|$.
        \item Let $Z_4 = \langle x \rangle$. Let $D_8 = \langle r, s \rangle$ and $Q_8 = \langle i, j \rangle$ be given by their usual generators and relations. Let $Z_4 * D_8$ be the central product of $Z_4$ and $D_8$ which identifies $x^2$ and $r^2$ (i.e., $Z_1 = \langle x^2 \rangle$, $Z_2 = \langle r^2 \rangle$ and the isomorphism is $x^2 \mapsto r^2$) and let $Z_4 * Q_8$ be the central product of $Z_4$ and $Q_8$ which identifies $x^2$ and $-1$. Prove that $Z_4 * D_8 \cong Z_4 * Q_8$.
    \end{subproblems}
\end{specialexercise}

\begin{solalph}
    \item Consider the maps $\pi_A : A \to A * B$ and $\pi_B : B \to A * B$ defined by $\pi_A(a) = (a, 1)Z$ and $\pi_B(b) = (1, b)Z$. It is clear that $\pi_A$ and $\pi_B$ are homomorphisms. Moreover, $\ker\pi_A$ consists of the elements $a \in A$ such that $(a, 1) \in Z$. Then $y_i = 1$, and since $x_i \mapsto y_i$ is an isomorphism, we have $x_i = 1$, so $a = 1$. Thus, $\ker\pi_A$ is trivial, and $\pi_A$ is injective. Similarly, $\ker\pi_B$ is trivial, and $\pi_B$ is injective. Therefore, the images of $A$ and $B$ in $A * B$ are isomorphic to $A$ and $B$, respectively. Moreover, the images of $A$ and $B$ in $A * B$ intersect in the subgroup $\{(x_i, 1)Z \mid x_i \in Z_1\}$, which is isomorphic to $Z_1$. Finally, since $Z$ has $|Z_1|$ elements, we have
    \[|A * B| = \frac{|A||B|}{|Z_1|}.\]
    \item Let $Z_D$ denote the $Z$ in the quotient for $Z_4 * D_8$ and let $Z_Q$ denote the $Z$ in the quotient for $Z_4 * Q_8$. Let $x$ generate $Z_4$ in $Z_4 * D_8$ and $y$ generate $Z_4$ in $Z_4 * Q_8$. Let $r$ and $s$ be the generators of $D_8$ in $Z_4 * D_8$ and let $i$ and $j$ be the generators of $Q_8$ in $Z_4 * Q_8$. Let $\b x$ denote the coset $(x, 1)Z_D$ in $Z_4 * D_8$, and use similar notation for the elements $r, s, y, i, j$. Let $\b z = \b j\b y\inv$. Consider the map $\phi : Z_4 * D_8 \to Z_4 * Q_8$ defined by 
    \[\phi(\b x^\alpha\b r^\beta\b s^\gamma) = \b y^\alpha\b i^\beta\b z^\gamma.\]
    To see that this is well-defined, suppose $\b x^\alpha\b r^\beta\b s^\gamma = \b 1$ in $Z_4 * D_8$. Then $\b r^\beta\b s^\gamma$ is in the image of $Z_4$ in $Z_4 * D_8$, so $\b r^\beta\b s^\gamma = \b x^{2\delta}$ for some $\delta$. Since $r^2$ is central in $D_8$, we have $\b r^\beta\b s^\gamma = \b r^{\beta + 2\delta}\b s^\gamma$. If $\gamma = 0$, then $\b r^{\beta + 2\delta} = \b 1$, so $4 \mid \beta + 2\delta$. If $\gamma = 1$, then $\b r^{\beta + 2\delta}\b s = \b 1$, so $4 \mid \beta + 2\delta - 2$. In either case, we have $4 \mid \beta + 2\delta$, so $\b r^\beta\b s^\gamma = \b x^{2\delta} = \b 1$. Thus, $\alpha$ is even, and we have $\phi(\b x^\alpha\b r^\beta\b s^\gamma) = \b y^\alpha\b i^\beta\b z^\gamma = \b 1$. Therefore, $\phi$ is well-defined. It is clear that $\phi$ is a homomorphism. Moreover, if $\phi(\b x^\alpha\b r^\beta\b s^\gamma) = \b 1$, then $\alpha$ is even and $\beta$ and $\gamma$ are zero, so $\phi$ is injective. Finally, since $|Z_4 * D_8| = |Z_4 * Q_8|$, we conclude that $\phi$ is an isomorphism.
\end{solalph}

\begin{exercise}\label{ex5.1.13}
    Give presentations for the groups $Z_4 * D_8$ and $Z_4 * Q_8$ constructed in the preceding exercise.
\end{exercise}

\begin{sol}
    To construct the presentations, we note that $\b x$ should satisfy the relations of $Z_4$, $\b r$ and $\b s$ should satisfy the relations of $D_8$, and $\b i$ and $\b z$ should satisfy the relations of $Q_8$. Moreover, we have the relations $\b x^2 = \b r^2$ in $Z_4 * D_8$, since $\b x^2\b r^{-2} = \b 1$. Lastly, $\b x$ is central in that $\b x \b r = \b r\b x$ and $\b x\b s = \b s\b x$ in $Z_4 * D_8$. Then a presentation for $Z_4 * D_8$ is
    \[\gen{\b x, \b r, \b s \mid \b r^4 = \b s^2 = \b 1, \b x^2 = \b r^2, \b s\b r\b s = \b r\inv, \b x\b r = \b r\b x, \b x\b s = \b s\b x}.\]
    To find a presentation for $Z_4 * Q_8$, we may use the isomorphism $\phi$ from the preceding exercise to rewrite the presentation for $Z_4 * D_8$ in terms of $\b y, \b i, \b z$. In particular, we have $\b y^2 = \b i^2\b z^2$, $\b z^2 = \b 1$, and $\b y\b i = \b i\b y\b z$. Thus, a presentation for $Z_4 * Q_8$ is
    \[\gen{\b y, \b i, \b z \mid \b i^4 = \b z^2 = \b 1, \b y^2 = \b i^2, \b z\b i \b z = \b i\inv, \b y\b i = \b i\b y, \b y\b z = \b z\b y}. \qh\]
\end{sol}

\begin{exercise}\label{ex5.1.14}
    Let $G = A_1 \times A_2 \times \cdots \times A_n$, and for each $i$, let $B_i$ be a normal subgroup of $A_i$. Prove that $B_1 \times B_2 \times \cdots \times B_n \nsub G$, and that
    \[(A_1 \times A_2 \times \cdots \times A_n)/(B_1 \times B_2 \times \cdots \times B_n) \cong (A_1/B_1) \times (A_2/B_2) \times \cdots \times (A_n/B_n).\]
\end{exercise}

\begin{sol}
    Let $G = A_1 \times A_2 \times \cdots \times A_n$ and let $B_i \nsub A_i$ for all $i$. Let $H = B_1 \times B_2 \times \cdots \times B_n$. It is clear that $H$ is a subgroup of $G$. Moreover, since each $B_i$ is normal in $A_i$, we have that for any $g = (a_1, a_2, \ldots, a_n) \in G$ and any $h = (b_1, b_2, \ldots, b_n) \in H$, we have
    \[ghg\inv = (a_1b_1a_1\inv, a_2b_2a_2\inv, \ldots, a_nb_na_n\inv) \in H\]
    since $a_ib_ia_i\inv \in B_i$ for all $i$. Thus, $H$ is normal in $G$.

    Consider the map $\phi : G \to (A_1/B_1) \times (A_2/B_2) \times \cdots \times (A_n/B_n)$ defined by
    \[\phi((a_1, a_2, \ldots, a_n)) = (a_1B_1, a_2B_2, \ldots, a_nB_n).\]
    It is clear that $\phi$ is a homomorphism. Moreover, $\phi$ is surjective since for any $(a_1B_1, a_2B_2, \ldots, a_nB_n)$, we have $\phi((a_1, a_2, \ldots, a_n)) = (a_1B_1, a_2B_2, \ldots, a_nB_n)$. Finally, the kernel of $\phi$ is the set of all $(a_1, a_2, \ldots, a_n)$ such that $a_i \in B_i$ for all $i$, which is exactly $H$. Thus, by the First Isomorphism Theorem, we have
    \[G/H \cong (A_1/B_1) \times (A_2/B_2) \times \cdots \times (A_n/B_n)\]
    as desired.
\end{sol}

\begin{exercise}\label{ex5.1.15}
    Let $I$ be any nonempty index set and let $(G_i, \star_i)$ be a group for each $i \in I$. The \textit{direct product} of the groups $G_i$, $i \in I$ is the set $G = \prod_{i \in I} G_i$ (the Cartesian product of the $G_i$'s) with a binary operation defined as follows: if $\prod a_i$ and $\prod b_i$ are elements of $G$, then their product in $G$ is given by
    \[\left(\prod_{i \in I} a_i\right)\left(\prod_{i \in I} b_i\right) = \prod_{i \in I} (a_i \star_i b_i)\]
    (i.e., the group operation in the direct product is defined componentwise).
    \begin{subproblems}
        \item Show that this binary operation is well defined and associative.
        \item Show that the element $\prod 1_i$ satisfies the axiom for the identity of $G$, where $1_i$ is the identity of $G_i$ for all $i$.
        \item Show that the element $\prod a_i^{-1}$ is the inverse of $\prod a_i$, where the inverse of each component element $a_i$ is taken in the group $G_i$.
    \end{subproblems}
    Conclude that the direct product is a group. (Note that if $I = \set{1, 2, \ldots, n}$, this definition of the direct product is the same as the $n$-tuple definition in the text.)
\end{exercise}

\begin{solalph}
    \item To show well-definedness, let $\prod \alpha_i = \prod a_i$ and $\prod \beta_i = \prod b_i$. Then $\alpha_i = a_i$ and $\beta_i = b_i$ for all $i$, so
    \[\left(\prod_{i \in I} \alpha_i\right)\left(\prod_{i \in I} \beta_i\right) = \prod_{i \in I} (\alpha_i \star_i \beta_i) = \prod_{i \in I} (a_i \star_i b_i)\]
    as desired. To show associativity, let $\prod a_i$, $\prod b_i$, and $\prod c_i$ be elements of $G$. Then
    \begin{align*}
        \left(\prod_{i \in I} a_i\right)\left(\left(\prod_{i \in I} b_i\right)\left(\prod_{i \in I} c_i\right)\right) & = \left(\prod_{i \in I} a_i\right)\left(\prod_{i \in I} (b_i \star_i c_i)\right) \\
        & = \prod_{i \in I} (a_i \star_i (b_i \star_i c_i)) \\
        & = \prod_{i \in I} ((a_i \star_i b_i) \star_i c_i) \\
        & = \left(\prod_{i \in I} (a_i \star_i b_i)\right)\left(\prod_{i \in I} c_i\right) \\
        & = \left(\left(\prod_{i \in I} a_i\right)\left(\prod_{i \in I} b_i\right)\right)\left(\prod_{i \in I} c_i\right)
    \end{align*}
    as desired.
    \item Let $\prod a_i$ be an element of $G$. Then
    \[\left(\prod_{i \in I} 1_i\right)\left(\prod_{i \in I} a_i\right) = \prod_{i \in I} (1_i \star_i a_i) = \prod_{i \in I} a_i\]
    and
    \[\left(\prod_{i \in I} a_i\right)\left(\prod_{i \in I} 1_i\right) = \prod_{i \in I} (a_i \star_i 1_i) = \prod_{i \in I} a_i.\]
    Thus, $\prod 1_i$ is the identity element of $G$.
    \item Let $\prod a_i$ be an element of $G$. Then
    \[\left(\prod_{i \in I} a_i\right)\left(\prod_{i \in I} a_i^{-1}\right) = \prod_{i \in I} (a_i \star_i a_i^{-1}) = \prod_{i \in I} 1_i\]
    and
    \[\left(\prod_{i \in I} a_i^{-1}\right)\left(\prod_{i \in I} a_i\right) = \prod_{i \in I} (a_i^{-1} \star_i a_i) = \prod_{i \in I} 1_i.\]
    Thus, $\prod a_i^{-1}$ is the inverse of $\prod a_i$. We conclude that $G$ is a group with inverse and identity as described above.
\end{solalph}

\begin{exercise}\label{ex5.1.16}
    State and prove the generalization of Proposition 2 to arbitrary direct products.
\end{exercise}

\begin{sol}
    Proposition 2 stated for direct arbitrary groups is as follows: if $G = \prod_{i \in I} G_i$ is a direct product of groups, then
    \begin{enumerate}
        \item Fix each $i$ with $j$-th coordinate where $j \neq i$ is the identity element of $G_j$ for $1 \leq j \leq n$. Then
        \[\wh G_i = \set{(1, 1, \ldots, 1, g_i, 1, \ldots) \mid g_i \in G_i}\] is a normal subgroup of $G$ that is isomorphic to $G_i$, or that
        \[G/\wh G_i \cong \prod_{j \neq i} G_j.\]
        \item For fixed $i$, define $\pi_i : G \to G_i$ by $\pi_i((g_j)_{j \in I}) = g_i$. Then $\pi_i$ is a surjective homomorphism with kernel $\wh G_i$, so $G/\wh G_i \cong G_i$.
        \item If $x \in G_i$ and $y \in G_j$ for $i \neq j$, then $x$ and $y$ commute in $G$.
    \end{enumerate}
    The proof is similar to the proof of Proposition 2 in the text.
    \begin{enumerate}
        \item It is clear that $\wh G_i$ is a subgroup of $G$. Moreover, for any $\prod_{j \in I} g_j \in G$ and any $h = \prod_{j \neq i} 1_j h_i \in \wh G_i$, we have
        \[\left(\prod_{j \in I} g_j\right)h\left(\prod_{j \in I} g_j\right)\inv = \prod_{j \neq i} 1_j (g_i h_i g_i\inv) \in \wh G_i\]
        since $g_i h_i g_i\inv \in G_i$. Thus, $\wh G_i$ is normal in $G$. Consider the map $\phi : G_i \to \wh G_i$ defined by $\phi(g_i) = \prod_{j \neq i} 1_j g_i$. It is clear that $\phi$ is a homomorphism. Moreover, if $\phi(g_i) = \prod_{j \neq i} 1_j g_i = \prod_{j \neq i} 1_j$, then $g_i$ is the identity element of $G_i$, so $\phi$ is injective. Finally, for any $\prod_{j \neq i} 1_j h_i \in \wh G_i$, we have $\phi(h_i) = \prod_{j \neq i} 1_j h_i$, so $\phi$ is surjective. Thus, $\wh G_i$ is isomorphic to $G_i$. 
        
        Consider the map
        \[\psi : G \to \prod_{j \neq i} G_j \text{ defined by } \psi\left(\prod_{j \in I} g_j\right) = \prod_{j \neq i} g_j.\]
        It is clear that $\psi$ is a homomorphism. Moreover, the kernel of $\psi$ is $\wh G_i$, so by the First Isomorphism Theorem, we have
        \[G/\wh G_i \cong \prod_{j \neq i} G_j.\]
        \item For fixed $i$, it is clear that $\pi_i$ is a homomorphism. Moreover, $\pi_i$ is surjective since for any $g_i \in G_i$, we have $\pi_i(\prod_{j \in I} 1_j g_i) = g_i$. Finally, the kernel of $\pi_i$ is $\wh G_i$, so by the First Isomorphism Theorem, we have
        \[G/\wh G_i \cong G_i.\]
        \item Let $x = \prod_{j \neq i} 1_j x_i$ and $y = \prod_{j \neq k} 1_j y_k$ for some $i \neq k$. Then
        \[xy = \prod_{j \neq i, k} 1_j x_i y_k = \prod_{j \neq i, k} 1_j y_k x_i = yx\]
        as desired.
    \end{enumerate}
\end{sol}

\begin{exercise}\label{ex5.1.17}
    Let $I$ be any nonempty index set and let $G_i$ be a group for each $i \in I$. The \textit{restricted direct product} or \textit{direct sum} of the groups $G_i$ is the set of elements of the direct product which are the identity in all but finitely many components, that is, the set of all elements $\prod a_i \in \prod_{i \in I} G_i$ such that $a_i = 1_i$ for all but a finite number of $i \in I$.
    \begin{subproblems}
        \item Prove that the restricted direct product is a subgroup of the direct product.
        \item Prove that the restricted direct product is normal in the direct product.
        \item Let $I = \mathbb{Z}^+$ and let $p_i$ be the $i$th integer prime. Show that if $G_i = \mathbb{Z}/p_i\mathbb{Z}$ for all $i \in \mathbb{Z}^+$, then every element of the restricted direct product of the $G_i$'s has finite order but $\prod_{i \in \mathbb{Z}^+} G_i$ has elements of infinite order. Show that in this example the restricted direct product is the torsion subgroup of the direct product (cf. \hyperref[ex2.1.16]{Exercise 2.1.16}).
    \end{subproblems}
\end{exercise}

\begin{solalph}
    \item Let $G = \prod_{i \in I} G_i$ and let $H$ be the restricted direct product of the $G_i$'s. It is clear that $H$ is a subset of $G$. Moreover, if $\prod a_i$ and $\prod b_i$ are elements of $H$, then there are only finitely many $i$ such that $a_i \neq 1_i$ and only finitely many $i$ such that $b_i \neq 1_i$. Thus, there are only finitely many $i$ such that $a_i b_i \neq 1_i$, so $\prod a_i \prod b_i = \prod (a_i b_i)$ is an element of $H$. Finally, if $\prod a_i$ is an element of $H$, then there are only finitely many $i$ such that $a_i \neq 1_i$, so there are only finitely many $i$ such that $a_i^{-1} \neq 1_i$, so $\prod a_i^{-1}$ is an element of $H$. Thus, $H \leq G$.
    \item Let $G = \prod_{i \in I} G_i$ and let $H$ be the restricted direct product of the $G_i$'s. Let $\prod a_i$ be an element of $G$ and let $\prod b_i$ be an element of $H$. Then there are only finitely many $i$ such that $b_i \neq 1_i$, so there are only finitely many $i$ such that $a_i b_i a_i^{-1} \neq 1_i$, so $\prod a_i \prod b_i \prod a_i^{-1}$ is an element of $H$. Thus, $H \nsub G$.
    \item Let $G = \prod_{i \in \mathbb{Z}^+} G_i$ where $G_i = \mathbb{Z}/p_i\mathbb{Z}$ for all $i$. Let $H$ be the restricted direct product of the $G_i$'s. If $\prod a_i$ is an element of $H$, then there are only finitely many $i$ such that $a_i \neq 1_i$, so there are only finitely many $i$ such that $a_i$ has order greater than 1, so $\prod a_i$ has finite order. On the other hand, if $\prod a_i$ is an element of $G$ such that $a_i$ is a generator of $G_i$ for all $i$, then $\prod a_i$ has infinite order since for any positive integer $n$, we have $(\prod a_i)^n = \prod a_i^n$, and there are infinitely many $i$ such that $a_i^n \neq 1_i$.
    
    To show $H = \tor(G)$, we note that $H \subseteq \tor(G)$ since every element of $H$ has finite order. Conversely, if $\prod a_i$ is an element of $\tor(G)$, then there is some positive integer $n$ such that $(\prod a_i)^n = \prod a_i^n = \prod 1_i$, so $a_i^n = 1_i$ for all $i$. Since $G_i$ is cyclic of prime order, this implies that $a_i$ is either the identity or a generator of $G_i$. If there are infinitely many $i$ such that $a_i$ is a generator of $G_i$, then $\prod a_i$ has infinite order, which is a contradiction. Thus, there are only finitely many $i$ such that $a_i$ is a generator of $G_i$, so there are only finitely many $i$ such that $a_i \neq 1_i$, so $\prod a_i$ is an element of $H$. Therefore, $\tor(G) \subseteq H$, and we conclude that $\tor(G) = H$.
\end{solalph}

\begin{exercise}\label{ex5.1.18}
    In each of (a) to (e) give an example of a group with the specified properties:
    \begin{subproblems}
        \item an infinite group in which every element has order 1 or 2
        \item an infinite group in which every element has finite order but for each positive integer $n$ there is an element of order $n$
        \item a group with an element of infinite order and an element of order 2
        \item a group $G$ such that every finite group is isomorphic to some subgroup of $G$
        \item a non-trivial group $G$ such that $G \cong G \times G$.
    \end{subproblems}
\end{exercise}

\begin{solalph}
    \item Consider the infinite direct product $\intmod[2] \times \intmod[2] \times \cdots$. Every element of this group is a tuple of 0's and 1's, and the order of each element is either 1 (if the tuple is the identity) or 2 (if the tuple has at least one 1).
    \item Let $G = \intmod[2] \times \intmod[3] \times \intmod[4] \times \cdots$. Every element of $G$ has finite order since it is a tuple of elements from finite cyclic groups. Moreover, for each positive integer $n$, the element that is the identity in all components except the $n$-th component, which is a generator of $\intmod[n]$, has order $n$.
    \item Consider the group $\intmod[2] \times \mathbb{Z}$. The element $(1, 0)$ has order 2, and the element $(0, 1)$ has infinite order.
    \item By Cayley's Theorem, every finite group is isomorphic to a subgroup of some symmetric group $S_n$. We may then consider the group $G = S_1 \times S_2 \times S_3 \times \cdots$. For any finite group $H$, there is some $n$ such that $H$ is isomorphic to a subgroup of $S_n$, so $H$ is isomorphic to a subgroup of $G$.
    \item Let $G = \z \times \z \times \cdots$. Consider the map $\phi : G \to G \times G$ defined by
    \[\phi((a_1, a_2, a_3, \ldots)) = ((a_1, a_3, a_5, \ldots), (a_2, a_4, a_6, \ldots)).\]
    It is clear that $\phi$ is a homomorphism. Moreover, if $\phi((a_1, a_2, a_3, \ldots)) = ((0, 0, 0, \ldots), (0, 0, 0, \ldots))$, then $a_i = 0$ for all $i$, so $\phi$ is injective. Finally, let $((b_1, b_2, b_3, \ldots), (c_1, c_2, c_3, \ldots))$ be an element of $G \times G$. Then 
    \[\phi((b_1, c_1, b_2, c_2, b_3, c_3, \ldots)) = ((b_1, b_2, b_3, \ldots), (c_1, c_2, c_3, \ldots)),\]
    so $\phi$ is surjective. Thus, $\phi$ is an isomorphism and $G \cong G \times G$.
\end{solalph}

\newpage

\subsection{The Fundamental Theorem of Finitely Generated Abelian Groups}

\begin{exercise}\label{ex5.2.1}
    In each parts (a) to (e), give the number of non-isomorphic abelian groups of the specified order---do not list the groups:

    \noindent (a) order 100, \quad (b) order 576, \quad (c) order 1155, \quad (d) order 42875, \quad (e) order 2704.
\end{exercise}

\begin{solalph}
    \item $100 = 2^2 \cdot 5^2$; 2 partitions of 2, hence $2 \cdot 2 = 4$ non-isomorphic abelian groups of order 100.
    \item $576 = 2^6 \cdot 3^2$; 11 partitions of 6, hence $11 \cdot 2 = 22$ non-isomorphic abelian groups of order 576.
    \item $1155 = 3 \cdot 5 \cdot 7 \cdot 11$. Only one abelian group of order 1155.
    \item $42875 = 5^3 \cdot 7^3$; 3 partitions of 3, hence $3 \cdot 3 = 9$ non-isomorphic abelian groups of order 42875.
    \item $2704 = 2^4 \cdot 13^2$; 5 partitions of 4, hence $5 \cdot 2 = 10$ non-isomorphic abelian groups of order 2704.
\end{solalph}

\begin{exercise}\label{ex5.2.2}
    In each of parts (a) to (e), give the lists of invariant factors for all abelian groups of the specified order:

    \noindent (a) group 270, \quad (b) order 9801, \quad (c) order 320, \quad (d) order 105, \quad (e) order 44100.
\end{exercise}

\begin{solalph}
    \item $270 = 2 \cdot 3^3 \cdot 5$. The invariant factors are $\set{270}, \set{90, 3}, \set{30, 3, 3}$.
    \item $9801 = 3^4 \cdot 11^2$. The invariant factors are $\set{9801}, \set{891, 11}, \set{3267, 3}, \set{297, 33}, \set{1089, 9}, \set{99, 99}, \\
    \set{1089, 3, 3}, \set{99, 33, 3}, \set{363, 3, 3, 3}, \set{33, 33, 3, 3}$.
    \item $320 = 2^6 \cdot 5$. The invariant factors are $\set{320}, \set{160, 2}, \set{80, 4}, \set{80, 2, 2}, \set{40, 8}, \set{40, 4, 2}, \set{40, 2, 2, 2}, \\
    \set{20, 4, 4}, \set{20, 4, 2, 2}, \set{20, 2, 2, 2, 2}, \set{10, 2, 2, 2, 2, 2}$.
    \item $105 = 3 \cdot 5 \cdot 7$. The invariant factors are $\set{105}$.
    \item $44100 = 2^2 \cdot 3^2 \cdot 5^2 \cdot 7^2$. The invariant factors are $\set{44100}, \set{22050, 2}, \set{14700, 3}, \set{8820, 5}, \set{6300, 7}, \\
    \set{4410, 10}, \set{3150, 14}, \set{2205, 15}, \set{1800, 21}, \set{1260, 35}, \set{1470, 30}, \set{1050, 42}, \set{630, 70}, \set{420, 105},\\
    \set{210, 210}$.
\end{solalph}

\begin{exercise}\label{ex5.2.3}
    In each of the parts (a) to (e), give the lists of elementary divisors for all abelian groups of the specified order, and then match each list with the corresponding list of invariant factors found in the preceding exercise:

    \noindent (a) group 270, \quad (b) order 9801, \quad (c) order 320, \quad (d) order 105, \quad (e) order 44100.
\end{exercise}

\begin{solalph}
    \item $\set{2, 27, 5}, \set{2, 9, 5, 3}, \set{2, 3, 5, 3, 3}$.
    \item $\set{81, 121}, \set{81, 11, 11}, \set{27, 121, 3}, \set{27, 11, 3, 11}, \set{9, 121, 9}, \set{9, 11, 9, 11}, \set{9, 121, 3, 3}, \set{9, 11, 3, 11, 3}, \\
    \set{3, 121, 3, 3, 3}, \set{3, 11, 3, 11, 3, 3}$.
    \item $\set{64, 5}, \set{32, 5, 2}, \set{16, 5, 4}, \set{16, 5, 2, 2}, \set{8, 5, 8}, \set{8, 5, 4, 2}, \set{8, 5, 2, 2, 2}, \set{4, 5, 4, 4}, \set{4, 5, 4, 2, 2}, \\
    \set{4, 5, 2, 2, 2, 2}, \set{2, 5, 2, 2, 2, 2, 2}$.
    \item $\set{3, 5, 7}$.
    \item $\set{4, 9, 25, 49}, \set{2, 9, 25, 49, 2}, \set{4, 3, 25, 49, 3}, \set{4, 9, 5, 49, 5}, \set{4, 9, 25, 7, 7}, \set{2, 3, 25, 49, 2, 3}, \set{2, 9, 5, 49, 2, 5}, \\
    \set{2, 9, 25, 7, 2, 7}, \set{4, 3, 5, 49, 3, 5}, \set{2, 3, 25, 49, 2, 3}, \set{2, 9, 5, 49, 2, 7}, \set{4, 3, 25, 7, 3, 7}, \set{2, 3, 5, 49, 2, 3, 5}, \\
    \set{2, 3, 25, 7, 2, 3, 7}, \set{4, 3, 5, 7, 3, 5, 7}, \set{2, 3, 5, 7, 2, 3, 5, 7}$.
\end{solalph}

\begin{exercise}\label{ex5.2.4}
    In each of parts (a) to (d), determine which pairs of abelian groups listed are isomorphic (here the expression $\set{a_1, a_2, \ldots, a_k}$ denotes the abelian group $Z_{a_1} \times Z_{a_2} \times \cdots \times Z_{a_k}$).
    \begin{subproblems}
        \item $\{4, 9\}$, \quad $\{6, 6\}$, \quad $\{8, 3\}$, \quad $\{9, 4\}$, \quad $\{6, 4\}$, \quad $\{64\}$.
        \item $\{2^2, 2 \cdot 3^2\}$, \quad $\{2^2 \cdot 3, 2 \cdot 3\}$, \quad $\{2^3 \cdot 3^2\}$, \quad $\{2^2 \cdot 3^2, 2\}$.
        \item $\{5^2 \cdot 7^2,\ 3^2 \cdot 5 \cdot 7\}$, \quad $\{3^2 \cdot 5^2 \cdot 7,\ 5 \cdot 7^2\}$, \quad $\{3 \cdot 5^2,\ 7^2,\ 3 \cdot 5 \cdot 7\}$,
        \quad $\{5^2 \cdot 7,\ 3^2 \cdot 5,\ 7^2\}$.
        \item $\{2^2 \cdot 5 \cdot 7,\ 2^3 \cdot 5^3,\ 2 \cdot 5^2\}$, \quad $\{2^3 \cdot 5^3 \cdot 7,\ 2^3 \cdot 5^3\}$, \quad $\{2^2,\ 2 \cdot 7,\ 2^3,\ 5^3,\ 5^3\}, \quad \{2 \cdot 5^3,\ 2^2 \cdot 5^3,\ 2^3,\ 7\}$.
    \end{subproblems}
\end{exercise}

\begin{solalph}
    \item $\set{4, 9}$ and $\set{9, 4}$ have the same elementary divisors, so they are isomorphic. In order, the other four groups have elementary divisors $\set{3, 2, 3, 2}$, $\set{8, 3}$, $\set{3, 2, 4}$, and $\set{64}$, so they are not isomorphic to each other or to $\set{4, 9}$.
    \item $\set{2^2, 2 \cdot 3^2}$ and $\set{2^2 \cdot 3^2, 2}$ have the same elementary divisors, so they are isomorphic. In order, the other two groups have elementary divisors $\set{2^2 \cdot 3, 2 \cdot 3}$ and $\set{2^3 \cdot 3^2}$, so they are not isomorphic to each other or to $\set{2^2, 2 \cdot 3^2}$.
    \item The first, second, and fourth groups have the same elementary divisors $\set{5^2, 7^2, 3^2, 5, 7}$, so they are isomorphic. The third group has elementary divisors $\set{3, 5^2, 7^2, 3, 5, 7}$, so it is not isomorphic to the other three groups.
    \item The third and fourth groups have the elementary divisors $\set{2, 125, 4, 125, 8, 7}$, so they are isomorphic.
\end{solalph}

\begin{exercise}\label{ex5.2.5}
    Let $G$ be a finite abelian group of type $(n_1, n_2, \ldots, n_t)$. Prove that $G$ contains an element of order $m$ if and only if $m \mid n_1$. Deduce that $G$ is of exponent $n_1$.
\end{exercise}

\begin{sol}
    \leavevmode
    \begin{ifandonlyif}
        \item[\rightimp] Let $g = (g_1, g_2, \ldots, g_t) \in G$ have order $m$. Then $m = \lcm(|g_1|, |g_2|, \ldots, |g_t|)$, so $|g_i| \mid m$ for all $i$. Since $G$ is of type $(n_1, n_2, \ldots, n_t)$, we have $|g_i| \mid n_i$ for all $i$, so $m$ divides $\lcm(n_1, n_2, \ldots, n_t) = n_1$.
        \item[\leftimp] If $m \mid n_1$, then there exists an element of order $m$ in the cyclic subgroup of order $n_1$.
    \end{ifandonlyif}
    Since $G$ contains an element of order $n_1$, the exponent of $G$ is at least $n_1$. On the other hand, if $g \in G$, then $|g| \mid n_1$, so the exponent of $G$ divides $n_1$. Thus, the exponent of $G$ is exactly $n_1$.
\end{sol}

\begin{exercise}\label{ex5.2.6}
    Prove that any finite group has a finite exponent. Give an example of an infinite group with finite exponent. Does a finite group of exponent $m$ always contain an element of order $m$?
\end{exercise}

\begin{sol}
    Let $G$ be a finite group. Then there are only finitely many elements of $G$, so there are only finitely many orders of elements in $G$. Let $m$ be the least common multiple of the orders of all elements of $G$. Then for any $g \in G$, we have $g^m = e$, so $G$ has exponent dividing $m$.
    
    An example of an infinite group with finite exponent is the infinite direct product $\intmod[2] \times \intmod[2] \times \cdots$. Every element of this group has order 1 or 2, so the exponent of this group is 2.
    
    A finite group of exponent $m$ does not always contain an element of order $m$: consider $S_3$, which has exponent 6 but does not contain an element of order 6.
\end{sol}

\begin{exercise}\label{ex5.2.7}
    Let $p$ be a prime and let $A = \langle x_1 \rangle \times \langle x_2 \rangle \times \cdots \times \langle x_n \rangle$ be an abelian $p$-group, where $|x_i| = p^{\alpha_i} > 1$ for all $i$. Define the $p^{\text{th}}$-\textit{power map}
    \[\phi : A \to A \quad \text{by} \quad \phi : x \mapsto x^p.\]
    \begin{subproblems}
        \item Prove that $\phi$ is a homomorphism.
        \item Describe the image and kernel of $\phi$ in terms of the given generators.
        \item Prove both $\ker \phi$ and $A / \operatorname{im} \phi$ have rank $n$ (i.e., have the same rank as $A$) and prove these groups are both isomorphic to the elementary abelian group, $E_{p^n}$, of order $p^n$.
    \end{subproblems}
\end{exercise}

\begin{solalph}
    \item Let $x = (x_1^{s_1}, x_2^{s_2}, \ldots, x_n^{s_n})$ and $y = (x_1^{t_1}, x_2^{t_2}, \ldots, x_n^{t_n})$ be elements of $A$. Then
    \begin{align*}
        \phi(xy) & = \phi((x_1^{s_1 + t_1}, x_2^{s_2 + t_2}, \ldots, x_n^{s_n + t_n})) \\
        & = (x_1^{p(s_1 + t_1)}, x_2^{p(s_2 + t_2)}, \ldots, x_n^{p(s_n + t_n)}) \\
        & = (x_1^{ps_1} x_1^{pt_1}, x_2^{ps_2} x_2^{pt_2}, \ldots, x_n^{ps_n} x_n^{pt_n}) \\
        & = (x_1^{ps_1}, x_2^{ps_2}, \ldots, x_n^{ps_n})(x_1^{pt_1}, x_2^{pt_2}, \ldots, x_n^{pt_n}) \\
        & = \phi(x)\phi(y)
    \end{align*}
    as desired.
    \item $\phi$ collects all the $p$-powers, so $\im\phi = \langle x_1^p \rangle \times \langle x_2^p \rangle \times \cdots \times \langle x_n^p \rangle$. Moreover, $\ker\phi$ consists of all elements $(x_1^{s_1}, x_2^{s_2}, \ldots, x_n^{s_n})$ such that $x_1^{ps_1} = x_2^{ps_2} = \cdots = x_n^{ps_n} = e$, which implies that $ps_i$ is a multiple of $p^{\alpha_i}$ for all $i$, so $s_i$ is a multiple of $p^{\alpha_i - 1}$ for all $i$. Thus,
    \[\ker\phi = \left\langle x_1^{p^{\alpha_1-1}} \right\rangle \times \left\langle x_2^{p^{\alpha_2-1}} \right\rangle \times \cdots \times \left\langle x_n^{p^{\alpha_n-1}} \right\rangle.\]
    \item Note that $x_i^{p^{\alpha_i - 1}}$ has order $p$ for all $i$, so $\ker\phi$ is an elementary abelian group of order $p^n$. Moreover, note that each $\gen{x_i}$ is cyclic, hence every group is normal. Examining $\gen{x_i} / \gen{x_i^p}$, we see that $\gen{x_i} / \gen{x_i^p}$ is an elementary abelian group of order $p$. Since $A / \im\phi$ is the direct product of the $\gen{x_i} / \gen{x_i^p}$, we conclude that $A / \im\phi$ is an elementary abelian group of order $p^n$, by \hyperref[ex5.1.14]{Exercise 5.1.14}.
\end{solalph}

\begin{exercise}\label{ex5.2.8}
    Let $A$ be a finite abelian group (written multiplicatively) and let $p$ be a prime. Let
    \[A^p = \{a^p \mid a \in A\} \qquad \text{and} \qquad A_p = \{x \mid x^p = 1\}\]
    (so $A^p$ and $A_p$ are the image and kernel of the $p^{\text{th}}$-power map, respectively).
    \begin{subproblems}
        \item Prove that $A/A^p \cong A_p$. [Show that they are both elementary abelian and they have the same order.]
        \item Prove that the number of subgroups of $A$ of order $p$ equals the number of subgroups of $A$ of index $p$. [Reduce to the case where $A$ is an elementary abelian $p$-group.]
    \end{subproblems}
\end{exercise}

\begin{solalph}
    \item By Theorem 5.5, we may write $A$ as a direct product:
    \[A \cong A_{p_1} \times A_{p_2} \times \cdots \times A_{p_k},\]
    where each $A_{p_i}$ is a $p_i$-group for distinct primes $p_1, p_2, \ldots, p_k$. There are two cases:
    \begin{itemize}
        \item If $p \neq p_i$ for all $i$, then then $(p, p_i) = 1$ for all $i$, so the $p$-th power map is an automorphism of each $A_{p_i}$, so $A^p = A$ and $A_p = \{1\}$, so $A/A^p \cong A_p$.
        \item If $p = p_i$ for some $i$, then it is an automorphism for all $A_j$ for $j \neq i$, while the case for $A_{p_i}$ is handled by \hyperref[ex5.2.7]{Exercise 5.2.7}. Thus,
        \[A^p = A_{p_1} \times \cdots \times A_{p_{i-1}} \times A_{p_i}^p \times A_{p_{i+1}} \times \cdots \times A_{p_k}\]
        and
        \[A_p = \set 1 \times \cdots \times \set 1 \times (A_{p_i})_p \times \set 1 \times \cdots \times \set 1.\]
        We may again use \hyperref[ex5.1.14]{Exercise 5.1.14} along with the fact that $A_{p_i} / A_{p_i}^p \cong A_{p_i}^p$ to conclude that $A/A^p \cong A_p$.
    \end{itemize}
    \item Let $H \leq A$ such that $|A : H| = p$. Let $a \in A$. Then $aH \in A/H$. Since $|A : H| = p$, then $|aH| = 1$ or $p$. If $|aH| = 1$, then $a \in H$. If $|aH| = p$, then $a^pH = H$, or $a^p \in H$. In either case, we have $a^p \in H$, so $A^p \subseteq H$. By the Lattice Isomorphism Theorem, then subgroups of $A$ containing $A^p$ are in bijection with subgroups of $A/A^p$. By the previous part, $A/A^p \cong A_p$, so the number of index $p$ subgroups of $A$ equals the number of index $p$ subgroups of $A_p$.
    
    Now let $K \leq A$ where $|K| = p$. For any $k \in K$, we know $k^p = 1$, so $K \subseteq A_p$. Then subgroups of order $p$ in $A$ are exactly the subgroups of order $p$ in $A_p$, since $A_p \leq A$. Thus, the number of subgroups of order $p$ in $A$ equals the number of subgroups of order $p$ in $A_p$.

    Note that $A_p$ is an elementary abelian $p$-group, so it is isomorphic to $E_{p^n}$ for some $n$. By \hyperref[ex5.1.11]{Exercise 5.1.11}, the number of subgroups of order $p$ in $E_{p^n}$ is $\frac{p^n - 1}{p - 1}$. To see that the number of subgroups of index $p$ in $E_{p^n}$ is also $\frac{p^n - 1}{p - 1}$, we note that there are $p^n - 1$ surjective homomorphisms from $E_{p^n}$ to $E_p$, and each subgroup of index $p$ is the kernel of exactly $p - 1$ such homomorphisms, so there are $\frac{p^n - 1}{p - 1}$ subgroups of index $p$ in $E_{p^n}$.
\end{solalph}

\begin{exercise}\label{ex5.2.9}
    Let $A = Z_{60} \times Z_{45} \times Z_{12} \times Z_{36}$. Find the number of elements of order 2 and the number of subgroups of index 2 in $A$.
\end{exercise}

\begin{sol}
    Note that
    \begin{itemize}
        \item $Z_{60} \cong Z_4 \times Z_3 \times Z_5$,
        \item $Z_{45} \cong Z_9 \times Z_5$,
        \item $Z_{12} \cong Z_4 \times Z_3$, and
        \item $Z_{36} \cong Z_4 \times Z_9$.
    \end{itemize}
    Thus, the 2-power part of $A$ is $Z_4 \times Z_4 \times Z_4$. Elements of order 2 in $A$ are elements of the 2-power part as well as trivial parts of the other groups. The 2-power map sends $(a, b, c)$ to $(2a, 2b, 2c)$, so the kernel of the 2-power map is $\set{0, 2} \times \set{0, 2} \times \set{0, 2}$, which has 8 elements, 7 of which have order 2. Thus, there are 7 elements of order 2 in $A$.
\end{sol}

\begin{exercise}\label{ex5.2.10}
    Let $n$ and $k$ be positive integers and let $A$ be the free abelian group of rank $n$ (written additively). Prove that $A/kA$ is isomorphic to the direct product of $n$ copies of $\intmod[k]$ (here, $kA = \set{ka \mid a \in A}$). [See \hyperref[ex5.1.14]{Exercise 5.1.14}.]
\end{exercise}

\begin{sol}
    It is easy to see that $kA = k\z \times \cdots \times k\z$. Since $\z$ is abelian, then $k\z \nsub \z$. By \hyperref[ex5.1.14]{Exercise 5.1.14}, we have $A/kA \cong \intmod[k] \times \cdots \times \intmod[k].$
\end{sol}

\newpage

\begin{exercise}\label{ex5.2.11}
    Let $G$ be a non-trivial, finite abelian group of rank $t$.
    \begin{subproblems}
        \item Prove that the rank of $G$ equals the maximum of the ranks of its Sylow subgroups.
        \item Prove that $G$ can be generated by $t$ elements, but no subset with fewer than $t$ elements generates $G$. (One way of doing this is by using part(a) together with \hyperref[ex5.2.7]{Exercise 5.2.7}.)
    \end{subproblems}
\end{exercise}

\begin{solalph}
    \item Let $G$ be of type $(n_1, n_2, \ldots, n_t)$, and consider the invariant factor decomposition of $G$:
    \[G \cong Z_{n_1} \times Z_{n_2} \times \cdots \times Z_{n_t}.\]
    Let $p$ be a prime dividing the smallest factor $n_t$. Then the Sylow $p$-subgroup of $G$ is isomorphic to $Z_{p^{\alpha_1}} \times Z_{p^{\alpha_2}} \times \cdots \times Z_{p^{\alpha_t}}$, where $\alpha_i$ is the largest power of $p$ dividing $n_i$. Since $\alpha_t > 0$, the Sylow $p$-subgroup has rank $t$. On the other hand, if $q$ is any prime, then for any $Q \in \syl_q(G)$, we have
    \[Q \cong Z_{q^{\beta_1}} \times Z_{q^{\beta_2}} \times \cdots \times Z_{q^{\beta_t}},\]
    where $\beta_i \mid n_i$ for all $i$ with some potentially trivial factors, so the rank of $Q$ is at most $t$. Thus, the rank of $G$ equals the maximum of the ranks of its Sylow subgroups.
    \item It is clear that $G$ is generated by $t$ elements. Suppose $G = \gen{g_1, g_2, \ldots, g_{t-1}}$ and some $P \in \syl_p(G)$, which has rank $t$. Projecting $G$ onto $P$, we see that $P = \gen{\pi(g_1), \pi(g_2), \ldots, \pi(g_{t-1})}$. However, letting $\phi$ be the $p$-power map on $P$ as defined in \hyperref[ex5.2.7]{Exercise 5.2.7}, we have that $P/\im\phi \cong E_{p^t}$, so $P/\im\phi$ has rank $t$. Since $\pi(g_1), \pi(g_2), \ldots, \pi(g_{t-1})$ generate $P/\im\phi$, we have a contradiction. Thus, no subset with fewer than $t$ elements generates $G$.
\end{solalph}

\begin{specialexercise}\label{ex5.2.12}
    Let $n$ and $m$ be positive integers with $d = (n, m)$. Let $Z_n = \gen x$ and $Z_m = \gen y$. Let $A$ be the central product of $\gen x$ and $\gen y$ with an element of order $d$ identified, which has presentation
    \[\gen{x, y \mid x^n = y^m = 1, xy = yx, x^{n/d} = y^{m/d}}.\]
    Describe $A$ as a direct product of two cyclic groups.
    \down
    \noindent\textbf{Note}: While the problem specifically calls for two direct products, the proof shows that for any general $n, m$, $A$ is actually cyclic with order $Z_{nm/d}$, so the second direct product is trivial. I'm not sure if this is a typo, but just know that the problem will have a trivial direct product in the end.
\end{specialexercise}

\begin{sol}
    Since we have an element of order $d$ identified, it is easy to see that $|A| = nm/d$. By \hyperref[ex5.2.5]{Exercise 5.2.5}, we know that $\exp(Z_n \times Z_m) = \lcm(n, m) = nm/d$. Note that $|x| = n$ and $|y| = m$ in $A$ still, for otherwise there would be a contradiction to the order of $A$ as we would not be able to account for all the elements of $A$.

    Again, by \hyperref[ex5.2.5]{Exercise 5.2.5}, we know that $A$ contains an element of order $n$ iff $n \mid \exp(A)$, and $A$ contains an element of order $m$ iff $m \mid \exp(A)$. Since $\exp(A) \mid |A| = nm/d$, we have that $\exp(A)$ is a common multiple of $n$ and $m$, so $\exp(A)$ is a multiple of $\lcm(n, m) = nm/d$, hence $\lcm(n, m) \mid \exp(A)$. On the other hand, $\exp(A) \mid |A| = nm/d$ since the order of any element divides the order of the group, so $\exp(A) \mid \lcm(n, m)$. Thus, $\exp(A) = \lcm(n, m)$. We now conclude that $A$ is cyclic, so $A \cong Z_{nm/d} \times Z_1$.
\end{sol}

\newpage

\begin{exercise}\label{ex5.2.13}
    Let $A = \gen{x_1} \times \gen{x_2} \times \cdots \times \gen{x_r}$ be a finite abelian group with $|x_i| = n_i$ for $1 \leq i \leq r$. Find a presentation for $A$. Prove that if $G$ is any group containing commuting elements $g_1, \ldots, g_r$ such that $g_i^{n_i} = 1$ for $1 \leq i \leq r$, then there is a unique homomorphism from $A$ to $G$ which sends $x_i$ to $g_i$ for all $i$.
\end{exercise}

\begin{sol}
    A presentation for $A$ is given by
    \[\gen{x_1, x_2, \ldots, x_r \mid x_i^{n_i} = 1 \text{ for } 1 \leq i \leq r, x_ix_j = x_jx_i \text{ for } 1 \leq i < j \leq r}.\]
    Now, observe that every element of $A$ is of the form $x_1^{s_1} x_2^{s_2} \cdots x_r^{s_r}$ for some integers $s_1, s_2, \ldots, s_r$. Define $\phi : A \to G$ by $\phi(x_1^{s_1} x_2^{s_2} \cdots x_r^{s_r}) = g_1^{s_1} g_2^{s_2} \cdots g_r^{s_r}$. Since the $g_i$ commute, $\phi$ is well-defined. Moreover, $\phi$ is a homomorphism since
    \begin{align*}
        \phi((x_1^{s_1} x_2^{s_2} \cdots x_r^{s_r})(x_1^{t_1} x_2^{t_2} \cdots x_r^{t_r})) & = \phi(x_1^{s_1 + t_1} x_2^{s_2 + t_2} \cdots x_r^{s_r + t_r}) \\
        & = g_1^{s_1 + t_1} g_2^{s_2 + t_2} \cdots g_r^{s_r + t_r} \\
        & = (g_1^{s_1} g_2^{s_2} \cdots g_r^{s_r})(g_1^{t_1} g_2^{t_2} \cdots g_r^{t_r}) \\
        & = \phi(x_1^{s_1} x_2^{s_2} \cdots x_r^{s_r})\phi(x_1^{t_1} x_2^{t_2} \cdots x_r^{t_r})
    \end{align*}
    for all integers $s_1, s_2, \ldots, s_r$ and $t_1, t_2, \ldots, t_r$. Finally, $\phi(x_i) = g_i$ for all $i$ by definition. To see that $\phi$ is the unique homomorphism from $A$ to $G$ sending $x_i$ to $g_i$ for all $i$, let $\psi : A \to G$ be any such homomorphism. Then $\psi(x_1^{s_1} x_2^{s_2} \cdots x_r^{s_r}) = \psi(x_1)^{s_1} \psi(x_2)^{s_2} \cdots \psi(x_r)^{s_r} = g_1^{s_1} g_2^{s_2} \cdots g_r^{s_r}$ for all integers $s_1, s_2, \ldots, s_r$, so $\psi = \phi$.
\end{sol}

\begin{specialexercise}\label{ex5.2.14}
    For any group $G$ define the \textit{dual group} of $G$ (denoted $\wh{G}$) to be the set of all homomorphisms from $G$ into the multiplicative group of roots of unity in $\mathbb{C}$. Define a group operation in $\wh{G}$ by pointwise multiplication of functions: if $\chi, \psi$ are homomorphisms from $G$ into the group of roots of unity then $\chi\psi$ is the homomorphism given by $(\chi\psi)(g) = \chi(g)\psi(g)$ for all $g \in G$, where the latter multiplication takes place in $\mathbb{C}$.
    \begin{subproblems}
        \item Show that this operation on $\wh{G}$ makes $\wh{G}$ into an abelian group. [Show that the identity is the map $g \mapsto 1$ for all $g \in G$ and the inverse of $\chi \in \wh{G}$ is the map $g \mapsto \chi(g)^{-1}$.]
        \item If $G$ is a finite abelian group, prove that $\wh{G} \cong G$. [Write $G$ as $\langle x_1 \rangle \times \cdots \times \langle x_r \rangle$ and if $n_i = |x_i|$ define $\chi_i$ to be the homomorphism which sends $x_i$ to $e^{2\pi i/n_i}$ and sends $x_j$ to $1$, for all $j \neq i$. Prove $\chi_i$ has order $n_i$ in $\wh{G}$ and $\wh{G} = \langle \chi_1 \rangle \times \cdots \times \langle \chi_r \rangle$.]
    \end{subproblems}
\end{specialexercise}

\begin{solalph}
    \item Multiplication in $\c$ is already associative and commutative, so pointwise multiplication of functions is also associative and commutative. The map $\ep : G \to \c$ given by $\ep(g) = 1$ for all $g \in G$ is a homomorphism, since $\ep(gh) = 1 = 1 \cdot 1 = \ep(g)\ep(h)$ for all $g, h \in G$, and 1 is a root of unity. Moreover, for any $\chi \in \wh G$, then $(\chi\ep)(g) = \chi(g)\ep(g) = \chi(g) \cdot 1 = \chi(g)$ for all $g \in G$. Thus, $\ep$ is the identity element of $\wh{G}$. Finally, if $\chi \in \wh{G}$, then the map $\chi^{-1} : G \to \c$ given by $\chi^{-1}(g) = \chi(g)^{-1}$ is a homomorphism, since $\chi^{-1}(gh) = \chi(gh)^{-1} = (\chi(g)\chi(h))^{-1} = \chi(g)^{-1}\chi(h)^{-1} = \chi^{-1}(g)\chi^{-1}(h)$ for all $g, h \in G$, and $\chi(g)^{-1}$ is a root of unity since $\chi(g)$ is a root of unity. Moreover, $\chi\chi^{-1} = \ep$, so $\chi^{-1}$ is the inverse of $\chi$ in $\wh{G}$.
    \item To see that $\chi_i$ has order $n_i$ in $\wh G$, let $g = x_1^{s_1} \cdots x_r^{s_r} \in G$. Then
    \[\chi_i(g) = \chi_i(x_1^{s_1} \cdots x_r^{s_r}) = \chi_i(x_1)^{s_1} \cdots \chi_i(x_r)^{s_r} = e^{2\pi i s_i / n_i},\]
    where all terms $\chi_i(x_j)^{s_j}$ for $j \neq i$ are 1. Thus, $\chi_i^k(g) = e^{2\pi i k s_i / n_i}$ for all integers $k$, so $\chi_i^k = \ep$ if and only if $n_i \mid k$. Thus, $\chi_i$ has order $n_i$ in $\wh{G}$. Now, note by \hyperref[ex5.2.13]{Exercise 5.2.13} that there exists a unique homomorphism $\phi : G \to \wh G$ such that $\phi(x_i) = \chi_i$ for all $i$. We now show that $\phi$ is an isomorphism.

    Let $g = x_1^{s_1} \cdots x_r^{s_r} \in G$ such that $\phi(g) = \ep$. Then for any $x_i$, we have that
    \[\phi(g)(x_i) = (\chi_1^{s_1} \cdots \chi_r^{s_r})(x_i) = \chi_i^{s_i}(x_i) = e^{2\pi i s_i / n_i} = 1.\]
    Thus, $n_i \mid s_i$ for all $i$, so $g = x_1^{s_1} \cdots x_r^{s_r} = 1$. Therefore, $\phi$ is injective. Similarly, for any $\chi \in \wh G$, and $x_i^{n_i} = 1$ for some $i$, then $\chi(x_i)^{n_i} = \chi(x_i^{n_i}) = \chi(1) = 1$, so $\chi(x_i)$ is an $n_i$-th root of unity, so $\chi(x_i) = e^{2\pi i s_i / n_i}$ for some integer $s_i$. Thus, $\chi = \phi(x_1^{s_1} \cdots x_r^{s_r})$, so $\phi$ is surjective. Therefore, $\phi$ is an isomorphism, so $\wh{G} \cong G$.
\end{solalph}

\begin{specialexercise}\label{ex5.2.15}
    Let $G = \gen x \times \gen y$, where $|x| = 8$ and $|y| = 4$.
    \begin{subproblems}
        \item Find all pairs $a, b$ in $G$ such that $G = \gen a \times \gen b$ (where $a$ and $b$ are expressed in terms of $x$ and $y$).
        \item Let $H = \gen{x^2y, y^2} \cong Z_4 \times Z_2$. Prove that there are no elements $a, b$ of $G$ such that $G = \gen a \times \gen b$ and $H = \gen{a^2} \times \gen{b^2}$ (i.e., one cannot pick direct product generators for $G$ in such a way that some powers of these are direct product generators for $H$).
    \end{subproblems}
\end{specialexercise}

\begin{solalph}
    \item From the description of $G$, we see that $G$ is of type $(8, 4)$, so we have the following conditions for $G = \gen a \times \gen b$:
    \begin{itemize}
        \item $|a| = 8$ and $|b| = 4$,
        \item $\gen a \cap \gen b = 1$.
    \end{itemize}
    Let $a = x^s y^t$ and $b = x^u y^v$ for some integers $s, t, u, v$. Then $|a| = \lcm(8/\gcd(8, s), 4/\gcd(4, t))$ and $|b| = \lcm(8/\gcd(8, u), 4/\gcd(4, v))$. For $|a|$ to be 8, we must have $\gcd(8, s) = 1$, implying that $s$ must be odd, while $t$ may range freely. For $|b|$ to be 4, then $8/\gcd(8, u) \mid 4$ and $\max(8/\gcd(8, u), 4/\gcd(4, v)) = 4$, so we have the following cases:
    \begin{enumerate}[label=(\arabic*)]
        \item If $\gcd(8, u) = 8$, then $u = 0$ and $v$ must be odd.
        \item If $\gcd(8, u) = 4$, then $u = 4$ and $v$ must be odd.
        \item If $\gcd(8, u) = 2$, then $u = 2, 6$ and $v$ can be free.
    \end{enumerate}
    To have trivial intersection, we must have that $x^{ks} y^{kt} = x^{\ell u} y^{\ell v} = 1$, or $ks \equiv \ell u \bmod 8$ and $kt \equiv \ell v \bmod 4$ implies that $k \equiv 0 \bmod 8$ and $\ell \equiv 0 \bmod 4$, which is seen from $a^k = 1$ and $b^\ell = 1$ iff $8 \mid k$ and $4 \mid \ell$. Thus, we check the cases for $|b|$:
    \begin{enumerate}[label=(\arabic*)]
        \item $u = 0$ and $v$ is odd: We have $ks \equiv 0 \bmod 8$ and $kt \equiv \ell v \bmod 4$. Since $s$ is odd, then $k \equiv 0 \bmod 8$. Then $kt \equiv 0 \bmod 4$, hence $\ell v \equiv 0 \bmod 4$. Since $v$ is odd, then $\ell \equiv 0 \bmod 4$, so we have trivial intersection.
        \item $u = 4$ and $v$ is odd: We have $ks \equiv 4\ell \bmod 8$ and $kt \equiv \ell v \bmod 4$. Since $s$ is odd, then $k \equiv 4\ell s\inv \bmod 8$. Note that $k \equiv 0 \bmod 4$, so $kt \equiv 0 \bmod 4$, hence $\ell v \equiv 0 \bmod 4$. Since $v$ is odd, then $\ell \equiv 0 \bmod 4$, so we have trivial intersection.
        \item $u = 2$ and $v$ free (the case $u = 6$ and $v$ free is similar): We have $ks \equiv 2\ell \bmod 8$ and $kt \equiv \ell v \bmod 4$. Since $s$ is odd, then $k \equiv 2 \ell s\inv \bmod 8$. We know that intersection will be non-trivial when $k \not\equiv 0 \bmod 8$ and $\ell \not\equiv 0 \bmod 4$. By inspection, we see that $\ell = 2$ gives $k \equiv 4s\inv \equiv 4 \bmod 8$ and $4t \equiv 2v \bmod 4$, which has non-trivial solution when $v$ is even. Inspecting odd $v$, we see that $k \equiv 2\ell s\inv \bmod 8$ reduced mod 4 gives $k \equiv 2 \ell s\inv \bmod 4$. Then $2 \ell s\inv \equiv \ell v \bmod 4$ implies that $\ell(2s\inv - v) \equiv 0 \bmod 4$. Since $v$ is odd, then $2s\inv - v$ is odd, so $\ell \equiv 0 \bmod 4$, so we have trivial intersection.
    \end{enumerate}
    Therefore, the pairs $a, b$ such that $G = \gen a \times \gen b$ are given by
    \begin{itemize}
        \item $a = x^s y^t$ and $b = y^v$ for some odd integers $s, v$ and any integer $t$,
        \item $a = x^s y^t$ and $b = x^4 y^v$ for some odd integers $s, v$ and any integer $t$,
        \item $a = x^s y^t$ and $b = x^2 y^v$ for some odd integer $s$, any integer $t$, and any odd integer $v$, or
        \item $a = x^s y^t$ and $b = x^6 y^v$ for some odd integer $s$, any integer $t$, and any odd integer $v$.
    \end{itemize}
    \item We begin with the elements of $H$:
    \[H = \set{1, x^2y, x^4y^2, x^6y^3, y^2, x^2y^3, x^4y, x^6y^2}\]
    Note that $H = \gen{x^2y} \times \gen{y^2}$, so it suffices to check whether or not $a^2$ generates $\gen{x^2y}$ from the choices of $a$ from part (a). Note that $a^2 = x^{2s} y^{2t}$, which implies that $y$ will always have an even exponent. However, the only element with even exponent on $y$ from $\gen{x^2y}$ is with the term $x^4y^2$, which would imply that $\gen{x^4y^2} = \gen{x^2y}$, which is a contradiction since $x^2y$ has order 4 while $x^4y^2$ has order 2. Thus, there are no elements $a, b$ of $G$ such that $G = \gen a \times \gen b$ and $H = \gen{a^2} \times \gen{b^2}$.
\end{solalph}

\newpage

\begin{exercise}\label{ex5.2.16}
    Prove that no finitely generated abelian group is divisible (cf. \hyperref[ex2.4.19]{Exercise 2.4.19}).
\end{exercise}

\begin{sol}
    Let $A$ be a finitely generated abelian group. By Theorem 5.5, we can write $A$ as a direct product of cyclic groups:
    \[A \cong \z^r \times Z_{n_1} \times Z_{n_2} \times \cdots \times Z_{n_t},\]
    where $n_i > 1$ for all $i$ and $r \geq 0$. For $r > 0$, observe that $\z$ is not divisible, since there is no $x \in \z$ such that $2x = 1$. For $r = 0$, observe that $A$ is a finite abelian group. By \hyperref[ex2.4.19]{Exercise 2.4.19}, finite abelian groups are not divisible. Thus, no finitely generated abelian group is divisible.
\end{sol}

\newpage

\subsection{Table of Groups of Small Order}

\begin{exercise}\label{ex5.3.1}
    Prove that $D_{16}, Z_2 \times D_8, Z_2 \times Q_8, Z_4 * D_8, QD_{16}$, and $M$ are non-isomorphic, non-abelian groups of order 16 (where $Z_4 * D_8$ is described in \hyperref[ex5.1.12]{Exercise 5.1.12} and $QD_{16}$ and $M$ are described in \hyperref[ex2.5.11]{Exercise 2.5.11} and \hyperref[ex2.5.14]{Exercise 2.5.14} respectively).
\end{exercise}

\begin{sol}
    We do a ``decision tree''-like process to show that these are non-isomorphic:
    \begin{itemize}
        \item $D_{16}$
        \begin{itemize}
            \item $Z_2 \times D_8$ has no element of order 8, while $D_{16}$ does.
            \item $Z_2 \times Q_8$ has no element of order 8, while $D_{16}$ does.
            \item $Z_4 * D_8$ has no elements of order 8, since quotienting preserves only divisibility of orders from $Z_4 \times D_8$, which has no elements of order 8, while $D_{16}$ does.
            \item In $QD_{16}$, we have $|\sigma^4| = 2$, and $|\sigma^k\tau| = 2$ when
            \[\sigma^k\tau\sigma^k\tau = \sigma^k\sigma^{3k}\tau\tau = \sigma^{4k},\]
            hence $\sigma^{4k} = 1$ when $k$ is even, hence there are only 5 elements of order 2 in $QD_{16}$, while there are 8 reflections of order 2, plus the central rotation $r^4$, for a total of 9 elements of order 2 in $D_{16}$.
            \item $M$ has $|u| = 2$, with
            \[uv^kuv^k = uuv^{5k}v^k = v^{6k},\]
            so $v^{6k} = 1$ when $8 \mid 6k$, or $4 \mid 3k$, or $4 \mid k$. Then only $k = 4$ works, giving $M$ a total of 3 elements of order 2.
        \end{itemize}
        \item $Z_2 \times D_8$
        \begin{itemize}
            \item $Z_2 \times Q_8$ does not have the same number of elements of order 2 as $Z_2 \times D_8$ does, since $D_8$ has 5 elements of order 2, while $Q_8$ has only 1 element of order 2.
            \item Note that $\bar{(x^k, q)} \in Z_4 * Q_8 \cong Z_4 * D_8$ has order 2 when $(x^k, q)^2 = (1, 1)$ or $(x^2, -1)$. Since we already identify $x^2$ with $-1$ in the quotient, then the case of $q^2 = 1$ and $2i \equiv 0 \bmod 4$ gives the same element, hence only 1 element of order 2. The other case of $q^2 = -1$ and $2i \equiv 2 \bmod 4$, or $i$ odd and $q = Q_8 \setminus Z(Q_8)$ gives 6 elements of order 2, for a total of 7 elements of order 2 in $Z_4 * D_8$, while $Z_2 \times D_8$ has 11 elements of order 2.
            \item $QD_{16}$ contains elements of order 8, while $Z_2 \times D_8$ does not.
            \item $M$ contains elements of order 8, while $Z_2 \times D_8$ does not.
        \end{itemize}
        \item $Z_2 \times Q_8$
        \begin{itemize}
            \item Note that $Z_2 \times Q_8$ only has 3 elements of order 2, while $Z_4 * D_8$ has 7 elements of order 2.
            \item $QD_{16}$ contains elements of order 8, while $Z_2 \times Q_8$ does not.
            \item $M$ contains elements of order 8, while $Z_2 \times Q_8$ does not.
        \end{itemize}
        \item $Z_4 * D_8$
        \begin{itemize}
            \item $QD_{16}$ contains elements of order 8, while $Z_4 * D_8$ does not.
            \item $M$ contains elements of order 8, while $Z_4 * D_8$ does not.
        \end{itemize}
        \item $QD_{16}$
        \begin{itemize}
            \item $M$ has only 3 elements of order 2, while $QD_{16}$ has 5 elements of order 2. \qh
        \end{itemize}
    \end{itemize}
\end{sol}

\newpage

\subsection{Recognizing Direct Products}

Let $G$ be a group.

\begin{exercise}\label{ex5.4.1}
    Prove that if $x, y \in G$, then $[y, x] = [x, y]\inv$. Deduce that for any subsets $A$ and $B$ of $G$, $[A, B] = [B, A]$ (recall that $[A, B]$ is the \textit{subgroup} of $G$ generated by the commutators $[a, b]$).
\end{exercise}

\begin{sol}
    Note that
    \[[x, y]\inv = (x\inv y\inv xy)\inv = yx y\inv x\inv = [y, x].\]
    Hence, $[y, x] = [x, y]\inv$. For any subsets $A$ and $B$ of $G$, the subgroup $[A, B]$ is generated by the commutators $[a, b]$ for $a \in A$ and $b \in B$. Since $[a, b] = [b, a]\inv$, it follows that $[A, B] = [B, A]$.
\end{sol}

\begin{exercise}\label{ex5.4.2}
    Prove that a subgroup $H$ of $G$ is normal if and only if $[G, H] \leq H$.
\end{exercise}

\begin{sol}
    By \hyperref[ex5.4.1]{Exercise 5.4.1}, we know $[H, G] = [G, H]$. Then use Proposition 5.7(2).
\end{sol}

\begin{exercise}\label{ex5.4.3}
    Let $a, b, c \in G$. Prove that
    \begin{subproblems}
        \item $[a, bc] = [a, c](c\inv [a, b]c)$,
        \item $[ab, c] = (b\inv[a, c]b)[b, c]$.
    \end{subproblems}
\end{exercise}

\begin{solalph}
    \item $[a, bc] = a\inv (bc)\inv a(bc) = a\inv c\inv b\inv a b c = a\inv c\inv a c (c\inv a\inv b\inv a b c) = [a, c](c\inv [a, b]c)$.
    \item $[ab, c] = (ab)\inv c\inv (ab)c = b\inv a\inv c\inv a b c = b\inv a\inv c\inv a c (c\inv a\inv b\inv a b c) = (b\inv[a, c]b)[b, c]$. \qh
\end{solalph}

\begin{exercise}\label{ex5.4.4}
    Find the commutator subgroups of $S_4$ and $A_4$.
\end{exercise}

\begin{sol}
    Note that $A_4 \nsub S_4$ and $S_4/A_4 \cong Z_2$, which is abelian, so $[S_4, S_4] \leq A_4$. On the other hand, $A_4$ is generated by the 3-cycles, and $[S_4, S_4]$ contains all 3-cycles since for any 3-cycle $(a\ b\ c)$, we have that
    \[(a\ b\ c) = (a\ b)\inv (a\ c)\inv (a\ b) (a\ c) = [(a\ c), (a\ b)],\]
    so $A_4 \leq [S_4, S_4]$. Thus, $[S_4, S_4] = A_4$. 

    By \hyperref[ex3.5.9]{Exercise 3.5.9}, we know that $V_4 \nsub A_4$ and $A_4/V_4 \cong Z_3$, which is abelian, so $[A_4, A_4] \leq V_4$. On the other hand, $A_4$ is non-abelian, so $[A_4, A_4]$ is non-trivial. Since $V_4$ is the only non-trivial subgroup of $A_4$ that is normal, we must have $[A_4, A_4] = V_4$. \qh
\end{sol}

\begin{exercise}\label{ex5.4.5}
    Prove that $A_n$ is the commutator subgroup of $S_n$ for all $n \geq 5$.
\end{exercise}

\begin{sol}
    Similar argument as in the previous exercise, where $S_n/A_n \cong Z_2$.
\end{sol}

\newpage

\begin{exercise}\label{ex5.4.6}
    Exhibit a representative of each cycle type of $A_5$ as a commutator in $S_5$.
\end{exercise}

\begin{sol}
    The cycle types of $A_5$ are the identity, 3-cycles, products of two disjoint 2-cycles, and 5-cycles. The identity is the commutator of any element with itself, so we can take $[1, 1] = 1$. For a 3-cycle $(a\ b\ c)$, we have that
    \[(a\ b\ c) = (a\ b)\inv (a\ c)\inv (a\ b) (a\ c) = [(a\ c), (a\ b)].\]
    For a product of two disjoint 2-cycles $(a\ b)(c\ d)$, we have that
    \[(a\ b)(c\ d) = (a\ b\ d)\inv (b\ c\ d)\inv (a\ b\ d) (b\ c\ d) = [(a\ b\ d), (b\ c\ d)].\]
    For a 5-cycle $(a\ b\ c\ d\ e)$, we have that
    \[(a\ b\ c\ d\ e) = (a\ e\ d)\inv ((b\ e)(c\ d))\inv (a\ e\ d) ((b\ e)(c\ d)) = [(a\ e\ d), (b\ e)(c\ d)]. \qh\]
\end{sol}

\begin{exercise}\label{ex5.4.7}
    Prove that if $p$ is a prime and $P$ is a non-abelian group of order $p^3$, then $P' = Z(P)$.
\end{exercise}

\begin{sol}
    Let $P$ be a non-abelian group of order $p^3$. By the class equation, we have that $|Z(P)|$ divides $p^3$ and is greater than 1, so $|Z(P)| = p$ or $p^2$. If $|Z(P)| = p^2$, then $P/Z(P)$ is cyclic, so $P$ is abelian, which is a contradiction. Thus, $|Z(P)| = p$, so $P/Z(P)$ is of order $p^2$, hence abelian. Therefore, $P' \leq Z(P)$. On the other hand, if $P' < Z(P)$, then properness of $P'$ implies that $|P'| = 1$, so $P$ is abelian, which is a contradiction. Thus, $P' = Z(P)$.
\end{sol}

\begin{exercise}\label{ex5.4.8}
    Assume $x, y \in G$, and both $x$ and $y$ commute with $[x, y]$. Prove that for all $n \in \zp$,
    \[(xy)^n = x^n y^n [y, x]^{\frac{n(n-1)}{2}}.\]
\end{exercise}

\begin{sol}
    We first prove that for all $n \in \zp$,
    \[y^nx = xy^n [y, x]^n.\]
    We do this by induction on $n$. For $n = 1$, we have $yx = (xyy\inv x\inv)(yx) = xy[y, x]$. Now, assume that $y^nx = xy^n [y, x]^n$ for some $n \in \zp$. Then
    \[y^{n+1}x = y(y^nx) = y(xy^n [y, x]^n) = (yx)y^n [y, x]^n = xy[y, x]y^n [y, x]^n = xy^{n+1} [y, x]^{n+1}.\]
    Now, we prove the main statement by induction on $n$. For $n = 1$, we have $(xy)^n = xy = x^n y^n [y, x]^{\frac{n(n-1)}{2}}$. Now, assume that $(xy)^n = x^n y^n [y, x]^{\frac{n(n-1)}{2}}$ for some $n \in \zp$. Then
    \begin{align*}
        (xy)^{n+1} & = (xy)^n (xy) \\
        & = (x^n y^n [y, x]^{\frac{n(n-1)}{2}})(xy) \\
        & = x^n y^n [y, x]^{\frac{n(n-1)}{2}} x y \\
        & = x^n y^n x y [y, x]^{\frac{n(n-1)}{2}} \\
        & = x^{n+1} y^{n+1} [y, x]^{\frac{(n+1)n}{2}}. \qh
    \end{align*}
\end{sol}

\newpage

\begin{exercise}\label{ex5.4.9}
    Prove that if $p$ is an odd prime, and $P$ is a group of order $p^3$, then the $p$-th power map $x \mapsto x^p$ is a homomorphism of $P$ into $Z(P)$. If $P$ is not cyclic, show that the kernel of the $p$-th power map has order $p^2$ or $p^3$. Is the squaring map a homomorphism in non-abelian groups of order 8? Where is the oddness of $p$ needed in the above proof?
\end{exercise}

\begin{sol}
    Note that the abelian case is trivial, so we may assume $P$ is non-abelian. By \hyperref[ex5.4.7]{Exercise 5.4.7}, we have $P' = Z(P)$, so $P/Z(P)$ is abelian. Note that $[y, x] \in P' = Z(P)$ for all $x, y \in P$, so $x$ and $y$ both commute with $[y, x]$. Moreover, $|Z(P)| = p$ by \hyperref[ex5.4.7]{Exercise 5.4.7}, hence
    \[(xy)^p = x^p y^p [y, x]^{\frac{p(p-1)}{2}} = x^p y^p,\]
    where the last equality holds since $[y, x]$ has order dividing $p$ and $\frac{p(p-1)}{2}$ is divisible by $p$ when $p$ is odd. Thus, the $p$-th power map is a homomorphism of $P$ into $Z(P)$. To see that $x^p \in Z(P)$ for all $x \in P$, note that for any $y \in P$, we have
    \[y x^py\inv = (y x y\inv)^p = (x [y, x])^p = x^p [y, x]^p = x^p.\]
    Suppose now that $P$ is not cyclic. Since the map is to $Z(P)$, which has order $p$, the image of the map has order dividing $p$. If the image is trivial, then the kernel is all of $P$, so the kernel has order $p^3$. If the image is non-trivial, then the kernel has order $p^2$ by the First Isomorphism Theorem.

    For the squaring map in non-abelian groups of order 8, we have $Q_8$ and $D_8$. In $Q_8$, observe that $(ij)^2 = (k)^2 = -1$, while $i^2j^2 = 1$. In $D_8$, we similarly have $(rs)^2 = 1$, while $r^2s^2 = r^2$. Thus, the squaring map is not a homomorphism in non-abelian groups of order 8. Lastly, the oddness of $p$ is needed in the proof to ensure that $\frac{p(p-1)}{2}$ is divisible by $p$, which allows us to conclude that $[y, x]^{\frac{p(p-1)}{2}} = 1$.
\end{sol}

\begin{specialexercise}\label{ex5.4.10}
    Prove that a finite abelian group is the direct product of its Sylow subgroups.
\end{specialexercise}

\begin{sol}
    Inducting on the number of prime factors of the group, we first consider $n = 1$, or when $|G| = p_1^n$ for some prime $p_1$. Then $G$ is its own Sylow $p_1$-subgroup, so the statement holds. We now assume that any abelian with fewer than $n$ prime factors is the direct product of its Sylow subgroups, and let $G$ be an abelian group with $n$ prime factors. Put
    \[|G| = p_1^{\alpha_1} p_2^{\alpha_2} \cdots p_n^{\alpha_n}.\]
    Let $H$ be some subgroup of $G$ generated by the Sylow $p_i$ subgroups for $1 \leq i \leq n - 1$. By definition, every element of $H$ is generated by some element from a Sylow $p_i$-subgroup, so its order must divide the lcm $p_1^{\alpha_1} p_2^{\alpha_2} \cdots p_{n-1}^{\alpha_{n-1}}$, so $H$ is a proper subgroup of $G$. Moreover, $H$ contains all such Sylow $p_i$-subgroups for $1 \leq i \leq n - 1$, hence $|H| \geq p_1^{\alpha_1} p_2^{\alpha_2} \cdots p_{n-1}^{\alpha_{n-1}}$. Since $H$ is a proper subgroup of $G$, we must have $|H| = p_1^{\alpha_1} p_2^{\alpha_2} \cdots p_{n-1}^{\alpha_{n-1}}$. By the inductive hypothesis, $H$ is the direct product of its Sylow subgroups, so
    \[H = P_1 \times P_2 \times \cdots \times P_{n-1},\]
    where $P_i$ is the Sylow $p_i$-subgroup of $G$ for $1 \leq i \leq n - 1$. Since $G$ is abelian, we have that $P_n$ is normal in $G$, so $HP_n$ is a subgroup of $G$. Moreover, since $H$ and $P_n$ have trivial intersection because $p_n^{\alpha_n}$ and $p_1^{\alpha_1} p_2^{\alpha_2} \cdots p_{n-1}^{\alpha_{n-1}}$ are coprime, we have that $HP_n = H \times P_n$. Since $|HP_n| = |H||P_n| = |G|$, we must have $HP_n = G$, so
    \[G = P_1 \times P_2 \times \cdots \times P_n.\] \qh
\end{sol}

\newpage

\begin{specialexercise}\label{ex5.4.11}
    Prove that if $G = HK$ where $H$ and $K$ are characteristic subgroups of $G$ with $H \cap K = 1$, then $\aut(G) \cong \aut(H) \times \aut(K)$. Deduce that if $G$ is an abelian group of finite order, then $\aut(G)$ is isomorphic to the direct product of the automorphism groups of its Sylow subgroups.
\end{specialexercise}

\begin{sol}
    Let $\phi \in \aut(G)$. Since $H$ and $K$ are characteristic subgroups of $G$, we have $\phi(H) = H$ and $\phi(K) = K$. Thus, we can define $\phi_H = \phi|_H$ and $\phi_K = \phi|_K$, which are automorphisms of $H$ and $K$ respectively. We can then define a map $\Phi: \aut(G) \to \aut(H) \times \aut(K)$ by $\Phi(\phi) = (\phi_H, \phi_K)$. To see that $\Phi$ is a homomorphism, let $\phi, \psi \in \aut(G)$ be given. Then
    \[\Phi(\phi \circ \psi) = ((\phi \circ \psi)_H, (\phi \circ \psi)_K) = (\phi_H \circ \psi_H, \phi_K \circ \psi_K) = \Phi(\phi) \circ \Phi(\psi).\]
    To see that $\Phi$ is injective, let $\phi \in \aut(G)$ be given such that $\Phi(\phi) = (1_H, 1_K)$, where $1_H$ and $1_K$ are the identity automorphisms of $H$ and $K$ respectively. Then $\phi(h) = h$ for all $h \in H$ and $\phi(k) = k$ for all $k \in K$. Since every element of $G$ can be uniquely expressed as $hk$ for some $h \in H$ and $k \in K$, we have $\phi(hk) = \phi(h)\phi(k) = hk$, so $\phi$ is the identity automorphism of $G$. Thus, $\Phi$ is injective.

    To show that $\Phi$ is surjective, let $(\psi_H, \psi_K) \in \aut(H) \times \aut(K)$ be given. We can define a map $\psi: G \to G$ by $\psi(hk) = \psi_H(h)\psi_K(k)$ for all $h \in H$ and $k \in K$. Because $H \cap K = 1$, and $G = HK$, every element of $G$ can be uniquely expressed as $hk$ for some $h \in H$ and $k \in K$. We check that $\psi$ is a homomorphism:
    \[\psi(h_1k_1 \cdot h_2k_2) = \psi(h_1h_2k_1k_2) = \psi_H(h_1h_2)\psi_K(k_1k_2) = \psi_H(h_1)\psi_H(h_2)\psi_K(k_1)\psi_K(k_2) = \psi(h_1k_1)\psi(h_2k_2).\]
    where $k_1$ and $h_2$ commute since $H$ and $K$ are the factors of $G = H \times K$. Since $\psi_H$ and $\psi_K$ are automorphisms, $\psi$ is a bijection, hence an automorphism of $G$. Moreover, $\Phi(\psi) = (\psi_H, \psi_K)$, so $\Phi$ is surjective.

    Therefore, we have shown that $\Phi: \aut(G) \to \aut(H) \times \aut(K)$ is an isomorphism. For the second part of the problem, since a finite abelian group is the direct product of its Sylow subgroups by the previous exercise, and conjugation is trivial in an abelian group, we have that each Sylow subgroup is characteristic, so the automorphism group of the finite abelian group is isomorphic to the direct product of the automorphism groups of its Sylow subgroups.
\end{sol}

\begin{exercise}\label{ex5.4.12}
    Use Theorem 4.17 to describe the automorphism group of a finite cyclic group.
\end{exercise}

\begin{sol}
    From Proposition 4.17, we know that $\aut(Z_p) \cong Z_{p - 1}$, and $\aut(Z_{p^n}) \cong Z_{p^{n - 1}(p - 1)}$ for any prime $p$ and $n \geq 2$. For $n \geq 3$ and $p = 2$, we have that $\aut(Z_{2^n}) \cong Z_2 \times Z_{2^{n - 2}}$ that is not cyclic. Hence, for any finite cyclic group $Z_n$ with (specific) prime factorization $n = 2^{\alpha_1} p_2^{\alpha_2} \cdots p_k^{\alpha_k}$, we have by \hyperref[ex5.4.10]{Exercise 5.4.10} and the previous exercise that
    \[\aut(Z_n) \cong \aut(Z_{2^{\alpha_1}}) \times \aut(Z_{p_2^{\alpha_2}}) \times \cdots \times \aut(Z_{p_k^{\alpha_k}}).\]
    Note that $\aut(Z_{p_i^{\alpha_i}}) \cong Z_{p_i^{\alpha_i - 1}(p_i - 1)}$ is cyclic for $i \geq 2$, and $\aut(Z_{2^{\alpha_1}})$ is cyclic if $\alpha_1 \leq 2$ and is not cyclic if $\alpha_1 \geq 3$. Thus, the structure of $\aut(Z_n)$ depends on the exponent of 2 in the prime factorization of $n$. We can check the following cases:
    \begin{itemize}
        \item $\alpha_1 = 0$: Then the 2-part of $n$ is trivial, so $\aut(Z_n) \cong Z_{p_2^{\alpha_2 - 1}(p_2 - 1)} \times \cdots \times Z_{p_k^{\alpha_k - 1}(p_k - 1)}$ is cyclic if and only if $k = 1$, or the $p_i^{\alpha_i - 1}(p_i - 1)$ are pairwise coprime.
        \item $\alpha_1 = 1$: Then the 2-part of $n$ is $Z_2$ with $\aut(Z_2)$ trivial, so $\aut(Z_n) \cong Z_{p_2^{\alpha_2 - 1}(p_2 - 1)} \times \cdots \times Z_{p_k^{\alpha_k - 1}(p_k - 1)}$ is cyclic if and only if the orders $p_i^{\alpha_i - 1}(p_i - 1)$ are pairwise coprime and coprime to 2, which is equivalent to the condition in the previous case. Hence, the same condition as $\alpha_1 = 0$ applies here.
        \item $\alpha_1 = 2$: Then the 2-part of $n$ is $Z_4$, so $\aut(Z_n) \cong \aut(Z_4) \times Z_{p_2^{\alpha_2 - 1}(p_2 - 1)} \times \cdots \times Z_{p_k^{\alpha_k - 1}(p_k - 1)}$ is cyclic if and only if $k = 1$, similar to $\alpha_1 = 0$, or the $p_i^{\alpha_i - 1}(p_i - 1)$ are pairwise coprime and coprime to 2.
        \item $\alpha_1 \geq 3$: Then the 2-part of $n$ is $Z_2 \times Z_{2^{\alpha_1 - 2}}$, so $\aut(Z_n) \cong \aut(Z_{2^{\alpha_1}}) \times Z_{p_2^{\alpha_2 - 1}(p_2 - 1)} \times \cdots \times Z_{p_k^{\alpha_k - 1}(p_k - 1)}$ is not cyclic. \qh
    \end{itemize}
\end{sol}

\begin{exercise}\label{ex5.4.13}
    Prove that $D_{8n}$ is not isomorphic to $D_{4n} \times Z_2$ for any $n \geq 1$.
\end{exercise}

\begin{sol}
    Note that $Z(D_{8n}) = \set{1, r^{2n}}$ has order 2, while $Z(D_{4n} \times Z_2) = Z(D_{4n}) \times Z(Z_2)$ has order 4, so $D_{8n}$ cannot be isomorphic to $D_{4n} \times Z_2$ for any $n \geq 1$.
\end{sol}

\begin{exercise}\label{ex5.4.14}
    Let $G = \set{(a_{ij}) \in \gl_n(F) \mid a_{ij} = 0 \text{ for } i > j, \text{ and } a_{11} = a_{22} = \cdots = a_{nn}}$, where $F$ is a field, be the group of upper triangular matrices all of whose diagonal entries are equal. Prove that $G \cong D \times U$, where $D$ is the group of all nonzero multiples of the identity matrix, and $U$ is the group of upper triangular matrices with 1's down the diagonal.
\end{exercise}

\begin{sol}
    Note that $I_n \in D$. For any $\lambda, \mu \in F\unt$, then $(\lambda I_n)(\mu I_n)\inv = (\lambda I_n)(\mu\inv I_n) = (\lambda\mu\inv)I_n \in D$, so $D$ is a subgroup of $G$. For $A, B \in U$, we know that the product still has $AB$ upper triangular with 1's down the diagonal, so $U$ is a subgroup of $G$. To see that $D \nsub G$, let $A \in G$ be given. Then $A$ is upper triangular with all diagonal entries equal to some $\lambda \in F\unt$. For any $\mu I_n \in D$, we have
    \[A(\mu I_n)A\inv = \mu AI_n A\inv = \mu A A\inv = \mu I_n \in D,\]
    so $D$ is normal in $G$. To see that $U \nsub G$, let $A \in G$ be given. Then $A$ is upper triangular with all diagonal entries equal to some $\lambda \in F\unt$. We may factor out $\lambda = a_{11}$ from $A$ to get $A = a_{11} X$, where $X$ is upper triangular with 1's down the diagonal, so $X \in U$. For any $Y \in U$, we have
    \[A Y A\inv = a_{11} X Y X\inv a_{11}\inv = X Y X\inv,\]
    where $XYX \inv \in U$ since $U \leq G$, so $U$ is normal in $G$. If $Z \in D \cap U$, then $Z = \lambda I_n$ for some $\lambda \in F\unt$, and $Z$ is upper triangular with 1's down the diagonal, so $\lambda = 1$, hence $Z = I_n$. Thus, $D \cap U = 1$. Lastly, for any $A \in G$, we can factor out the common diagonal entry to get $A = a_{11} X$ for some $X \in U$, so $G = DU \cong D \times U$ by Theorem 5.9.
\end{sol}

\begin{exercise}\label{ex5.4.15}
    If $A$ and $B$ are normal subgroups of $G$ such that $G/A$ and $G/B$ are both abelian, prove that $G/(A \cap B)$ is abelian.
\end{exercise}

\begin{sol}
    $G/A$ abelian implies $G' \leq A$, and $G/B$ abelian implies $G' \leq B$, so $G' \leq A \cap B$. Hence, $G/(A \cap B)$ is abelian by Proposition 5.7(4).
\end{sol}

\begin{exercise}\label{ex5.4.16}
    Prove that if $K$ is a normal subgroup of $G$, then $K' \nsub G$.
\end{exercise}

\begin{sol}
    Let $K$ be a normal subgroup of $G$. By Proposition 5.7(2), we have $[G, K] \leq K$. Since $K' = [K, K] \leq [G, K]$, we have $K' \leq K$. Moreover, for any $g \in G$ and $x = c_1 c_2 \cdots c_n \in K'$, where $c_i = [k_i, \ell_i]$ for some $k_i, \ell_i \in K$, we have
    \[g x g\inv = g c_1 c_2 \cdots c_n g\inv = (g c_1 g\inv)(g c_2 g\inv) \cdots (g c_n g\inv),\]
    where $g c_i g\inv = g [k_i, \ell_i] g\inv = [g k_i g\inv, g \ell_i g\inv]$ since $K$ is normal in $G$, so $g x g\inv \in K'$. Thus, $K' \nsub G$.
\end{sol}

\begin{exercise}\label{ex5.4.17}
    If $K$ is a normal subgroup of $G$ and $K$ is cyclic, prove that $G' \leq C_G(K)$. [Recall that the automorphism group of a cyclic group is abelian.]
\end{exercise}

\begin{sol}
    Let $K$ be a normal subgroup of $G$ and $K$ be cyclic. Since $K$ is cyclic, we have $\aut(K)$ is abelian. Moreover, the homomorphism of $G$ into $\aut(K)$ given by conjugation has kernel $C_G(K)$, so $G/C_G(K)$ is isomorphic to a subgroup of $\aut(K)$, which is abelian. Hence, $G' \leq C_G(K)$ by Proposition 5.7(4).
\end{sol}

\begin{exercise}\label{ex5.4.18}
    Let $K_1, K_2, \ldots, K_n$ be non-abelian simple groups, and let $G = K_1 \times K_2 \times \cdots \times K_n$. Prove that every normal subgroup of $G$ is of the form $G_I$ for some subset $I$ of $\set{1, 2, \ldots, n}$ (where $G_I$ is defined in \hyperref[ex5.1.2]{Exercise 5.1.2}). [If $N \nsub G$ and $x = (a_1, \ldots, a_n) \in N$ with some $a_i \neq 1$, then show that there is some $g_i \in G_i$ not commuting with $a_i$. Show that $[(1, \ldots, g_i, \ldots, 1), x] \in K_i \cap N$, and deduce $K_i \leq N$.]
\end{exercise}

\begin{sol}
    Let $N$ be a normal subgroup of $G$. If $N = 1$, then $N = G_\emptyset$. If $N = G$, then $N = G_{\set{1, 2, \ldots, n}}$. Now, assume that $1 < N < G$. Then there exists some $x = (a_1, a_2, \ldots, a_n) \in N$ such that $a_i \neq 1$ for some $i$. Since $K_i$ is non-abelian simple, there exists some $g_i \in K_i$ such that $g_i$ does not commute with $a_i$. Consider the element $(1, \ldots, g_i, \ldots, 1) \in G$, where $g_i$ is in the $i$-th position. We have
    \[(1, \ldots, g_i, \ldots, 1) x (1, \ldots, g_i\inv, \ldots, 1) = (a_1, \ldots, g_i a_i g_i\inv, \ldots, a_n).\]
    Hence,
    \[(1, \ldots, g_i, \ldots, 1) x (1, \ldots, g_i\inv, \ldots, 1) x\inv = (1, \ldots, [g_i, a_i], \ldots, 1) \in N.\]
    Since $[g_i, a_i]$ is non-trivial and is an element of the simple group $K_i$, we have that $\langle [g_i, a_i]^{K_i} \rangle \nsub K_i$. However, simplicity of $K_i$ forces $\langle [g_i, a_i]^{K_i} \rangle = K_i$, so $K_i \leq N$. By repeating this argument for any other non-trivial component of an element of $N$, we can find some subset $I$ of $\set{1, 2, \ldots n}$ such that $G_I = K_{i_1} \times K_{i_2} \times \cdots K_{i_k} \leq N$, where $\set{i_1, i_2, \ldots i_k} = I$, which shows $G_I \leq N$. On the other hand, if $N$ contains some element $x = (a_1, a_2, \ldots, a_n)$ such that $a_j \neq 1$ for some $j \notin I$, then we can similarly find some element of $N$ that has non-trivial component in the $j$-th position, which contradicts that $K_j \not\leq N$. Thus, $N \leq G_I$. Therefore, $N = G_I$.
\end{sol}

\begin{exercise}\label{ex5.4.19}
    A group $H$ is called \textit{perfect} if $H' = H$ (i.e., $H$ equals its own commutator subgroup).
    \begin{subproblems}
        \item Prove that every non-abelian simple group is perfect.
        \item Prove that if $H$ and $K$ are perfect subgroups of a group $G$, then $\gen{H, K}$ is also perfect. Extend this to show that the subgroup of $G$ generated by any collection of perfect subgroups is perfect.
        \item Prove that any conjugate of a perfect subgroup is perfect.
        \item Prove that any group $G$ has a maximal perfect subgroup, and that this subgroup is normal.
    \end{subproblems}
\end{exercise}

\begin{solalph}
    \item Let $H$ be a non-abelian simple group. Since $H$ is non-abelian, we have $H' \neq 1$. Since $H$ is simple, the only normal subgroups of $H$ are $1$ and $H$, so we must have $H' = H$.
    \item Let $H$ and $K$ be perfect subgroups of a group $G$. We want to show that $\gen{H, K}' = \gen{H, K}$. Since $H$ and $K$ are perfect, we have $H' = H$ and $K' = K$, so $H \leq H'$ and $K \leq K'$. Moreover, since $H$ and $K$ are subgroups of $\gen{H, K}$, we have $H' \leq \gen{H, K}'$ and $K' \leq \gen{H, K}'$, so $H \leq \gen{H, K}'$ and $K \leq \gen{H, K}'$. Hence, $\gen{H, K} = \gen{H, K}'$. The extension to any collection of perfect subgroups follows by induction on the number of perfect subgroups in the collection.
    \item Let $H$ be a perfect subgroup of a group $G$, and let $g \in G$ be given. We want to show that $g\inv H g$ is perfect. Since $H$ is perfect, we have $H' = H$, so $g\inv H' g = g\inv H g$. Moreover, since conjugation is an automorphism of $G$, we have $g\inv H' g = (g\inv H g)'$, so $(g\inv H g)' = g\inv H g$. Hence, $g\inv H g$ is perfect.
    \item Let $\mathcal{P}$ be the collection of perfect subgroups of $G$. Since $1$ is a perfect subgroup of $G$, we have $\mathcal{P}$ is nonempty. Let $\mathcal{C}$ be a chain in $\mathcal{P}$, and let $H = \bigcup_{K \in \mathcal{C}} K$. To see that $H \leq G$, note that for any $x, y \in H$, then there exists $K_x, K_y \in \mathcal C$ such that $x \in K_x$ and $y \in K_y$. Without loss of generality, we may assume $K_x \leq K_y$, so $x, y \in K_y$, hence $xy\inv \in K_y \leq H$. Thus, $H$ is a subgroup of $G$. To see that $H \in \pp$, observe that for any $x \in H$, then $x \in K$ for some $K \in \mathcal{C}$, so $x \in K' = K$, hence $x \in H'$. Thus, $H \leq H'$. Since $H' \leq H$ trivially, then $H = H'$, hence $H \in \mathcal P$. By Zorn's Lemma, $\mathcal{P}$ has a maximal element $M$. 
    
    To see that $M$ is normal in $G$, let $g \in G$ be given. Then $g\inv M g$ is perfect by part (c), so $g\inv M g \in \mathcal{P}$. Since $M$ is perfect, we may use part (b) to get that $\gen{M, g\inv M g}$ is perfect, so $\gen{M, g\inv M g} \in \mathcal{P}$. Since $M \leq \gen{M, g\inv M g}$ and $M$ is maximal in $\mathcal{P}$, we must have $M = \gen{M, g\inv M g}$, so $g\inv M g \leq M$. By symmetry of the argument, we also have $M \leq g\inv M g$, so $g\inv M g = M$. Hence, $M$ is normal in $G$.
\end{solalph}

\begin{exercise}\label{ex5.4.20}
    Let $H(F)$ be the Heisenberg group over the field $F$, cf. \hyperref[ex1.4.11]{Exercise 1.4.11}. Find an explicit formula for the commutator $[X, Y]$, where $X, Y \in H(F)$, and show that the commutator subgroup of $H(F)$ equals the center of $H(F)$ (cf. \hyperref[ex2.2.14]{Exercise 2.2.14}).
\end{exercise}

\begin{sol}
    Let
    \[X = 
    \begin{bmatrix}
        1 & a & b \\
        0 & 1 & c \\
        0 & 0 & 1
    \end{bmatrix}, \quad
    Y = 
    \begin{bmatrix}
        1 & d & e \\
        0 & 1 & f \\
        0 & 0 & 1
    \end{bmatrix} \in H(F).\]
    Then we have
    \[X\inv =
    \begin{bmatrix}
        1 & -a & ac - b \\
        0 & 1 & -c \\
        0 & 0 & 1
    \end{bmatrix}, \quad
    Y\inv =
    \begin{bmatrix}
        1 & -d & df - e \\
        0 & 1 & -f \\
        0 & 0 & 1
    \end{bmatrix}.\]
    Hence,
    \begin{align*}
        [X, Y] & = X\inv Y\inv XY \\
        & =
        \begin{bmatrix}
            1 & -a & ac - b \\
            0 & 1 & -c \\
            0 & 0 & 1
        \end{bmatrix}
        \begin{bmatrix}
            1 & -d & df - e \\
            0 & 1 & -f \\
            0 & 0 & 1
        \end{bmatrix}
        \begin{bmatrix}
            1 & a & b \\
            0 & 1 & c \\
            0 & 0 & 1
        \end{bmatrix}
        \begin{bmatrix}
            1 & d & e \\
            0 & 1 & f \\
            0 & 0 & 1
        \end{bmatrix} \\
        & =
        \begin{bmatrix}
            1 & 0 & af - cd \\
            0 & 1 & 0 \\
            0 & 0 & 1
        \end{bmatrix}
    \end{align*}
    Recall that the center of $H(F)$ consists of all matrices of the form
    \[
    \begin{bmatrix}
        1 & 0 & z \\
        0 & 1 & 0 \\
        0 & 0 & 1
    \end{bmatrix}\]
    for some $z \in F$. Clearly, $af - cd \in F$, so $[X, Y] \in Z(H(F))$, hence $H(F)' \leq Z(H(F))$. On the other hand, for any $z \in F$, we may choose $a = 1$, $f = z$, and $c = d = 0$ to get that $Z(H(F)) \leq H(F)'$. Therefore, $H(F)' = Z(H(F))$.
\end{sol}

\newpage

\subsection{Semidirect Products}

Let $H$ and $K$ be groups, let $\phi : K \to \aut(H)$ be a homomorphism, and identify $H, K \leq G = H \semidir_\phi K$.

\begin{exercise}\label{ex5.5.1}
    Prove that $C_K(H) = \ker\phi$ (recall that $C_K(H) = C_G(H) \cap K$).
\end{exercise}

\begin{sol}
    Let $k \in C_K(H)$. Then $k \in K$ and $kh = hk$ for all $h \in H$. Then $\phi(k)(h) = khk\inv = h$ for all $h \in H$, so $k \in \ker\phi$. Conversely, let $k \in \ker\phi$. Then $\phi(k)(h) = khk\inv = h$ for all $h \in H$, so $kh = hk$ for all $h \in H$, hence $k \in C_K(H)$. Therefore, $C_K(H) = \ker\phi$.
\end{sol}

\begin{exercise}\label{ex5.5.2}
    Prove that $C_H(K) = N_H(K)$.
\end{exercise}

\begin{sol}
    It is clear that $C_H(K) \leq N_H(K)$. Now let $h \in N_H(K)$. Then $hKh\inv = K$, so for any $k \in K$, we have $hkh\inv \in K$. Since $H \cap K = 1$, and $[h, k] = hkh\inv k\inv \in H \cap K$, we must have $[h, k] = 1$, so $hk = kh$ for all $k \in K$. Hence, $h \in C_H(K)$, so $N_H(K) \leq C_H(K)$. Therefore, $C_H(K) = N_H(K)$.
\end{sol}

\begin{exercise}\label{ex5.5.3}
    In Example 1 following the proof of Proposition 5.11, prove that every element of $G - H$ has order 2. Prove that $G$ is abelian if and only if $h^2 = 1$ for all $h \in H$.
\end{exercise}

\begin{sol}
    Recall that $H \nsub G$. It is clear to see that $G/H = \set{H, Hk}$ for $k \in H$, since $H$ consists of elements $hk$ with trivial $k$ component. Then for any $hk \in Hk$, and noting the fact that $kh = h\inv k$ for all $h \in H$, we have
    \[(hk)^2 = h(kh)k = h(h\inv k)k = kk = 1.\]
    If $hk = 1$, then $h = k\inv \in H \cap K = 1$, so $hk = 1$ if and only if $h = k = 1$. Hence, every element of $G - H$ has order 2.

    If $G$ is abelian, then for any $h \in H$ and $k \in K$, we have $hk = kh$. Moreover, $kh = h\inv k$, hence $hk = h\inv k$, so $h^2 = 1$. Conversely, if $h^2 = 1$ for all $h \in H$, then any $h \in H$ and $k \in K$ commute, because $kh = h\inv k = hk$ by $h^2 = 1$. Then for any $h_1, h_2 \in H$ and $k_1, k_2 \in K$, we have
    \[(h_1k_1)(h_2k_2) = h_1(h_2k_1)k_2 = (h_2h_1)(k_2k_1) = (h_2k_2)(h_1k_1),\]
    so $G$ is abelian.
\end{sol}

\begin{exercise}\label{ex5.5.4}
    Let $p = 2$ and check that the construction of the two non-abelian groups of order $p^3$ is valid in this case. Prove that \textit{both} resulting groups are isomorphic to $D_8$.
\end{exercise}

\begin{sol}
    Let $U = \gen u \cong Z_4$ and $V = \gen v \cong Z_2$. Then $\aut(U) \cong Z_2$, so there is a unique non-trivial homomorphism $\phi : V \to \aut(U)$, which sends $v$ to the automorphism of $U$ that sends $u$ to $u\inv$. Then the semidirect product $U \semidir_\phi V$ has order 8, and is generated by $u$ and $v$ with relations
    \[u^4 = v^2 = 1, \quad vu = u\inv v.\]
    It is clear that $U \semidir_\phi V$ is non-abelian, since $vu \neq uv$. Let $r = u$ and $s = v$. Then $r^4 = s^2 = 1$, and $sr = vu = u\inv v = r\inv s$, so $U \semidir_\phi V$ has the same presentation as $D_8$, hence $U \semidir_\phi V \cong D_8$.

    Now let $W = \gen w \times \gen x \cong V_4$ and $Y = \gen y \cong Z_2$. Then $\aut(W) \cong S_3$. Since $S_3$ contains 3 transpositions that are conjugate to each other, choose some $\phi : Y \to \aut(W)$ that sends $y$ to a transposition in $\aut(W)$. Then the semidirect product $W \semidir_\phi Y$ has order 8, and is generated by $w$, $x$, and $y$ with relations
    \[w^2 = x^2 = y^2 = 1, \quad yw = x y, \quad yx = w y.\]
    It is clear that $W \semidir_\phi Y$ is non-abelian, since $yw \neq wy$. Let $r = yw$ and $s = y$. Then $s^2 = 1$,
    \[r^4 = (r^2)^2 = ((yw)^2)^2 = (xyyw)^2 = (xw)^2 = x^2w^2 = 1,\]
    and
    \[sr = y(yw) = w = w\inv = w\inv y\inv y = (yw)\inv y = r\inv s,\]
    so $W \semidir_\phi Y$ has the same presentation as $D_8$, hence $W \semidir_\phi Y \cong D_8$.
\end{sol}

\begin{exercise}\label{ex5.5.5}
    Let $G = \hol(Z_2 \times Z_2)$.
    \begin{subproblems}
        \item Prove that $G = H \semidir K$, where $H = Z_2 \times Z_2$ and $K \cong S_3$. Deduce that $|G| = 24$.
        \item Prove that $G$ is isomorphic to $S_4$. [Obtain a homomorphism from $G$ into $S_4$ by letting $G$ act on the left cosets of $K$. Use \hyperref[ex5.5.1]{Exercise 5.5.1} to show that this representation is faithful.]
    \end{subproblems}
\end{exercise}

\begin{solalph}
    \item Note that $\aut(H) \cong \gl_2(\f_2) \cong S_3$ as it is a non-abelian group of order 6, by \hyperref[ex4.2.10]{Exercise 4.2.10}. It follows by definition that $G = H \semidir K$, where $K \cong S_3$. Hence, $|G| = |H||K| = 4 \cdot 6 = 24$.
    \item Let $G$ act on the left cosets of $K$ in $G$. From $|G : K| = 4$, we have a homomorphism $\phi : G \to S_4$. Moreover, we have some homomorphism $\psi : K \to \aut(H)$ such that $G \cong H \semidir_\psi K$. Lastly, we know from Theorem 4.3 that
    \[\ker\phi = \bigcap_{g \in G} gKg^{-1} \nsub G\]
    such that $\ker\phi \leq K$. Putting $N = \ker\phi$, then $N = \set{1, A_3, K}$. If $N = K$, then $K \nsub G$, which is a contradiction since $K$ is not normal in $G$ because $\psi$ is injective. If $N = A_3$, let $n \in N$ be a nonidentity element. Then $hnh\inv \in N$ for every $h$. However, \hyperref[ex5.5.1]{Exercise 5.5.1} says that $\ker\psi = C_K(H)$. Since $\psi$ is injective (an isomorphism, actually), then $C_K(H) = 1$, so $n$ does not commute with some $h_0 \in H$, i.e., $nh_0n\inv \neq h_0$. However, from $N \nsub G$, we know that $h_0nh_0\inv \in N$. 
    
    Observe that $h_0nh_0\inv = (h_0nh_0\inv n\inv)n = [h_0, n]n$. Observe that $h_0(nh_0\inv n\inv) \in H$ by definition of the semidirect product, so $[h_0, n] \in H$. If $[h_0, n] = 1$, then $h_0nh_0\inv = n \in N$, contradicting our choice of $h_0$. Hence $h_0nh_0\inv$ has non-trivial component in $H$, so $h_0nh_0\inv \notin N$, which is a contradiction. Therefore, $N = 1$, so $\phi$ is injective, hence $G \cong \phi(G) \leq S_4$. Since $|G| = |S_4| = 24$, we have $G \cong S_4$.
\end{solalph}

\begin{specialexercise}\label{ex5.5.6}
    Assume that $K$ is a cyclic group, $H$ is an arbitrary group, and $\phi_1$ and $\phi_2$ are homomorphisms from $K$ into $\aut(H)$ such that $\phi_1(K)$ and $\phi_2(K)$ are conjugate subgroups of $\aut(H)$. If $K$ is infinite, assume $\phi_1$ and $\phi_2$ are injective. Prove, by constructing an explicit isomorphism, that
    \[H \semidir_{\phi_1} K \cong H \semidir_{\phi_2} K\]
    (in particular, if the subgroups $\phi_1(K)$ and $\phi_2(K)$ are equal in $\aut(H)$, then the resulting semidirect products are isomorphic). [Suppose $\sigma\phi_1(K)\sigma\inv = \phi_2(K)$ so that for some $a \in \z$ we have $\sigma\phi_1(k)\sigma\inv = \phi_2(k)^a$ for all $k \in K$. Show that the map
    \[\psi : H \semidir_{\phi_1} K \to H \semidir_{\phi_2} K\]
    defined by $\psi((h,k)) = (\sigma(h), k^a)$ is a homomorphism. Show $\psi$ is bijective by constructing a 2-sided inverse.]
\end{specialexercise}

\begin{sol}
    Let $K = \gen k$ and $\sigma \in \aut(H)$ such that $\sigma\phi_1(K)\sigma\inv = \phi_2(K)$. Then there exists some $a \in \z$ such that $\sigma\phi_1(k)\sigma\inv = \phi_2(k)^a$ for all $k \in K$. Define the map
    \[\psi : H \semidir_{\phi_1} K \to H \semidir_{\phi_2} K\]
    by $\psi((h, k^m)) = (\sigma(h), k^{am})$ for all $h \in H$ and $m \in \z$. Then for any $(h_1, k^s), (h_2, k^t) \in H \semidir_{\phi_1} K$, we have
    \begin{align*}
        \psi((h_1, k^s)(h_2, k^t)) & = \psi((h_1\phi_1(k^s)(h_2), k^{s + t})) \\
        & = (\sigma(h_1\phi_1(k^s)(h_2)), k^{a(s + t)}) \\
        & = (\sigma(h_1)\sigma(\phi_1(k^s)(h_2)), k^{as}k^{at}) \\
        & = (\sigma(h_1)\phi_2(k^{as})(\sigma(h_2)), k^{as}k^{at}) \\
        & = (\sigma(h_1), k^{as})(\sigma(h_2), k^{at}) \\
        & = \psi((h_1, k^s))\psi((h_2, k^t)),
    \end{align*}
    so $\psi$ is a homomorphism. To see that $\psi$ is bijective, we split into two cases.
    \begin{enumerate}
        \item Suppose $K$ is infinite. By conjugate subgroups, we know that $\phi_2(K) = \phi_2(K^a)$, where $K^a = \set{k^a \mid k \in K}$. Since $\phi_2$ is injective, then $K^a = K$, forcing $a = \pm 1$. Then $\theta((h, k)) = (\sigma\inv(h), k^a)$ is a 2-sided inverse of $\psi$, so $\psi$ is bijective.
        \item We first prove a lemma: if $m \mid n$, then the map defined by
        \[\nu : \units[n] \to \units[m], \quad \nu(a) = a \bmod m\]
        is a surjective homomorphism. We first check that $\nu$ is well-defined: if $a \equiv b \bmod n$, then $n \mid (a - b)$, so $m \mid (a - b)$, hence $a \equiv b \bmod m$. Next, we check that $\nu$ is a homomorphism: for any $a, b \in \units[n]$, we have
        \[\nu(ab) = ab \bmod m = (a \bmod m)(b \bmod m) = \nu(a)\nu(b).\]
        Finally, we check that $\nu$ is surjective: let $a \in \units[m]$. If $a \in \units[n]$, then $\nu(a) = a$. If $a \not\in \units[n]$, then $(a, n) \neq 1$. Let $c$ be the product of primes that divide $n$ but not $a$, let $p$ be a prime that divides $n$, and put $b = a + cm$. We claim that $b \in \units[n]$ as follows:
        \begin{itemize}
            \item If $p \mid a$, then $p \nmid c$ by definition of $c$, hence $p \nmid b$.
            \item If $p \nmid a$, then $p \mid c$ by definition of $c$, hence $p \nmid b$.
        \end{itemize}
        Thus, $b \in \units[n]$, and $\nu(b) = b \bmod m = a + cm \bmod m = a$. Hence, $\nu$ is surjective.
        
        Suppose $K \cong Z_n$ is finite. Let $|\phi_1(K)| = |\phi_2(K)| = m$, where $m \mid n$. Note that $\gen{\phi_1(k)} = \phi_1(K)$ and $\gen{\phi_2(k)} = \phi_2(K)$, so $|\phi_1(k)| = |\phi_2(k)| = m$. Moreover, $\gen{\phi_2(k)^a} = \phi_2(K)$, so $|\phi_2(k)^a| = m$. It follows that $(a, m) = 1$. 
        
        Consider the surjective homomorphism $\nu : \units[n] \to \units[m]$ defined in the lemma above. Since $(a, m) = 1$, we have $a \in \units[m]$, so there exists some $b \in \units[n]$ such that $\nu(b) = b \bmod m = a$. Then $\theta((h, k)) = (\sigma\inv(h), k^b)$ is a 2-sided inverse of $\psi$, so $\psi$ is bijective. \qh
    \end{enumerate}
\end{sol}

\begin{specialexercise}\label{ex5.5.7}
    This exercise describes thirteen isomorphism types of groups of order 56. (It is not too difficult to show that every group of order 56 is isomorphic to one of these.)
    \begin{subproblems}
        \item Prove that there are three abelian groups of order 56.
        \item Prove that every group of order 56 has either a normal Sylow 2-subgroup or a normal Sylow 7-subgroup.
        \item Construct the following non-abelian groups of order 56 which have a normal Sylow 7-subgroup and whose Sylow 2-subgroup $S$ is as specified:
        \begin{itemize}
            \item one group when $S \cong Z_2 \times Z_2 \times Z_2$
            \item two non-isomorphic groups when $S \cong Z_4 \times Z_2$
            \item one group when $S \cong Z_8$
            \item two non-isomorphic groups when $S \cong Q_8$
            \item three non-isomorphic groups when $S \cong D_8$
        \end{itemize}
        [For a particular $S$, two groups are not isomorphic if the kernels of the maps from $S$ into $\aut(Z_7)$ are not isomorphic.]
        \item Let $G$ be a group of order 56 with a non-normal Sylow 7-subgroup. Prove that if $S$ is the Sylow 2-subgroup of $G$ then $S \cong Z_2 \times Z_2 \times Z_2$. [Let an element of order 7 act by conjugation on the seven nonidentity elements of $S$ and deduce that they all have the same order.]
        \item Prove that there is a unique group of order 56 with a non-normal Sylow 7-subgroup. [For existence use the fact that $|GL_3(\f_2)| = 168$; for uniqueness use \hyperref[ex5.5.6]{Exercise 5.5.6}.]
    \end{subproblems}
\end{specialexercise}

\begin{solalph}
    \item Note that $56 = 2^3 \cdot 7$. By the Fundamental Theorem of Finite Abelian Groups, the partitions of 3 are $3, 2 + 1, 1 + 1 + 1$, so the abelian groups of order 56 are of the form $Z_8 \times Z_7$, $Z_4 \times Z_2 \times Z_7$, and $Z_2 \times Z_2 \times Z_2 \times Z_7$.
    \item Let $n_7$ be the number of Sylow 7-subgroups of $G$. Then $n_7 \equiv 1 \bmod 7$ and $n_7 \mid 2^3$, so $n_7 = 1$ or $8$. If $n_7 = 1$, then the unique Sylow 7-subgroup is normal. Suppose $n_7 = 8$. Then there are $8 \cdot 6 = 48$ elements of order 7 in $G$, so the remaining $56 - 48 = 8$ elements of $G$ must be the identity and the elements of a Sylow 2-subgroup. Hence, the Sylow 2-subgroup is unique, so it is normal.
    \item Let $P \in \syl_7(G)$ such that $P \nsub G$. Then $P = \gen p \cong Z_7$, and $\aut(P) \cong Z_6$.
    \begin{itemize}
        \item If $S \cong Z_2 \times Z_2 \times Z_2$, put $S = \gen a \times \gen b \times \gen c$. Let $\phi : S \to \aut(P)$ be a homomorphism. We disregard $a, b, c$ being sent to $\id$ as that is the trivial homomorphism. Since $a, b, c$ have order 2, they can only be sent to elements of order 1 or 2 in $\aut(P)$. Put $\alpha \in \aut(P)$ such that $|\alpha| = 2$, where $\alpha$ is unique since $\aut(P)$ is cyclic. Then we have the following non-trivial homomorphisms:
        \[\phi_1 = 
        \begin{cases}
            a \mapsto \alpha, \\
            b \mapsto \id, \\
            c \mapsto \id,
        \end{cases} \quad
        \phi_2 =
        \begin{cases}
            a \mapsto \alpha, \\
            b \mapsto \alpha, \\
            c \mapsto \id,
        \end{cases} \quad
        \phi_3 =
        \begin{cases}
            a \mapsto \alpha, \\
            b \mapsto \alpha, \\
            c \mapsto \alpha,
        \end{cases}\]
        Note that $\ker\phi_1 = \gen b \times \gen c$, $\ker\phi_2 = \gen{ab, c}$, and $\ker\phi_3 = \gen{ab, ac}$. Note that these are all isomorphic to $V_4$, hence the resulting semidirect products are isomorphic by \hyperref[ex5.5.6]{Exercise 5.5.6}. Without loss of generality, consider $\phi_1$. Then $a$ acts on $P$ by sending $p$ to the inversion automorphism of $P$, or $apa\inv = p\inv$. It is clear that $b$ and $c$ centralize $P$ and commute with $a$. Observe that $\gen{p, a \mid apa\inv = p\inv, a^2 = p^7 = 1} \cong D_{14}$ and has trivial intersection with $\gen b \times \gen c$, so
        \[G \cong P \semidir_{\phi_1} S \cong D_{14} \times Z_2 \times Z_2.\]
        \item Put $S = Z_4 \times Z_2 \cong \gen a \times \gen b$. Let $\alpha \in \aut(P)$ be the unique element of order 2 again. Then we have the following non-trivial homomorphisms:
        \[\phi_1 =
        \begin{cases}
            a \mapsto \alpha, \\
            b \mapsto \id,
        \end{cases} \quad
        \phi_2 = 
        \begin{cases}
            a \mapsto \id, \\
            b \mapsto \alpha,
        \end{cases} \quad
        \phi_3 =
        \begin{cases}
            a \mapsto \alpha, \\
            b \mapsto \alpha,
        \end{cases}\]
        We have $\ker\phi_1 = \gen{a^2} \times \gen b, \ker\phi_2 = \gen a$, and $\ker\phi_3 = \gen{ab}$, where $|ab| = 4$. Note that $\ker\phi_1 \cong V_4$, while $\ker\phi_2 \cong \ker\phi_3 \cong Z_4$. We may consider $\phi_2$ first. Then $b$ acts on $P$ by sending $p$ to the inversion automorphism of $P$, or $bpb\inv = p\inv$. It is clear that $a$ centralizes $P$ and commutes with $b$. Observe that $\gen{p, b \mid bpb\inv = p\inv, b^2 = p^7 = 1} \cong D_{14}$ and has trivial intersection with $\gen a$, so
        \[G \cong P \semidir_{\phi_2} S \cong D_{14} \times Z_4.\]
        For $\phi_1$, we see that $\gen b$ centralizes $P$, and $a$ acts on $P$ by sending $p$ to the inversion automorphism of $P$. Moreover, consider the subgroup $P_0 = \gen{p, a \mid apa\inv = p\inv, a^4 = p^7 = 1}$. It is easy to see that $P \cap \gen a = 1$, so $|P_0| = |P||\gen a| = 28$. Moreover, $pPp\inv = P$, and $aPa\inv = P$ since $apa\inv = p\inv$ under the action of $\phi_1$ restricted to $\gen a$. Hence, $P_0 \cong Z_7 \semidir Z_4$. It is clear that $P_0$ is normal in $G$, and $\gen b$ centralizes $P_0$, so
        \[G \cong P \semidir_{\phi_1} S \cong (Z_7 \semidir Z_4) \times Z_2.\]
        \item For $S \cong Z_8$, let $S = \gen s$. Since $\aut(P) \cong Z_6$, and $(8, 6) = 2$, then the only non-trivial homomorphism $\phi : S \to \aut(P)$ is the one that sends $s$ to the inversion automorphism $\alpha$ of $P$. Moreover, cyclicity of $S$ guarantees uniqueness of the homomorphism. Then $sps\inv = p\inv$, and
        \[G \cong P \semidir_\phi S \cong Z_7 \semidir Z_8.\]
        \item Let $S \cong Q_8 = \gen{i, j} = \set{1, i, i^2, i^3, j, ij, i^2j, i^3j}$. If we consider the trivial homomorphism, then $G \cong P \times S \cong Z_7 \times Q_8$. Now consider the non-trivial homomorphism that sends some generator of $S$ to $\alpha$. We have three cases:
        \[\phi_1 = 
        \begin{cases}
            i \mapsto \alpha, \\
            j \mapsto \id,
        \end{cases} \quad
        \phi_2 = 
        \begin{cases}
            i \mapsto \id, \\
            j \mapsto \alpha,
        \end{cases} \quad
        \phi_3 =
        \begin{cases}
            i \mapsto \alpha, \\
            j \mapsto \alpha,
        \end{cases}\]
        We have $\ker\phi_1 = \gen j, \ker\phi_2 = \gen i$, and $\ker\phi_3 = \gen{ij}$. These are all isomorphic to $Z_4$, so the semidirect products will be isomorphic to each other. Pick $\phi_1$ without loss of generality. Then $i$ acts on $P$, and $ipi\inv = p\inv$. Unlike the other cases, $j$ does not centralize $P$, since $jij\inv = i\inv$, so $jpj\inv = p\inv$. Hence, we have
        \[G \cong P \semidir_{\phi_1} S \cong Z_7 \semidir Q_8 = \set{p, i, j \mid i^4 = j^4 = 1, i^2 = j^2, iji\inv = j\inv, iji\inv = jpj\inv = p\inv}\]
        Note that non-triviality of the homomorphism implies that some $s \in S$ does not centralize $P$. If $G$ could be written as a direct product, then $S$ would centralize $P$, which is a contradiction. Hence, $G$ is not a direct product.
        \item Let $S \cong D_8 = \gen{r, s} = \set{1, r, r^2, r^3, s, rs, r^2s, r^3s}$. If we consider the trivial homomorphism, then $G \cong P \times S \cong Z_7 \times D_8$. Now consider the non-trivial homomorphism that sends some generator of $S$ to $\alpha$. We have three cases:
        \[\phi_1 = 
        \begin{cases}
            r \mapsto \alpha, \\
            s \mapsto \id,
        \end{cases} \quad
        \phi_2 =
        \begin{cases}
            r \mapsto \id, \\
            s \mapsto \alpha,
        \end{cases} \quad
        \phi_3 =
        \begin{cases}
            r \mapsto \alpha, \\
            s \mapsto \alpha,
        \end{cases}\]
        We have $\ker\phi_1 = \gen{r^2, s}, \ker\phi_2 = \gen r$, and $\ker\phi_3 = \gen{sr, r^2}$. Then $\ker\phi_1 \cong \ker\phi_3 \cong V_4$, while $\ker\phi_2 \cong Z_4$. Hence, the semidirect products corresponding to $\phi_1$ and $\phi_3$ are isomorphic to each other, but not to the semidirect product corresponding to $\phi_2$. Pick $\phi_1$ without loss of generality. Then $r$ acts on $P$, and $rpr\inv = p\inv$. Moreover, $s$ centralizes $P$, hence
        \[G \cong P \semidir_{\phi_1} S \cong Z_7 \semidir D_8 = \set{p, r, s \mid p^7 = r^4 = s^2 = 1, srs\inv = r\inv, rpr\inv = p\inv}.\]
        If we pick $\phi_2$, then $s$ acts on $P$, and $sps\inv = p\inv$. Moreover, note that $r$ centralizes $P$, hence $|rp| = 28$. Since $s(rp)s = srps\inv = (srs\inv)(sps\inv) = r\inv p\inv = (rp)\inv$, then
        \[G \cong P \semidir_{\phi_2} S \cong D_{56}.\]
    \end{itemize}
    \item Suppose $G$ has a non-normal Sylow 7-subgroup. We know that $S \in \syl_2(G)$ must be normal by part (b). Let $P \in \syl_7(G)$, and put $P = \gen p$. Let $p$ act by conjugation on the seven nonidentity elements of $S$. Note that the stabilizer of any nonidentity element $s \in S$ is a subgroup of $P$, hence $|G_s| = 1$ or 7. If $|G_s| = 7$, then the orbit is trivial, hence $psp\inv = s$, or $p = sps\inv \in P$, which contradicts that $|s|$ is a power of 2. Then $|G_s| = 1$, so the orbit $\oo_s$ of $s$ has size 7. Since $|S - 1| = 7$, then $\oo_s$ contains all nonidentity elements of $S$. Hence, all nonidentity elements of $S$ have the same order.
    
    If they all had order 8, then $s^2$ would have order 4, but $(s^2)^4 = s^8 = 1$. A similar reasoning eliminates the possibility of order 4. Hence, all nonidentity elements of $S$ have order 2, so $S \cong Z_2 \times Z_2 \times Z_2$.
    \item Let $H = Z_2 \times Z_2 \times Z_2$ and $K = \gen k \cong Z_7$. Then $\aut(H) \cong \gl_3(\f_2)$, which has order 168. Let $\phi : K \to \aut(H)$ be a homomorphism. Since $|K| = 7$, then $|\phi(K)| = 1$ or 7. If $|\phi(K)| = 1$, then the resulting semidirect product is abelian, which is a contradiction. Hence, $|\phi(K)| = 7$. Since $\aut(H)$ has order 168, then there exists a unique subgroup of order 7 in $\aut(H)$, call it $\psi$. Then $\phi : K \to \aut(H)$ given by $\phi(k) = \psi$ is a homomorphism, and the resulting semidirect product $H \semidir_\phi K$ has order $|H||K| = 8 \cdot 7 = 56$. We now claim that $K$ is not normal in $H \semidir_\phi K$.
    
    If $K \nsub G$, then $[s, k] \in K$ for all $s \in S$. Moreover, $[s, k] \in S$ since $S \nsub G$. Hence, $[s, k] \in S \cap K = 1$, so $[s, k] = 1$, or $sk = ks$ for all $s \in S$. Then $\phi(k) = \id$, which is a contradiction. Hence, $K$ is not normal in $H \semidir_\phi K$.

    To see that $H \semidir_\phi K$ is unique, let $\phi_1, \phi_2 : K \to \aut(H)$ be two homomorphisms such that $|\phi_1(K)| = |\phi_2(K)| = 7$. Then $\phi_1(K)$ and $\phi_2(K)$ are conjugate subgroups of $\aut(H)$. Since $|\aut(H) = 2^3 \cdot 3 \cdot 7$, then $\phi_1(K), \phi_2(K) \in \syl_7(\aut(H))$, so they are conjugate by the Sylow theorems. Hence, by \hyperref[ex5.5.6]{Exercise 5.5.6}, we have $H \semidir_{\phi_1} K \cong H \semidir_{\phi_2} K$. Therefore, $H \semidir_\phi K$ is unique.
\end{solalph}

\begin{exercise}\label{ex5.5.8}
    Construct a non-abelian group of order 75. Classify all groups of order 75 (there are three of them). [Use \hyperref[ex5.5.6]{Exercise 5.5.6} to show that the non-abelian group is unique.] (The classification of groups of order $pq^2$, where $p$ and $q$ are primes with $p < q$ and $p$ not dividing $q - 1$, is quite similar.)
\end{exercise}

\begin{sol}
    Let $P \in \syl_5(G)$ and $Q \in \syl_3(G)$. Then $|P| = 25$ and $|Q| = 3$. By the Sylow theorems, we have $n_5 \equiv 1 \bmod 5$ and $n_5 \mid 3$, so $n_5 = 1$ or 3. If $n_5 = 1$, then the unique Sylow 5-subgroup is normal. If $n_5 = 3$, then there are $3 \cdot 24 = 72$ elements of order 5 in $G$, so the remaining $75 - 72 = 3$ elements of $G$ must be the identity and the elements of a Sylow 3-subgroup. Hence, the Sylow 3-subgroup is unique, so it is normal. In either case, we have a normal Sylow subgroup, so $G \cong P \semidir Q$ or $G \cong Q \semidir P$. Moreover, $P \cong Z_{25}$ or $Z_5 \times Z_5$, and $Q \cong Z_3$. We then have two cases to consider:
    \begin{itemize}
        \item Suppose $P \cong Z_{25}$. Then $\aut(P) \cong \units[25]$, which has order 20. Since $|Q| = 3$, then the only homomorphism $\phi : Q \to \aut(P)$ is the trivial homomorphism, so $G \cong P \semidir Q \cong P \times Q \cong Z_{25} \times Z_3$.
        
        If we consider $G \cong Q \semidir P$, then we have a homomorphism $\phi : P \to \aut(Q)$. Since $|\aut(Q)| = 2$, then the only homomorphism $\phi : P \to \aut(Q)$ is the trivial homomorphism, so $G \cong Q \semidir P \cong Q \times P \cong Z_3 \times Z_{25}$.
        \item Suppose $P \cong Z_5 \times Z_5$ with $\aut(P) \cong \gl_2(\f_5)$, which has order 480. Since $|Q| = 3$, then there exists a non-trivial homomorphism $\phi : Q \to \aut(P)$, which is unique since $|\aut(P)| = 480$ and $(480, 3) = 3$. We may use similar reasoning as the previous exercise in that $\phi_1(Q)$ and $\phi_2(Q)$ are conjugate subgroups of $\aut(P)$, and $\phi_1(Q), \phi_2(Q) \in \syl_3(\aut(P))$, so they are conjugate by the Sylow theorems. Hence, by \hyperref[ex5.5.6]{Exercise 5.5.6}, we have $P \semidir_{\phi_1} Q \cong P \semidir_{\phi_2} Q$. Therefore, $P \semidir_\phi Q$ is unique, and $G \cong P \semidir_\phi Q$ is non-abelian.
        
        If we consider $G \cong Q \semidir P$, then we have a homomorphism $\phi : P \to \aut(Q)$. Since $|\aut(Q)| = 2$, then the only homomorphism $\phi : P \to \aut(Q)$ is the trivial homomorphism, so $G \cong Q \semidir P \cong Q \times P \cong Z_3 \times (Z_5 \times Z_5)$.
    \end{itemize}
\end{sol}

\begin{exercise}\label{ex5.5.9}
    Show that the matrix
    \[
    \begin{bmatrix}
        0 & -1 \\
        1 & 4
    \end{bmatrix}\]
    is an element of order 5 in $\gl_2(\f_{19})$. Use this matrix to construct a non-abelian group of order 1805, and give a presentation of this group. Classify groups of order 1805 (there are three isomorphism types). [Use \hyperref[ex5.5.6]{Exercise 5.5.6} to prove uniqueness of the non-abelian group.] (A general method for finding elements of prime order in $\gl_n(\f_p)$ is described in the exercises in Section 12.2; this particular matrix of order 5 in $\gl_2(\f_{19})$ appears in \hyperref[ex12.2.16]{Exercise 12.2.16} of that section as an illustration of the method.)
\end{exercise}

\begin{sol}
    Let $M$ be the given matrix. It is easy to calculate that
    \[M^5 = 
    \begin{bmatrix}
       -56 & -209 \\
       209 & 780
    \end{bmatrix},\]
    where, once entries are taken mod 19, results in $I_2$, hence $|M| = 5$. Moreover, we know that $1805 = 19^2 \cdot 5$. Since $n_{19} \equiv 1 \bmod 19$ and $n_{19} \mid 5$, then $n_{19} = 1$, so $P \in \syl_{19}(G)$ is normal. Moreover, we have $P \cong Z_{361}$ or $P \cong Z_{19} \times Z_{19}$. If $P \cong Z_{361}$, then $\aut(P) \cong \units[361]$, which has order 19 * 18 = 342. Since $|Q| = 5$, then the only homomorphism $\phi : Q \to \aut(P)$ is the trivial homomorphism, so $G \cong P \semidir Q \cong P \times Q \cong Z_{361} \times Z_5$.

    If $P \cong Z_{19} \times Z_{19}$, then $\aut(P) \cong \gl_2(\f_{19})$, which has order 123120. If we take $\phi$ to be the trivial homomorphism, then $G \cong Z_{19} \times Z_{19} \times Z_5$. Now, since $|Q| = 5$, then there exists a non-trivial homomorphism $\phi : Q \to \aut(P)$, which is unique since $|\aut(P)| = 123120$ and $(123120, 5) = 5$. We may use similar reasoning as the previous exercise in that $\phi_1(Q)$ and $\phi_2(Q)$ are conjugate subgroups of $\aut(P)$, and $\phi_1(Q), \phi_2(Q) \in \syl_5(\aut(P))$, so they are conjugate by the Sylow theorems. Hence, by \hyperref[ex5.5.6]{Exercise 5.5.6}, we have $P \semidir_{\phi_1} Q \cong P \semidir_{\phi_2} Q$. Therefore, $P \semidir_\phi Q$ is unique, and $G \cong P \semidir_\phi Q$ is non-abelian.
\end{sol}

\begin{specialexercise}\label{ex5.5.10}
    This exercise classifies the groups of order 147 (there are six isomorphism types).
    \begin{subproblems}
        \item Prove that there are two abelian groups of order 147.
        \item Prove that every group of order 147 has a normal Sylow 7-subgroup.
        \item Prove that there is a unique non-abelian group of order 147 whose Sylow 7-subgroup is cyclic.
        \item Let
        \[t_1 = 
        \begin{bmatrix}
            2 & 0 \\
            0 & 1
        \end{bmatrix}, \quad
        t_2 = 
        \begin{bmatrix}
            1 & 0 \\
            0 & 2
        \end{bmatrix}\]
        be elements of $\gl_2(\f_7)$. Prove $P = \gen{t_1, t_2}$ is a Sylow 3-subgroup of $\gl_2(\f_7)$, and that $P \cong Z_3 \times Z_3$. Deduce that every subgroup of order 3 is conjugate in $\gl_2(\f_7)$ to a subgroup of $P$.
        \item By Example 3 in Section 1, the group $P$ has four subgroups of order 3, and these are: $P_1 = \gen{t_1}, P_2 = \gen{t_2}, P_3 = \gen{t_1t_2}$, and $P_4 = \gen{t_1t_2^2}$. For $i = 1, 2, 3, 4$, let $G_i = (Z_7 \times Z_7) \semidir_{\phi_i} P_i$, where $\phi_i$ is an isomorphism of $Z_3$ with the subgroup $P_i$ of $\aut(Z_7 \times Z_7)$. For each $i$, describe $G_i$ in terms of generators and relations. Deduce that $G_1 \cong G_2$.
        \item Prove that $G_1$ is not isomorphic to either $G_3$ or $G_4$. [Show that the center of $G_1$ has order 7, whereas the centers of $G_3$ and $G_4$ are trivial.]
        \item Prove that $G_3$ is not isomorphic to $G_4$.
        \item Classify the groups of order 147 by showing that the six non-isomorphic groups described above (two from part (a), one from part (c), and $G_1, G_3$, and $G_4$) are all the groups of order 147. [Use \hyperref[ex5.5.6]{Exercise 5.5.6} and part (d).] (The classification of groups of order $pq^2$, where $p$ and $q$ are primes with $p < q$ and $p \mid q - 1$, is quite similar.)
    \end{subproblems}
\end{specialexercise}

\begin{solalph}
    \item Note that $147 = 3 \cdot 7^2$. By the Fundamental Theorem of Finite Abelian Groups, the partitions of 2 are $2, 1 + 1$, so the abelian groups of order 147 are of the form $Z_{49} \times Z_3$ and $Z_7 \times Z_7 \times Z_3$.
    \item Let $n_7$ be the number of Sylow 7-subgroups of $G$. Then $n_7 \equiv 1 \bmod 7$ and $n_7 \mid 3$, so $n_7 = 1$. Hence, the unique Sylow 7-subgroup is normal.
    \item Let $P \cong Z_{49}$ and $Q \cong Z_3$. Then $\aut(P) \cong \units[49]$, which has order $7 \cdot 6 = 42$. Since $|Q| = 3$, then there exists non-trivial homomorphisms $\phi : Q \to \aut(P)$, where $q \in Q = \gen q$ is mapped to a generator of the unique subgroup of order 3 in $\aut(P)$, say $\alpha$. Then either $q \mapsto \alpha$ or $q \mapsto \alpha^2$. Let $\phi_1, \phi_2 : Q \to \aut(P)$ be the two homomorphisms defined by $\phi_1(q) = \alpha$ and $\phi_2(q) = \alpha^2$. Since $\phi_1(Q) = \phi_2(Q) = \gen\alpha$, then by \hyperref[ex5.5.6]{Exercise 5.5.6}, we have $P \semidir_{\phi_1} Q \cong P \semidir_{\phi_2} Q$. Hence, there is a unique non-abelian group of order 147 whose Sylow 7-subgroup is cyclic.
    \item Recall that $|\gl_2(\f_7)| = (7^2 - 1)(7^2 - 7) = 48 \cdot 42 = 2016 = 2^5 \cdot 3^2 \cdot 7$. Observe that
    \[t_1^3 = 
    \begin{bmatrix}
        8 & 0 \\
        0 & 1
    \end{bmatrix}, \quad 
    t_2^3 = 
    \begin{bmatrix}
        1 & 0 \\
        0 & 8
    \end{bmatrix}, \quad
    t_1t_2 = 
    \begin{bmatrix}
        2 & 0 \\
        0 & 2
    \end{bmatrix} = t_2t_1.\]
    Taking entries mod 7, then $t_1^3 = t_2^3 = I_2$, so $|t_1| = |t_2| = 3$. Moreover, $t_1t_2 = t_2t_1$, so $P = \gen{t_1}\gen{t_2}$. To see that this is isomorphic to $\gen{t_1} \times \gen{t_2}$, note that $\gen{t_1} \cap \gen{t_2} = 1$. Moreover, commutativity shows that ${t_1} \nsub P$ and ${t_2} \nsub P$, so $P \cong \gen{t_1} \times \gen{t_2} \cong Z_3 \times Z_3$. Since $|P| = 9$, then $P$ is a Sylow 3-subgroup of $\gl_2(\f_7)$.

    Let $R$ be a subgroup of order 3 in $\gl_2(\f_7)$. Then $R \leq gPg\inv$ for some $g \in \gl_2(\f_7)$ by the Sylow theorems. Then $gRg\inv \leq P$, so $gRg\inv$ is a subgroup of order 3 in $P$. Hence, every subgroup of order 3 is conjugate in $\gl_2(\f_7)$ to a subgroup of $P$.
    \item Let $Z_7 \times Z_7 = \gen a \times \gen b$. Then $\aut(Z_7 \times Z_7) \cong \gl_2(\f_7)$, and $P_i$ is a subgroup of $\aut(Z_7 \times Z_7)$ for $i = 1, 2, 3, 4$. Let $\phi_i : Z_3 \to P_i$ be an isomorphism. We check each $i$:
    \begin{itemize}
        \item For $i = 1$, we have $\phi_1 : Z_3 \to P_1$ given by $\phi_1(q) = t_1$. Then $qaq\inv = \phi_1(q)(a) = t_1(a) = a^2$, and $qbq\inv = \phi_1(q)(b) = t_1(b) = b$. Hence, we have for a na\"ive presentation of $G_1$:
        \[G_1 = \set{a, b, q \mid a^7 = b^7 = q^3 = 1, ab = ba, qaq\inv = a^2, qbq\inv = b}.\]
        However, note that $b$ is central in $G_1$ as it commutes with both $a$ and $q$. It follows that $\gen b \nsub G_1$ with $|\gen b| = 7$. Observe that $\gen b \cap \gen{a, q} = 1$ by construction. To see that $\gen{a, q} \nsub G_1$, observe that $b\gen{a, q}b\inv = \gen{a, q}$ because $b \in Z(G_1)$. Then $G_1 \cong Z_7 \times \gen{a, q}$.

        If we consider a homomorphism $\psi : Z_3 \to \aut(Z_7)$, then $q \mapsto \alpha$ or $q \mapsto \alpha^2$ such that $\alpha$ generates the unique subgroup of order 3 in $\aut(Z_7)$. In either case, we have $\gen{a, q} \cong Z_7 \semidir_\psi Z_3$ as it satisfies the relations $qaq\inv = a^2$ in the na\"ive presentation. Hence, we have
        \[G_1 \cong Z_7 \times (Z_7 \semidir_\psi Z_3).\]
        \item For $i = 2$, simply note that $t_2$ takes the place of $t_1$ in the previous case, $aq = qa$, and $qbq\inv = b^2$. Hence, we have
        \[G_2 = \set{a, b, q \mid a^7 = b^7 = q^3 = 1, ab = ba, qaq\inv = a, qbq\inv = b^2}.\]
        By similar reasoning as the previous case, we have $G_2 \cong Z_7 \times (Z_7 \semidir_\psi Z_3)$, so $G_1 \cong G_2$.
        \item For $i = 3$, we have $\phi_3 : Z_3 \to P_3$ given by $\phi_3(q) = t_1t_2$. Then $qaq\inv = \phi_3(q)(a) = t_1t_2(a) = t_1(a) = a^2$, and $qbq\inv = \phi_3(q)(b) = t_1t_2(b) = t_1(b^2) = b^2$. A presentation of $G_3$ is
        \[G_3 = \set{a, b, q \mid a^7 = b^7 = q^3 = 1, ab = ba, qaq\inv = a^2, qbq\inv = b^2}.\]
        \item For $i = 4$, we have $\phi_4 : Z_3 \to P_4$ given by $\phi_4(q) = t_1t_2^2$. Then $qaq\inv = \phi_4(q)(a) = t_1t_2^2(a) = t_1(a) = a^2$, and $qbq\inv = \phi_4(q)(b) = t_1t_2^2(b) = t_1(b^4) = b^4$. Then a presentation of $G_4$ is
        \[G_4 = \set{a, b, q \mid a^7 = b^7 = q^3 = 1, ab = ba, qaq\inv = a^2, qbq\inv = b^4}.\]
    \end{itemize}
    \item The prior discussion showed that $Z(G_1) = \gen b$, so $|Z(G_1)| = 7$. Now, consider $G_3$, and let $z \in Z(G_3)$. We first consider when $z \in Z(G_3) \cap (\gen a \times \gen b)$. Then $z = a^ib^j$ for some $0 \leq i, j \leq 6$. Since $z$ commutes with $q$, then 
    \[a^ib^j = z = qzq\inv = qa^ib^jq\inv = qaq\inv qbq\inv = a^{2i}b^{2j}.\]
    Then this implies that $2i \equiv i \bmod 7$ and $2j \equiv j \bmod 7$, so $i = j =0$. Hence, $Z(G_3) \cap (\gen a \times \gen b) = 1$. Now, consider when $z \in Z(G_3) - (\gen a \times \gen b)$. Then $z = a^ib^jq^k$ for some $0 \leq i, j \leq 6$ and $1 \leq k \leq 2$. Since $z$ commutes with $a$, then 
    \[a^ib^jq^k = z = aza\inv = aa^ib^jq^ka\inv = a^{i+1}b^jq^ka\inv = a^{i+1}b^jq^ka^{-1}.\]
    However, since $q$ does not centralize $\gen a$, then this is a contradiction. Hence, $Z(G_3) - (\gen a \times \gen b) = 1$, so $Z(G_3) = 1$. A similar argument shows that $Z(G_4) = 1$.
    \item To show that $G_3 \not\cong G_4$, we first observe that $G_3$ and $G_4$ have the same subgroups of order 7, since $Z_7 \times Z_7$ is the unique Sylow 7-subgroup of both groups. Hence, $G_3$ and $G_4$ have the list of subgroups
    \[\ss = \set{\gen a, \gen b, \gen{ab}, \gen{ab^2}, \gen{ab^3}, \gen{ab^4}, \gen{ab^5}, \gen{ab^6}}.\]
    Let $q$ act on $\ss$ by conjugation. In $G_3$, for any $\gen{ab^i} \in \ss$, then
    \[q(ab^i)q\inv = (qaq\inv)(qb^iq\inv) = a^2b^{2i}.\]
    It is clear that $a^2b^{2i} \in \gen{ab^i}$, so $q$ stabilizes $\gen{ab^i}$. Hence, $q$ stabilizes all subgroups of order 7 in $G_3$, and each subgroup of order 7 is normal in $G_3$. In $G_4$, consider $\gen{ab} \in \ss$. Then
    \[q(ab)q\inv = (qaq\inv)(qbq\inv) = a^2b^4 \notin \gen{ab},\]
    so $q$ does not stabilize $\gen{ab}$. Hence, $\gen{ab}$ is not normal in $G_4$. Since $G_3$ and $G_4$ have different numbers of normal subgroups of order 7, then $G_3 \not\cong G_4$.
    \item By part (b), every group of order 147 has a normal Sylow 7-subgroup. Let $P \in \syl_7(G)$ and $Q \in \syl_3(G)$. Since $P \cap Q = 1$ because (49, 3) = 1, and $|G| = |P||Q|$, then $G = PQ$, and $G \cong P \semidir Q$. If the homomorphism is trivial, then we have the two abelian groups from part (a). If the homomorphism is non-trivial and $P$ is cyclic, then we have the unique non-abelian group from part (c). If the homomorphism is non-trivial and $P$ is not cyclic, then $\phi(Q)$ is a subgroup of order 3 in $\aut(P) \cong \gl_2(\f_7)$. By part (d), every subgroup of order 3 is conjugate in $\gl_2(\f_7)$ to a subgroup of $P$, so by \hyperref[ex5.5.6]{Exercise 5.5.6} we have the three non-abelian groups $G_1, G_3$, and $G_4$ from part (e). Hence, the six non-isomorphic groups described above are all the groups of order 147.
\end{solalph}

\begin{exercise}\label{ex5.5.11}
    Classify groups of order 28 (there are four isomorphism types).
\end{exercise}

\begin{sol}
    Note $28 = 2^2 \cdot 7$. We have the abelian groups $Z_7 \times Z_4$ and $Z_7 \times Z_2 \times Z_2$. Let $P \in \syl_7(G)$ and $S \in \syl_2(G)$. Then $|P| = 7$ and $|S| = 4$. By the Sylow theorems, we have $n_7 \equiv 1 \bmod 7$ and $n_7 \mid 4$, so $n_7 = 1$. Hence, the unique Sylow 7-subgroup is normal. Moreover, we have $P \cong Z_7$ and $S \cong Z_4$ or $Z_2 \times Z_2$. We then have two cases to consider:
    \begin{itemize}
        \item Suppose $S \cong Z_4 = \gen s$. If we consider the trivial homomorphism, then $G \cong P \times S \cong Z_7 \times Z_4$. Now consider the non-trivial homomorphism that sends $s$ to $\alpha$, a generator of the unique subgroup of order 2 in $\aut(P)$. Then $s$ acts on $P$, and $sps\inv = p\inv$. Hence, we have
        \[G \cong P \semidir S \cong Z_7 \semidir Z_4 = \set{p, s \mid p^7 = s^4 = 1, sps\inv = p\inv}.\]
        \item Suppose $S \cong Z_2 \times Z_2 = \gen s \times \gen t$. If we consider the trivial homomorphism, then $G \cong P \times S \cong Z_7 \times (Z_2 \times Z_2)$. Now consider a non-trivial homomorphism that sends some generator of $S$ to $\alpha$. We have the three homomorphisms
        \[\theta_1 = 
        \begin{cases}
            s \mapsto \alpha, \\
            t \mapsto \id,
        \end{cases} \quad
        \theta_2 = 
        \begin{cases}
            s \mapsto \id, \\
            t \mapsto \alpha,
        \end{cases} \quad
        \theta_3 = 
        \begin{cases}
            s \mapsto \alpha, \\
            t \mapsto \alpha,
        \end{cases}\]
        Observe that $\ker\theta_1 = \gen t, \ker\theta_2 = \gen s$, and $\ker\theta_3 = \gen{st}$. Hence, the semidirect products corresponding to $\theta_1, \theta_2$, and $\theta_3$ are all isomorphic to each other. Pick $\theta_1$ without loss of generality. Then $s$ acts on $P$, and $sps\inv = p\inv$. Moreover, $t$ centralizes $P$, hence $|tp| = 14$. Since
        \[s(tp)s\inv = stps\inv = (sts\inv)(sps\inv) = tp\inv = (tp)\inv,\]
        then
        \[G \cong P \semidir_{\theta_1} S \cong \set{p, s, t \mid p^7 = s^2 = t^2 = 1, st = ts, sps\inv = p\inv, tpt\inv = p} \cong D_{28}. \qh\] 
    \end{itemize}
\end{sol}

\begin{exercise}\label{ex5.5.12}
    Classify the groups of order 20 (there are five isomorphism types).
\end{exercise}

\begin{sol}
    We have the abelian groups $Z_5 \times Z_4$ and $Z_5 \times Z_2 \times Z_2$. Let $P \in \syl_5(G)$ and $S \in \syl_2(G)$. Then $|P| = 5$ and $|S| = 4$. By the Sylow theorems, we have $n_5 \equiv 1 \bmod 5$ and $n_5 \mid 4$, so $n_5 = 1$. Hence, the unique Sylow 5-subgroup is normal. Moreover, we have $P \cong Z_5$ and $S \cong Z_4$ or $Z_2 \times Z_2$. We then have two cases to consider:
    \begin{itemize}
        \item If $S \cong Z_4 = \gen s$, then $s$ has four candidates to be sent to: $\id, \alpha$, $\alpha^2$, or $\alpha^3$, where $\aut(Z_5) = \gen\alpha$. Define the homomorphisms
        \[\phi_1 : s \mapsto \id, \quad \phi_2 : s \mapsto \alpha, \quad \phi_3 : s \mapsto \alpha^2, \quad \phi_4 : s \mapsto \alpha^3\]
        Since $\phi_1$ is trivial, then $G \cong P \semidir_{\phi_1} S \cong P \times S \cong Z_5 \times Z_4$. Also observe that $\ker\phi_2 = \ker\phi_4 = 1$ while $\ker\phi_3 = \gen{s^2}$. Then the resulting semidirect products from $\phi_2$ and $\phi_4$ are isomorphic, since $\phi_2(S) = \phi_4(S) = \gen\alpha$, so by \hyperref[ex5.5.6]{Exercise 5.5.6}, we have $P \semidir_{\phi_2} S \cong P \semidir_{\phi_4} S$. Moreover, $\ker\phi_3 = \gen{s^2}$, which is the unique subgroup of order 2 in $\aut(P)$, so $P \semidir_{\phi_3} S$ is not isomorphic to $P \semidir_{\phi_2} S$. Hence, we have two non-abelian groups of order 20, and we have
        \[G \cong P \semidir_{\phi_2} S \cong \set{p, s \mid p^5 = s^4 = 1, sps\inv = p^2} \cong F_{20}.\]
        In the case of $\phi_3$, we have $sps\inv = p\inv$. Then
        \[G \cong P \semidir_{\phi_3} S \cong \set{p, s \mid p^5 = s^4 = 1, sps\inv = p\inv} \cong Z_5 \semidir Z_4.\]
        \item If $S \cong Z_2 \times Z_2 = \gen s \times \gen t$, then we have the trivial homomorphism $\theta_1 : S \to \aut(P)$, which gives $G \cong P \semidir_{\theta_1} S \cong P \times S \cong Z_5 \times (Z_2 \times Z_2)$. If $\alpha^2$ is the generator of the unique subgroup of order 2 in $\aut(P)$, then we have the three non-trivial homomorphisms
        \[\theta_2 = 
        \begin{cases}
            s \mapsto \alpha^2, \\
            t \mapsto \id,
        \end{cases} \quad
        \theta_3 = 
        \begin{cases}
            s \mapsto \id, \\
            t \mapsto \alpha^2,
        \end{cases} \quad
        \theta_4 =
        \begin{cases}
            s \mapsto \alpha^2, \\
            t \mapsto \alpha^2,
        \end{cases}\]
        This is the same situation as in the previous exercise, since $\ker\theta_2 = \gen t, \ker\theta_3 = \gen s$, and $\ker\theta_4 = \gen{st}$. Hence, the semidirect products corresponding to $\theta_2, \theta_3$, and $\theta_4$ are all isomorphic to each other. Pick $\theta_2$ without loss of generality. Then $s$ acts on $P$, and $sps\inv = p\inv$. Moreover, $t$ centralizes $P$, hence $|tp| = 10$. Since
        \[s(tp)s\inv = stps\inv = (sts\inv)(sps\inv) = tp\inv = (tp)\inv,\]
        then
        \[G \cong P \semidir_{\theta_2} S \cong \set{p, s, t \mid p^5 = s^2 = t^2 = 1, st = ts, sps\inv = p\inv, tpt\inv = p} \cong D_{20}. \qh\]
    \end{itemize}
\end{sol}

\begin{exercise}\label{ex5.5.13}
    Classify groups of order $4p$, where $p$ is a prime greater than 3. [There are four isomorphism types when $p \equiv 3 \bmod 4$, and five isomorphism types when $p \equiv 1 \bmod 4$.]
\end{exercise}

\begin{sol}
    We have a framework for each case of $p$, namely \hyperref[ex5.5.11]{Exercise 5.5.11} when $p \equiv 3 \bmod 4$, and \hyperref[ex5.5.12]{Exercise 5.5.12} when $p \equiv 1 \bmod 4$. The only difference is that we have $P \in \syl_p(G)$ and $S \in \syl_2(G)$, where $|P| = p$ and $|S| = 4$. By the Sylow theorems, we have $n_p \equiv 1 \bmod p$ and $n_p \mid 4$, so $n_p = 1$. Hence, the unique Sylow $p$-subgroup is normal. Moreover, we have $P = \gen q \cong Z_p$ and $S \cong Z_4$ or $Z_2 \times Z_2$. We then have two cases to consider:
    \begin{itemize}
        \item If $S \cong Z_4 = \gen s$, then $s$ has different candidates to be sent to in $\aut(P) = \units[p]$, which has order $p - 1$. If $p \equiv 3 \bmod 4$, then $p - 1 \equiv 2 \bmod 4$, hence $\phi(s)$ can only be sent to $\id$ or the unique element of order 2 in $\aut(P)$. If we have the latter case, then $sqs\inv = q\inv$, and we have
        \[G \cong P \semidir S \cong Z_p \semidir Z_4 = \set{q, s \mid q^p = s^4 = 1, sqs\inv = q\inv}.\]
        
        If $p \equiv 1 \bmod 4$, then $p - 1 \equiv 0 \bmod 4$, hence $\phi(s)$ can be sent to $\id$ or one of the two elements of order 4 in $\aut(P)$ (which are isomorphic since their images generate the same subgroup), or the unique element of order 2 in $\aut(P)$. If $\aut(P) = \gen\alpha$, let $\rho = \alpha^{(p - 1)/4}$, which has order 4. Then we have the three non-trivial homomorphisms
        \[\phi_1 : s \mapsto \rho, \quad \phi_2 : s \mapsto \rho^3, \quad \phi_3 : s \mapsto \rho^2.\]
        We know that $\phi_1$ and $\phi_2$ give isomorphic semidirect products, since $\phi_1(S) = \phi_2(S) = \gen\rho$. Hence, we have two non-abelian groups of order $4p$, and we have
        \[G \cong P \semidir_{\phi_1} S \cong \set{q, s \mid q^p = s^4 = 1, sqs\inv = q^k}.\]
        for some $k \in \units[p]$ such that $|k| = 4$. In the case of $\phi_3$, we have $sqs\inv = q\inv$. Then
        \[G \cong P \semidir_{\phi_3} S \cong \set{q, s \mid q^p = s^4 = 1, sqs\inv = q\inv}.\]
        \item If $S \cong Z_2 \times Z_2 = \gen s \times \gen t$, then we have the trivial homomorphism $\theta_1 : S \to \aut(P)$, which gives $G \cong P \semidir_{\theta_1} S \cong P \times S \cong Z_p \times (Z_2 \times Z_2)$. If $\rho^2$ is the generator of the unique subgroup of order 2 in $\aut(P)$, then we have the three non-trivial homomorphisms
        \[\theta_2 = 
        \begin{cases}
            s \mapsto \rho^2, \\
            t \mapsto \id,
        \end{cases} \quad
        \theta_3 =
        \begin{cases}
            s \mapsto \id, \\
            t \mapsto \rho^2,
        \end{cases} \quad
        \theta_4 =
        \begin{cases}
            s \mapsto \rho^2, \\
            t \mapsto \rho^2,
        \end{cases}\]
        Again, we know that $\ker\theta_2 = \gen t, \ker\theta_3 = \gen s$, and $\ker\theta_4 = \gen{st}$. Hence, the semidirect products corresponding to $\theta_2, \theta_3$, and $\theta_4$ are all isomorphic to each other. Pick $\theta_2$ without loss of generality. Then $s$ acts on $P$, and $sqs\inv = q\inv$. Moreover, $t$ centralizes $P$, hence $|tq| = 2p$. Since
        \[s(tq)s\inv = stqs\inv = (sts\inv)(sqs\inv) = tq\inv = (tq)\inv,\]
        then
        \[G \cong P \semidir_{\theta_2} S \cong \set{q, s, t \mid (tq)^2 = s^2 = 1, st = ts, sqs\inv = q\inv, tqt\inv = q} \cong D_{4p}. \qh\]
    \end{itemize}
\end{sol}

\begin{specialexercise}\label{ex5.5.14}
    This exercise classifies the groups of order 60 (there are thirteen isomorphism types). Let $G$ be a group of order 60, let $P$ be a Sylow 5-subgroup of $G$, and let $Q$ be a Sylow 3-subgroup of $G$.
    \begin{subproblems}
        \item Prove that if $P$ is not normal in $G$, then $G \cong A_5$. [See Section 4.5.]
        \item Prove that if $P \nsub G$ but $Q$ is not normal in $G$, then $G \cong A_4 \times Z_5$. [Show in this case that $P \leq Z(G), G/P \cong A_4$, a Sylow 2-subgroup $T$ of $G$ is normal, and $TQ \cong A_4$.]
        \item Prove that if both $P$ and $Q$ are normal in $G$, then $G \cong Z_{15} \semidir T$, where $T \cong Z_4$ or $Z_2 \times Z_2$. Show in this case that there are six isomorphism types when $T$ is cyclic (one abelian), and there are five isomorphism types when $T$ is the Klein 4-group (one abelian). [Use the same ideas as in the classification of groups of order 30 and 20.]
    \end{subproblems}
\end{specialexercise}

\begin{solalph}
    \item Clear, from Corollary 4.22.
    \item If $Q$ is not normal in $G$, then $n_3 \equiv 1 \bmod 3$ and $n_3 \mid 20$ implies that $n_3 = 4$ or 10. If $n_3 = 10$, then since $P \nsub G$, we know that $PQ_i$ is a subgroup of $G$ that is cyclic. But then $PQ_i \cap PQ_j = 5$ intersect at $P$, which means there would be $10 \cdot 10$ elements in $G$, which would be a contradiction. Hence, $n_3 = 4$. Again, each $PQ_i$ intersect non-trivially at $P$, and $P$ commutes with the non-identity elements of each $PQ_i$. Then there are $5 + 10 \cdot 4 = 45$ elements in $C_G(P)$. Since $C_G(P) \leq G$, then $|C_G(P)| = 60$ so that $C_G(P) = G$, showing that $P \leq Z(G)$.
    
    Consider $G/P$, where $|G/P| = 12$. If we consider $\pi$ to be the natural projection map from $G$ to $G/P$, we claim that the map $\psi : \syl_3(G) \to \syl_3(G/P)$ given by $\psi(Q_i) \to \pi(Q_i)$ is a bijection. Note that this is a well-defined map, since $\pi(Q_i) = Q_iP/P$ is a subgroup of $G/P$, and since $Q_iP \leq G$ with order 15, then $Q_iP/P$ has order 3, so $\pi(Q_i) \in \syl_3(G/P)$. To see that $\psi$ is injective, suppose that $\psi(Q_i) = \psi(Q_j)$, then $Q_iP/P = Q_jP/P$, which implies that $Q_iP = Q_jP$. Since $|Q_iP| = |Q_jP| = 15$, and $n_3(Q_iP) \equiv 1 \bmod 3$ and $n_3(Q_iP) \mid 5$, then $n_3(Q_iP) = 1$. Hence, $Q_i = Q_j$, so $\psi$ is injective. Since $n_3(G) = 4$, then $|\syl_3(G/P)| \geq 4$. By injectivity, it must be that $\psi$ is a bijection, hence $n_3(G/P) = 4$. We now consider the classification of groups of order 12, which are $Z_{12}, Z_2 \times Z_6, D_{12}, A_4$, and $Z_3 \semidir Z_4$.
    \begin{itemize}
        \item We know $Z_{12}$ and $Z_2 \times Z_6$ are abelian, so any group of order 3 in these groups is normal, which is a contradiction.
        \item For $D_{12} \cong S_3 \times Z_2$, let $\alpha \in S_3$ be a 3-cycle. Then $\gen{(\alpha, 1)}$ is a subgroup of order 3. Recall that conjugating $\alpha$ by any cycle in $S_3$ gives $\alpha$ or $\alpha^2$, hence $\gen{(\alpha, 1)}$ is normal in $D_{12}$, which is a contradiction.
        \item For $Z_3 \semidir Z_4$, we know that $Z_3$ is normal in this group, which is a contradiction.
    \end{itemize}
    Then the only possibility is that $G/P \cong A_4$.

    Let $T_1, T_2 \in \syl_2(G)$ be distinct, and consider $T_1P$ and $T_2P$. Since $P \leq Z(G)$, then $T_iP \cong T_i \times P$ for $i = 1, 2$. If there was some $t \in T_1$ such that $t \in T_2P$, observe that $|t|$ is divisible by 4. Since all non-identity elements of $P$ have order 5, then $t \in T_2$. But $T_1$ and $T_2$ are distinct by assumption, hence there must be some $t^* \in T_1$ such that $t^* \not\in T_2P$. Then $T_1P \neq T_2P$, and $T_1P/P, T_2P/P \in \syl_2(G/P)$. Since $n_2(G/P) = 1$, then $T_1P/P = T_2P/P$, which implies that $T_1P = T_2P$. But this is a contradiction, so there is a unique Sylow 2-subgroup of $G$, which we call $T$.

    Consider $TQ$. It is clear that $T \cap Q = 1$, and $|TQ| = 12$. If $Q \nsub TQ$, then since $P \leq Z(G)$ also normalizes $Q$, then $Q \nsub G$, which is a contradiction. Then $n_3(TQ) \neq 1$, and by a similar argument as above, we conclude that $TQ \cong A_4$. If we now consider $\phi : TQ \to \aut(P)$, note that $\phi$ is trivial since $TQ$ is generated by 3-cycles, while $\aut(P) \cong Z_4$ has no elements of order 3. Then $G \cong P \times TQ \cong Z_5 \times A_4$.
    \item If $P$ and $Q$ are both normal in $G$, then we have $PQ \nsub G$. Since $P \cap Q = 1$, then $PQ \cong P \times Q \cong Z_5 \times Z_3 \cong Z_{15}$. Put $P = \gen p$ and $Q = \gen q$. Let $T \in \syl_2(G)$, and consider a homomorphism from $T$ to $\aut(Z_{15})$, where $\aut(Z_{15}) \cong Z_4 \times Z_2$. From the text, we know that there are exactly 3 automorphisms of order 2:
    \[\alpha_1 = 
    \begin{cases}
        p \mapsto p\inv \\
        q \mapsto q
    \end{cases} \quad
    \alpha_2 = 
    \begin{cases}
        p \mapsto p \\
        q \mapsto q\inv
    \end{cases} \quad
    \alpha_3 = 
    \begin{cases}
        p \mapsto p\inv \\
        q \mapsto q\inv
    \end{cases}\]
    Moreover, there are 4 automorphisms of order 4:
    \[\beta_1 = 
    \begin{cases}
        p \mapsto p^2 \\
        q \mapsto q
    \end{cases} \quad
    \beta_2 =
    \begin{cases}
        p \mapsto p^2 \\
        q \mapsto q\inv
    \end{cases} \quad
    \beta_3 =
    \begin{cases}
        p \mapsto p^3 \\
        q \mapsto q
    \end{cases} \quad
    \beta_4 =
    \begin{cases}
        p \mapsto p^3 \\
        q \mapsto q\inv
    \end{cases}\]
    However, note that $\gen{\beta_1} = \gen{\beta_3}$ and $\gen{\beta_2} = \gen{\beta_4}$, so there are only two distinct subgroups of order 4 in $\aut(Z_{15})$. Without loss of generality, we may consider the automorphisms $\alpha_1, \alpha_2, \alpha_3, \beta_1$, and $\beta_2$. We then have the following cases to consider for the Sylow 2-subgroup $T$:
    \begin{itemize}
        \item Suppose $T = \gen t \cong Z_4$. Then we have the trivial homomorphism, which gives $G \cong Z_{15} \times Z_4$. Then we have the non-trivial homomorphisms
        \[\phi_1 : t \mapsto \alpha_1, \quad \phi_2 : t \mapsto \alpha_2, \quad \phi_3 : t \mapsto \alpha_3, \quad \phi_4 : t \mapsto \beta_1, \quad \phi_5 : t \mapsto \beta_2\]
        \begin{enumerate}
            \item For $G_1 = Z_{15} \semidir_{\phi_1} Z_4$, we have $tpt\inv = p\inv$ and $tqt\inv = q$. Since $q \in G_1$ is centralized, we have the factorization that $\gen{p, t} \cong Z_5 \semidir Z_4$. However, note that $\aut(Z_5) \cong Z_4$, so we have the presentation
            \[G_1 \cong \gen{p, t} \times \gen q \cong \hol(Z_5) \times Z_3.\]
            \item For $G_2 = Z_{15} \semidir_{\phi_2} Z_4$, we have $tpt\inv = p$ and $tqt\inv = q\inv$. Since $p \in G_2$ is centralized, we have the factorization
            \[G_2 \cong \gen{q, t} \times \gen p \cong (Z_3 \semidir Z_4) \times Z_5.\]
            \item For $G_3 = Z_{15} \semidir_{\phi_3} Z_4$, we have $tpt\inv = p\inv$ and $tqt\inv = q\inv$. Then
            \[G_3 \cong Z_{15} \semidir Z_4 = \set{p, q, t \mid p^5 = q^3 = t^4 = 1, pq = qp, tpt\inv = p\inv, tqt\inv = q\inv}.\]
            \item For $G_4 = Z_{15} \semidir_{\phi_4} Z_4$, we have $tpt\inv = p^2$ and $tqt\inv = q$. Since $q \in G_4$ is centralized, we have the factorization
            \[G_4 \cong \gen{p, t} \times \gen q \cong F_{20} \times Z_3.\]
            \item For $G_5 = Z_{15} \semidir_{\phi_5} Z_4$, we have $tpt\inv = p^2$ and $tqt\inv = q\inv$. Then
            \[G_5 \cong \set{p, q, t \mid p^5 = q^3 = t^4 = 1, pq = qp, tpt\inv = p^2, tqt\inv = q\inv} = F_{20} \semidir Z_3.\]
        \end{enumerate}
        \item Suppose $T = \gen u \times \gen v \cong Z_2 \times Z_2$. While this may seem like a lot of cases, observe that $\beta_i^2 = \alpha_1$ for all $i = 1, 2, 3, 4$, since either $q \mapsto q\inv \mapsto q$, or $p \mapsto p^2 \mapsto p^4 = p\inv$ ($p^3$ is handled similarly). Hence, we only need to consider the homomorphisms that send $u$ and $v$ to the the images of the $\alpha_i$'s in $\aut(Z_{15})$. Moreover, we know by \hyperref[ex5.5.6]{Exercise 5.5.6} that homomorphisms from the cyclic subgroups of $T$, i.e., $\gen u, \gen v, \gen{uv}$, to $\aut(Z_{15})$ that have conjugate (same, since $\aut(Z_{15})$ is abelian) images give isomorphic semidirect products. For completeness, we will list all potential non-trivial homomorphisms, but group them by their images in $\aut(Z_{15})$:
        \begin{itemize}
            \item With image $\gen{\alpha_1}$, we have the three homomorphisms
            \[\theta_1 = 
            \begin{cases}
                u \mapsto \alpha_1, \\
                v \mapsto \id,
            \end{cases} \quad
            \theta_2 =
            \begin{cases}
                u \mapsto \id, \\
                v \mapsto \alpha_1,
            \end{cases} \quad
            \theta_3 =
            \begin{cases}
                u \mapsto \alpha_1, \\
                v \mapsto \alpha_1,
            \end{cases}\]
            \item With image $\gen{\alpha_2}$, we have the three homomorphisms
            \[\rho_1 = 
            \begin{cases}
                u \mapsto \alpha_2, \\
                v \mapsto \id,
            \end{cases} \quad
            \rho_2 =
            \begin{cases}
                u \mapsto \id, \\
                v \mapsto \alpha_2,
            \end{cases} \quad
            \rho_3 =
            \begin{cases}
                u \mapsto \alpha_2, \\
                v \mapsto \alpha_2,
            \end{cases}\]
            \item With image $\gen{\alpha_3}$, we have the three homomorphisms
            \[\sigma_1 = 
            \begin{cases}
                u \mapsto \alpha_3, \\
                v \mapsto \id,
            \end{cases} \quad
            \sigma_2 =
            \begin{cases}
                u \mapsto \id, \\
                v \mapsto \alpha_3,
            \end{cases} \quad
            \sigma_3 =
            \begin{cases}
                u \mapsto \alpha_3, \\
                v \mapsto \alpha_3,
            \end{cases}\]
            \item With image $V_4 = \gen{\alpha_1, \alpha_2}$ (where $\alpha_1\alpha_2 = \alpha_3$), we have the homomorphisms
            \[\tau_1 = 
            \begin{cases}
                u \mapsto \alpha_1, \\
                v \mapsto \alpha_2,
            \end{cases} \hspace{-2pt}
            \tau_2 =
            \begin{cases}
                u \mapsto \alpha_2, \\
                v \mapsto \alpha_1,
            \end{cases} \hspace{-2pt}
            \tau_3 =
            \begin{cases}
                u \mapsto \alpha_1, \\
                v \mapsto \alpha_3,
            \end{cases} \hspace{-2pt}
            \tau_4 =
            \begin{cases}
                u \mapsto \alpha_3, \\
                v \mapsto \alpha_1,
            \end{cases} \hspace{-2pt}
            \tau_5 = 
            \begin{cases}
                u \mapsto \alpha_2, \\
                v \mapsto \alpha_3,
            \end{cases} \hspace{-2pt}
            \tau_6 =
            \begin{cases}
                u \mapsto \alpha_3, \\
                v \mapsto \alpha_2,
            \end{cases}\]
        \end{itemize}
        Hence, we may consider the semidirect products corresponding to $\theta_1, \rho_1, \sigma_1$, and $\tau_1$ without loss of generality. Then we have
        \begin{enumerate}
            \item For $G_\theta = Z_{15} \semidir_{\theta_1} (Z_2 \times Z_2)$, we have $u$ acting on $Z_{15}$ by $\alpha_1$, and $v$ centralizing $Z_{15}$. Then we have the factorization $\gen v \times \gen{p, q, u}$. Observe that $q$ is centralized by $u$, and $p$ is inverted by $u$. Hence, we have
            \[G_\theta \cong \gen v \times \gen q \times \gen{p, u} \cong Z_2 \times Z_3 \times D_{10} \cong Z_6 \times D_{10}.\]
            \item For $G_\rho = Z_{15} \semidir_{\rho_1} (Z_2 \times Z_2)$, we have $u$ acting on $Z_{15}$ by $\alpha_2$, and $v$ centralizing $Z_{15}$. Then we have the factorization $\gen v \times \gen{p, q, u}$. Observe that $p$ is centralized by $u$, and $q$ is inverted by $u$. Hence, we have
            \[G_\rho \cong \gen v \times \gen p \times \gen{q, u} \cong Z_2 \times Z_5 \times D_6 \cong Z_{10} \times D_6.\]
            \item For $G_\sigma = Z_{15} \semidir_{\sigma_1} (Z_2 \times Z_2)$, we have $u$ acting on $Z_{15}$ by $\alpha_3$, and $v$ centralizing $Z_{15}$. Then we have the factorization $\gen v \times \gen{p, q, u}$. Observe that both $p$ and $q$ are inverted by $u$. Hence, we have
            \[G_\sigma \cong \gen v \times \gen{p, q, u} \cong Z_2 \times D_{30}.\]
            However, we can simplify this further by noting that $u$ inverts the element $pqv$, which has order 30. Hence, we have
            \[G_\sigma \cong \gen{pqv, u} \cong D_{60}.\]
            \item For $G_\tau = Z_{15} \semidir_{\tau_1} (Z_2 \times Z_2)$, we have $u$ acting on $Z_{15}$ by $\alpha_1$, and $v$ acting on $Z_{15}$ by $\alpha_2$. Then we have the factorization $\gen{p, q, u, v}$. Observe that $p$ is inverted by $u$ and centralized by $v$, while $q$ is inverted by $v$ and centralized by $u$. Hence, we have
            \[G_\tau \cong \gen{p, u} \times \gen{q, v} \cong D_{10} \times D_6 \cong D_{10} \times S_6. \qh\]
        \end{enumerate}
    \end{itemize}
\end{solalph}

\begin{exercise}\label{ex5.5.15}
    Let $p$ be an odd prime. Prove that every element of order 2 in $\gl_2(\f_p)$ is conjugate to a diagonal matrix with $\pm 1$'s on the diagonal. Classify the groups of order $2p^2$. [If $A$ is a $2 \times 2$ matrix with $A^2 = I$ and $v_1, v_2$ is a basis for the underlying vector space, look at $A$ acting on the vectors $w_1 = v_1 + v_2$ and $w_2 = v_1 - v_2$.]
\end{exercise}

\begin{sol}
    If $A = \pm I_2$, then $A$ is already diagonal. We now prove that if every nonzero vector of a 2-dimensional vector space is an eigenvector of $A$, then $A$ is a scalar matrix. To see this, let $e_1$ and $e_2$ be eigenvectors of $A$ corresponding to eigenvalues $\lambda_1$ and $\lambda_2$, respectively. Consider the vector $e_1 + e_2$, which is also an eigenvector of $A$ with some eigenvalue $\lambda$. Then
    \[
    A(e_1 + e_2) = \lambda (e_1 + e_2) \implies \lambda_1 e_1 + \lambda_2 e_2 = \lambda e_1 + \lambda e_2 \implies (\lambda_1 - \lambda) e_1 + (\lambda_2 - \lambda) e_2 = 0.
    \]
    Since $e_1$ and $e_2$ are linearly independent, we must have $\lambda_1 = \lambda = \lambda_2$. Hence, $A$ is a scalar matrix.

    We now take the contrapositive of this, and assume that $A$ is not a scalar matrix. Then there exists a nonzero vector $v_1$ that is not an eigenvector of $A$. Let $v_2 = Av_1$. If $v_2 = \lambda v_1$, then $Av_1 = \lambda v_1$, contradicting that $v_1$ is not an eigenvector of $A$. Hence, $v_1$ and $v_2$ are linearly independent and form a basis of the vector space.

    Let $w_1 = v_1 + v_2$ and $w_2 = v_1 - v_2$. Then
    \begin{align*}
        Aw_1 & = A(v_1 + v_2) = Av_1 + Av_2 = v_2 + A v_2 = v_2 + A^2 v_1 = v_2 + v_1 = w_1, \\
        Aw_2 & = A(v_1 - v_2) = Av_1 - Av_2 = v_2 - A v_2 = v_2 - A^2 v_1 = v_2 - v_1 = -w_2.
    \end{align*}
    Let $P = \begin{bmatrix} w_1 & w_2 \end{bmatrix}$. Then
    \[P\inv AP = 
    \begin{bmatrix}
        1 & 0 \\
        0 & -1
    \end{bmatrix}.\]
    Hence, $A$ is similar to the diagonal matrix with entries $1$ and $-1$. 

    Let $|G| = 2p^2$, $P \in \syl_p(G)$, and $Q \in \syl_2(G)$. Since $n_p \equiv 1 \bmod p$ and $n_p \mid 2$, then $p$ being odd implies that $n_p = 1$, hence $P \nsub G$. Moreover, $|P| = p^2$ implies that $P \cong Z_{p^2}$, or $Z_p \times Z_p$. Put $Q = \gen q$. We split by cases:
    \begin{itemize}
        \item The abelian case is $Z_{p^2} \times Z_2$. If $P = \gen x \cong Z_{p^2}$, then $\aut(P) \cong Z_{p(p - 1)}$. Since $|q| = 2$ and $p$ is odd, then $p - 1$ is even and $\aut(P)$ has even order, hence it has a unique subgroup of order 2 generated by some $\alpha \in \aut(P)$, namely the inversion map. We then have $qxq\inv = x\inv$. It follows that
        \[G \cong \gen{x, q \mid x^{p^2} = q^2 = 1, qxq\inv = x\inv} \cong D_{2p^2}.\]
        \item The abelian case is $Z_p \times Z_p \times Z_2$. If $P = \gen x \times \gen y \cong Z_p \times Z_p$, then $\aut(P) \cong \gl_2(\f_p)$. Then there are three homomorphisms from $Q$ to $\aut(P)$, namely
        \[\phi_1 = 
        \begin{cases}
            x \mapsto x, \\
            y \mapsto y
        \end{cases} \quad
        \phi_2 = 
        \begin{cases}
            x \mapsto x\inv, \\
            y \mapsto y
        \end{cases} \quad
        \phi_3 = 
        \begin{cases}
            x \mapsto x, \\
            y \mapsto y\inv
        \end{cases} \quad
        \phi_4 = 
        \begin{cases}
            x \mapsto x\inv, \\
            y \mapsto y\inv
        \end{cases} \]
        Note that $\phi_1 = \id$, hence that is the same abelian case as before. Moreover, if we let $S$ be the exchange matrix, then $S\phi_2(q)S\inv = \phi_3(q)$. Hence the semidirect products corresponding to $\phi_2$ and $\phi_3$ are isomorphic. We then have the following non-isomorphic groups:
        \begin{itemize}
            \item $G_2 \cong (Z_p \times Z_p) \semidir_{\phi_2} Z_2$ has $y$ centralized, while $x$ is inverted. We then have $\gen y \times \gen{q, x}$, where $\gen{q, x} \cong D_{2p}$. Hence,
            \[G_2 \cong Z_p \times D_{2p}.\]
            \item $G_4 \cong (Z_p \times Z_p) \semidir_{\phi_4} Z_2$ has both $x$ and $y$ inverted by $q$. Hence, it has presentation
            \[G_4 \cong \gen{x, y, q \mid x^p = y^p = q^2 = 1, qxq\inv = x\inv, qyq\inv = y\inv} \cong (Z_p \times Z_p) \semidir Z_2. \qh\]
        \end{itemize}
    \end{itemize}
\end{sol}

\begin{exercise}\label{ex5.5.16}
    Show that there are exactly 4 distinct homomorphisms from $Z_2$ into $\aut(Z_8)$. Prove that the resulting semidirect products are the groups $Z_8 \times Z_2, D_{16}$, the quasidihedral group $QD_{16}$ (cf. \hyperref[ex2.5.11]{Exercise 2.5.11}), and the modular group $M$ (cf. \hyperref[ex2.5.14]{Exercise 2.5.14}). 
\end{exercise}

\begin{sol}
    Let $Z_8 = \gen a$ and $Z_2 = \gen b$. Note that $\aut(Z_8) \cong \units[8] = \{1, 3, 5, 7\} \cong V_4$. Hence, there are exactly 4 distinct homomorphisms from $Z_2$ into $\aut(Z_8)$, given by
    \[\phi_1 : b \mapsto a, \quad \phi_2 : b \mapsto a^3, \quad \phi_3 : b \mapsto a^5, \quad \phi_4 : b \mapsto a^7.\]
    The resulting semidirect products are:
    \begin{itemize}
        \item $G_1 \cong Z_8 \semidir_{\phi_1} Z_2 \cong Z_8 \times Z_2$, since $\phi_1$ is the trivial homomorphism.
        \item $G_2 \cong Z_8 \semidir_{\phi_2} Z_2$ has the automorphism $bab\inv = a^3$, hence $G_2 \cong QD_{16}$.
        \item $G_3 \cong Z_8 \semidir_{\phi_3} Z_2$ has the automorphism $bab\inv = a^5$, hence $G_3 \cong M$.
        \item $G_4 \cong Z_8 \semidir_{\phi_4} Z_2$ has the inversion map $bab\inv = a\inv$, hence $G_4 \cong D_{16}$.
    \end{itemize}
    By \hyperref[ex5.3.1]{Exercise 5.3.1}, we know that the three non-abelian groups are pairwise non-isomorphic. Moreover, $Z_8 \times Z_2$ is abelian and the others are non-abelian, so it is also non-isomorphic to the others.
\end{sol}

\begin{exercise}\label{ex5.5.17}
    Show that for any $n \geq 3$, there are exactly 4 distinct homomorphisms from $Z_2$ into $\aut(Z_{2^n})$. Prove that the resulting semidirect groups give 4 non-isomorphic groups of order $2^{n + 1}$. [Recall \hyperref[ex2.3.21]{Exercise 2.3.21} through \hyperref[ex2.3.23]{Exercise 2.3.23}.] (These four groups together with the cyclic group and the generalized quaternion group $Q_{2^{n + 1}}$ are all the groups of order $2^{n + 1}$ which possess a cyclic subgroup of index 2.)
\end{exercise}

\begin{specialexercise}\label{ex5.5.18}
    Show that if $H$ is any group, then there is a group $G$ that contains $H$ as a normal subgroup with the property that for every automorphism $\sigma$ of $H$, there is an element $g \in G$ such that conjugation by $g$ when restricted to $H$ is the given automorphism $\sigma$, i.e., every automorphism of $H$ is obtained as an inner automorphism of $G$ restricted to $H$.
\end{specialexercise}

\begin{sol}
    Let $H$ be a group, and let $G = H \semidir \aut(H)$, where the action is given by the natural action of $\aut(H)$ on $H$. Then $H$ is a normal subgroup of $G$. Let $\sigma \in \aut(H)$. Let $g = (1, \sigma) \in G$. Then for any $(h, 1) \in \wt H$, where $\wt H \cong H$, then
    \[(1, \sigma)(h, 1)(1, \sigma)\inv = ( \sigma(h), 1).\]
    Hence conjugation by $g$ restricted to $H$ gives the automorphism $\sigma$. We remark that $G = \hol(H)$.
\end{sol}

\begin{exercise}\label{ex5.5.19}
    Let $H$ be a group of order $n$, let $K = \aut(H)$, and form $G = \hol(H) = H \semidir K$ (where $\phi$ is the identity homomorphism). Let $G$ act by left multiplication on the left cosets of $K$ in $G$, and let $\pi$ be the associated permutation representation $\pi : G \to S_n$.
    \begin{subproblems}
        \item Prove that the elements of $H$ are coset representatives for the left cosets of $K$ in $G$, and with this choice of coset representatives, $\pi$ restricted to $H$ is the regular representation of $H$.
        \item Prove $\pi(G)$ is the normalizer in $S_n$ of $\pi(H)$. Deduce that, under the regular representation of any finite group $H$ of order $n$, the normalizer in $S_n$ of the image of $H$ is isomorphic to $\hol(H)$. [Show $|G| = |N_{S_n}(\pi(H))|$ using \hyperref[ex5.5.1]{Exercise 5.5.1} and \hyperref[ex5.5.2]{Exercise 5.5.2}.]
        \item Deduce that the normalizer of the group generated by an $n$-cycle in $S_n$ is isomorphic to $\hol(Z_n)$ and has order $n\phi(n)$.
    \end{subproblems}
\end{exercise}

\begin{exercise}\label{ex5.5.20}
    Let $p$ be an odd prime. Prove that if $P$ is a non-cyclic $p$ group, then $P$ contains a normal subgroup $U$ with $U \cong Z_p \times Z_p$. Deduce that, for odd primes $p$, a $p$-group that contains a unique subgroup of order $p$ is cyclic. (For $p = 2$, it is a theorem that the generalized quaternion groups $Q_{2^n}$ are the only non-cyclic 2-groups which contain a unique subgroup of order 2.) [Proceed by induction on $|P|$. Let $Z$ be a subgroup of order $p$ in $Z(P)$, and let $\bbar P = P/Z$. If $\bbar P$ is cyclic, then $P$ is abelian by \hyperref[ex3.1.36]{Exercise 3.1.36}---show that the result is true for abelian groups. When $\bbar P$ is not cyclic, use induction to produce a normal subgroup $\bbar H$ of $\bbar P$ with $\bbar H \cong Z_p \times Z_p$. Let $H$ be the complete preimage of $\bbar H$ in $P$, so $|H| = p^3$. Let $H_0 = \set{x \in H \mid x^p = 1}$, so that $H_0$ is a characteristic subgroup of $H$ of order $p^2$ or $p^3$ by \hyperref[ex5.4.9]{Exercise 5.4.9}. Show that a suitable subgroup of $H_0$ gives the desired normal subgroup $U$].
\end{exercise}

\begin{exercise}\label{ex5.5.21}
    Let $p$ be an odd prime, and let $P$ be a $p$-group. Prove that if every subgroup of $P$ is normal, then $P$ is abelian. (Note that $Q_8$ is a non-abelian 2-group with this property, so the result is false for $p = 2$.) [Use the preceding exercises and \hyperref[ex5.4.15]{Exercise 5.4.15}.]
\end{exercise}

\begin{exercise}\label{ex5.5.22}
    Let $F$ be a field, let $n$ be a positive integer, and let $G$ be a group of upper triangular matrices in $\gl_n(F)$ (cf. \hyperref[ex2.1.16]{Exercise 2.1.16})
    \begin{subproblems}
        \item Prove that $G$ is the semidirect product $U \semidir D$, where $U$ is the set of upper triangular matrices with 1's down the diagonal (cf \hyperref[ex2.1.17]{Exercise 2.1.17}), and $D$ is the set of diagonal matrices in $\gl_n(F)$.
        \item Let $n = 2$. Recall that $U \cong F$, and $D \cong F\unt \times F\unt$ (cf \hyperref[ex3.1.11]{Exercise 3.1.11}). Describe the homomorphism from $D$ into $\aut(U)$ explicitly in terms of these isomorphisms (i.e., show how each element of $F\unt \times F\unt$ acts as an automorphism on $F$).
    \end{subproblems}
\end{exercise}

\begin{exercise}\label{ex5.5.23}
    Let $K$ and $L$ be groups, let $n$ be a positive integer, let $\rho : K \to S_n$ be a homomorphism and let $H$ be the direct product of $n$ copies of $L$. In \hyperref[ex5.1.8]{Exercise 5.1.8}, an injective homomorphism $\psi$ from $S_n$ into $\aut(H)$ was constructed by letting the elements of $S_n$ permute the $n$ factors of $H$. The composition $\psi \circ \rho$ is a homomorphism from $G$ into $\aut(H)$. The \emph{wreath product} of $L$ by $K$ is the semidirect product $H \semidir K$ with respect to this homomorphism and is denoted by $L \wr K$ (this wreath product depends on the choice of permutation representation $\rho$ of $K$---if none is given explicitly, $\rho$ is assumed to be the left regular representation of $K$).
    \begin{subproblems}
        \item Assume $K$ and $L$ are finite groups and $\rho$ is the left regular representation of $K$. Find $|L \wr K|$ in terms of $|K|$ and $|L|$.
        \item Let $p$ be a prime, let $K = L = Z_p$ and let $\rho$ be the left regular representation of $K$. Prove that $Z_p \wr Z_p$ is a non-abelian group of order $p^{p+1}$ and is isomorphic to a Sylow $p$-subgroup of $S_{p^2}$. [The $p$ copies of $Z_p$ whose direct product makes up $H$ may be represented by $p$ disjoint $p$-cycles; these are cyclically permuted by $K$.]
    \end{subproblems}
\end{exercise}

\begin{exercise}\label{ex5.5.24}
    Let $n$ be an integer $> 1$. Prove the following classification: every group of order $n$ is abelian if and only if $n = p_1^{\alpha_1} p_2^{\alpha_2} \cdots p_r^{\alpha_r}$, where $p_1, \ldots, p_r$ are distinct primes, $\alpha_i = 1$ or $2$ for all $i \in \{1, \ldots, r\}$ and $p_i$ does not divide $p_j^{\alpha_j} - 1$ for all $i$ and $j$. [See \hyperref[ex4.5.56]{Exercise 4.5.56}.]
\end{exercise}

\begin{exercise}\label{ex5.5.25}
    Let $H(\f_p)$ be the Heisenberg group over the finite field $\f_p = \z/p\z$ (cf. \hyperref[ex5.4.20]{Exercise 5.4.20}). Prove that $H(\f_2) \cong D_8$, and that $H(\f_p)$ has exponent $p$ and is isomorphic to the first non-abelian group in Example 7.
\end{exercise}